\documentclass[12pt]{report}


%Packages-------------------------------------
\usepackage{amsmath}
\usepackage{amsfonts}
\usepackage{amssymb}
\usepackage{amsthm}
\usepackage{graphics}
\usepackage{graphicx}
\usepackage{tikz}
%\usepackage{hyperref}
\usepackage[utf8]{inputenc}
\usepackage[T1]{fontenc}
\usepackage{amsmath}
\usepackage{amsfonts}
\usepackage{amssymb}
\usepackage[version=4]{mhchem}
\usepackage{stmaryrd}
\usepackage{bbold}
\usepackage{graphicx}
\usepackage[export]{adjustbox}
\graphicspath{ {./images/} }
\usepackage{mathrsfs}
\usepackage{CJKutf8}
%\usepackage{xtocnic}
\usepackage{xepersian}

\settextfont{XB Zar}

%Layout---------------------------------------

%\usepackage[top=2cm, bottom=2cm, left=2cm, right=2cm]{geometry}

%Commands-------------------------------------


%Theorems-------------------------------------
\newtheorem{thm}{قضیه}[chapter]
\newtheorem{cor}[thm]{نتیجه}
\newtheorem{defn}[thm]{تعریف}
\newtheorem{prop}[thm]{گزاره}
\newtheorem{lmm}[thm]{لم}
\newtheorem{conj}[thm]{حدس}
\newtheorem{exm}[thm]{مثال}
\newtheorem{rem}[thm]{تذکر}
\newtheorem{note}[thm]{یادداشت}
\newtheorem{alg}[thm]{الگوریتم}


\begin{document}

%Title page ---------------------------------------------

\begin{figure}
\centering
\includegraphics[height=2.5cm]{UT-Logo.pdf}
\end{figure}

\begin{center}
پردیس علوم
\\
دانشکده ریاضی، آمار و علوم کامپیوتر
\end{center}

\begin{center}
%%%%%%%%%%%%%
\end{center}

\begin{center}
\huge{بهبود روش وارسی مدل با استفاده از نظریه تعبیر مجرد}
\end{center}

\begin{center}
%%%
\end{center}

\begin{center}
نگارنده
\end{center}
\begin{center}
\textbf{
پویا پرتو
}
\end{center}

\begin{center}
\begin{tabular}{rr}
استاد راهنمای اول:& دکتر مجید علی‌زاده
\\
استاد راهنمای دوم: & دکتر مجتبی مجتهدی
\\

\end{tabular}
\end{center}

\vspace{3cm}
\begin{center}
پایاننامه برای دریافت درجه کارشناسی ارشد
\\
در رشته علوم کامپیوتر
\end{center}

\begin{center}
تاریخ دفاع
\end{center}

\pagestyle{empty}
\pagenumbering{}

\newpage
%\pagestyle{headings}
%\setcounter{page}{1}
%\pagenumbering{roman}
\pagestyle{plain}
\setcounter{page}{1}
\pagenumbering{harfi}
%Abstract page-------------------------------
\chapter*{}
\section*{چکیده}
روش وارسی مدل یک روش قابل اعتماد برای بررسی صحت عملکرد برنامه‌های کامپیوتری است. بیان‌های مختلف این روش از منطق موجهات بهره می‌برند که چندان برای برنامه نویسان شناخته شده نیستند. در این رساله سعی شده یک بیان جدید از روش وارسی مدل شرح داده شود و مورد بررسی قرار گیرد که در آن به کمک نظریه تعبیر مجرد به جای منطق موجهات از عبارات منظم استفاده شده.

در این نوشته ابتدا به بیان مفاهیم اولیه‌ برای بیان جدید روش وارسی مدل پرداخته‌ایم سپس این روش را به سه صورت متفاوت بیان کرده‌ایم. صورت اول ساختار خاصی ندارد و صرفا روش در ادبیاتی نو بیان شده، صورت دوم ساختار عبارات منظم را به صورتبندی‌اش اضافه کرده و صورت سوم ساختار برنامه‌ها را به صورتبندی اضافه کرده تا روشْ به پیاده سازی نزدیک شود و معادل بودن این سه صورت نشان داده شده.\\\\\\


\textbf{کلمات کلیدی:}
وارسی مدل، نظریه تعبیر مجرد، معناشناسی دلالتی، پیوند گالوا، درستی یابی صوری، تحلیل ایستا، درستی یابی برنامه‌های کامپیوتری
%Dedication page-----------------------------
\chapter*{تقدیم به}
\section*{تقدیم به}

%Acknowledgement page------------------------
\chapter*{سپاسگزاری}
\section*{سپاسگزاری}

%copyright + declaration page
%\chapter*{اصالت اثر}
%\section*{اصالت اثر}
% هیچ قسمت از این پایان‌نامه، پیش از این در هیج موسسه تحصیلات عالی برای دریافت درجه تحصیلی استفاده نشده است. همچنین، هیچ قسمت از این پایان‌نامه برگردان فارسی تمامی یا قسمتی از یک اثر دیگر علمی (مانند مقاله، پایان‌نامه، و غیره) به زبانی دیگر نمی‌باشد. ارائه این پایان‌نامه توسط نگارنده به معاونت آموزشی (معاونت پژوهشی و تحصیلات تکمیلی برای ارشد و دکتری) به منزله تعهد نگارنده به اصالت متن و محتوای ارائه شده بر اساس یک کار پژوهشی در مدت تحصیل در دانشگاه تهران می باشد. در صورت اثبات خلاف این امر، مدرک تحصیلی اخذ شده توسط این پایان‌نامه از دانشگاه تهران، معتبر نمی باشد. 
%
%\chapter*{حق مالکیت معنوی}
%\section*{حق مالکیت معنوی}
%حق مالکیت معنوی این اثر متعلق به دانشگاه تهران می باشد. استفاده از مطالب این پایان‌نامه در فعالیت های تحقیقاتی با ذکر منبع بلامانع می‌باشد. در صورت استفاده تجاری، مانند چاپ این پایان‌نامه، هماهنگی لازم و اجازه کتبی از دانشگاه و نگارنده پایان‌نامه الزامی می‌باشد.

%\listoffigures

%-------------------------
%\pagestyle{plain}
%\setcounter{page}{1}
%\pagenumbering{harfi}
%-------------------------


\chapter*{پیشگفتار }

با توجه به پیشرفت روز‌ افزون علوم کامپیوتر و ورود کاربرد‌های آن به زندگی روزمره، پیشرفت در روش‌های ساخت و نگه‌داری برنامه‌ها نیازی آشکار به نظر می‌رسد. یکی از مسائل مهم در این زمینه بررسی صحت کارکرد برنامه‌های نوشته‌شده است. عدم صحت کارکرد برنامه‌های نوشته‌شده بسته به حساسیت کاری که یک برنامه انجام می‌دهد می‌تواند تبعات مختلفی داشته‌باشد. پرتاب ناموفق آریان 5\cite{ariane}  ، از مدار خارج شدن مدار‌گرد مریخ \cite{orbiter} و  تصادف هلیکوپتر چینوک \cite{chinook} چند نمونه از تبعات بزرگ این قضیه در گذشته بوده اند، هرچند که به‌سادگی می‌توان فجایع دیگری را مشابه این اتفاقات در زندگی روزمره‌ی انسان‌ها متصور شد.
برای کشف صحت کارکرد برنامه‌های کامپیوتری روش های متفاوتی ابداع شده‌اند که در ادامه به طور مختصر از آن‌ها یاد می‌کنیم اما پیش از آن به یک خاصیت مشترک همه ی این روش‌ها می‌پردازیم که ناکامل بودن است. منظور از ناکامل بودن این است که با استفاده از هیچ یک از روش‌هایی که داریم نمی‌توانیم هر خاصیتی را برای هر برنامه‌ای بررسی کنیم. به عبارت دیگر استفاده از هر روشی، محدودیت‌هایی دارد. و البته قضیه رایس  \cite{rice} به ما این تضمین را داده که روش کاملی اصلا وجود ندارد. قضیه رایس( به طور غیر رسمی) بیان می‌کند که مسئله ی بررسی هر خاصیت غیر بدیهی، برای همه ی برنامه‌ها، تصمیم ناپذیر است. این دلیلی بر این شده که روش‌های مختلفی برای این کار درست شوند که هر کدام می‌توانند حالت‌های خاصی از مسأله را حل بکنند.\\





یک دسته‌بندی برای این روش‌ها تقسیم آن‌ها به دو دسته‌ی پویا و ایستا است. روش‌های پویا روش‌هایی هستند که در آن‌ها تست برنامه با اجرای برنامه همراه است و روش‌های ایستا بدون اجرای برنامه‌ها انجام می‌شوند.
روش‌های پویا معمولاً با اجرای حالات محدودی از برنامه، تصمیم می‌گیرند که برنامه‌ای که نوشته‌ای ایم، انتظاراتمان را برآورده می‌کند یا خیر. اگر این روش بتواند تشخیص دهد برنامه‌ای درست کار نمی‌کند می‌توانیم با اطمینان نتیجه بگیریم که برنامه غلط نوشته‌شده، اما اگر برنامه‌ای از تست‌های ساخته‌شده با این روش‌ها با موفقیت عبور کند، نمی‌توان اطمینان حاصل کرد که برنامه درست کار می‌کند، زیرا ممکن است حالتی مشکل‌زا از اجرای برنامه وجود داشته باشد که در تست ها نیامده باشد. در  کتاب \cite{art}  به توضیح این روش‌ها پرداخته شده. این دسته از روش‌ها از موضوع اصلی کار ما دور هستند.
روش‌های ایستا معمولاً روش‌هایی هستند که از نظریه‌های مختلف در منطق ریاضی به عنوان ابزار بهره می‌برند تا بدون اجرای خود برنامه‌ها در مورد صحت اجرای آن‌ها نتیجه‌گیری کنند. به همین دلیل به بخشی مهم و بزرگی از این دستورات که از منطق استفاده می‌کنند روش‌های صوری هم گفته‌می‌شود. از معروف ترین این روش‌ها وارسی‌ مدل، روش‌های استنتاجی و استفاده از نظریه تعبیر مجرد است.
در روش وارسی مدل، یک مدل صوری متناهی از برنامه‌ی موردبررسی می‌سازیم که  همه‌ی حالات اجرای برنامه با آن قابل‌توصیف است، سپس با استفاده از یک زبان صوری که بتواند در مورد مدل هایمان صحبت‌کند، ویژگی‌های مورد بررسیمان را بیان‌می‌کنیم و در نهایت صحت ویژگی‌های بیان‌شده را بررسی می کنیم. مقاله \cite{clarke} شروع این روش‌ها بوده که این کار را با استفاده از نوعی مدل کریپکی  \cite{kripke} و نوعی منطق زمانی به نام منطق زمانی خطی \cite{clarke} انجام داده که روشی است با دقت و البته هزینه‌ی محاسباتی بسیار بالا. \cite{logicincs} یک منبع بسیار مقدماتی در این زمینه و کتاب\cite{clarke2} یک مرجع سنتی در این زمینه است.
در روش‌های استنتاجی که شاید بتوان یکی از ابتدایی ترین آن‌ها را استفاده از منطق هور\cite{hoare} دانست، درستی کارکرد برنامه‌هایمان را با ارائه‌ی یک درخت اثبات در یک دستگاه استنتاجی که متناسب با زبان برنامه‌هایمان ساخته‌شده، نشان می‌‌دهیم. در این روش هم اگر بتوانیم درستی یک برنامه را اثبات کنیم، دیگر به‌طور تئوری، خیالی آسوده از درستی برنامه خواهیم داشت اما ساختن درخت اثبات در یک نظریه برهان می‌تواند چالش بر‌انگیز باشد چون این یک مسئله‌ی \lr{NP-Hard} است. در\cite{logicincs} به منطق هور به طور مقدماتی پرداخته شده. همین‌طور کتاب\cite{softwarefoundations} نیز به پیاده‌سازی منطق هور در زبان \lr{coq} پرداخته.   \lr{coq} نیز یک اثبات‌یار است که بر اساس یک نظریه نوع وابسته کار می‌کند. برای اطلاعات بیشتر در مورد چگونگی طرز کار این اثبات‌یار و تئوری بنیادین آن می شود به کتاب\cite{chlipala} مراجعه‌کرد. تئوری مورد شرح در\cite{dynamic} نیز می‌تواند در این مسیر به کار گرفته شود.
نظریه تعبیر مجرد\cite{cousot1} نیز یک نظریه ریاضیاتی است که در این بحث می‌توان گفت، به‌نوعی سعی می‌کند از روی معناشناسی یک برنامه‌ی کامپیوتری\cite{winskel} ، یک تقریب بسازد. منظور از تقریب، یک دستگاه کوچک‌تر از معناشناسی اصلی است که رفتارش زیرمجموعه‌ی رفتارهای دستگاه اصلی است. سعی بر این است که دستگاه جدیدی که می‌سازیم به لحاظ محاسباتی هم ساده‌تر از معناشناسی اصلی کار کند تا بتوانیم خواص آن را راحت‌تر بررسی کنیم. در این صورت هر نتیجه‌ای که در مورد خواص جدید بگیریم را می‌توانیم در مورد خود برنامه هم بیان کنیم اما می‌دانیم که در این صورت هم به همه‌ی حقایق دست پیدا نکرده ایم. در مورد این نظریه نیز به‌تازگی کتاب\cite{cousotbook} منتشر شده که حاصل نزدیک به 5 دهه کار مبدع این نظریه، پاتریک کوزو، است. همین‌طور\cite{statica} نیز در مورد پیاده‌سازی این نظریه بحث کرده.



\tableofcontents

%---------------------
\pagestyle{plain}
\setcounter{page}{1}
\pagenumbering{arabic}
%---------------------

\chapter{مقدمه}

در این فصل به عنوان مقدمه، روش وارسی مدل و نظریه تعبیر مجرد، به‌طور مختصر، معرفی شده‌اند. در فصل‌های بعدی، با هدف بهبود روش وارسی مدل، صورت جدیدی از این روش ارائه شده و مورد بررسی قرار گرفته است، بنابراین لازم است که ابتدا، این دو موضوع معرفی شوند.

\section{روش وارسی مدل}

روش وارسی مدل یک روش صوری است که برای درستی‌یابی سیستم‌های مختلف استفاده می‌شود. در این روش معمولا ابتدا یک ماشین حالات متناهی از روی سیستم مورد بررسی ساخته می‌شود، سپس بررسی‌هایی که قرار است روی سیستم اصلی انجام شوند، روی این ماشین( مدل) انجام می‌شود. 

از این روش در بررسی صحت کارکرد برنامه‌های کامپیوتری استفاده می‌شود، اما این تنها مورد استفاده‌ی این روش نیست. و هر سیستمی که قابلیت بیان شدن به صورت صوری را داشته باشد، با این روش قابل بررسی است. مثلا می‌توان از این روش برای بررسی صحت عملکرد یک برنامه برای قطارهای شهری استفاده کرد. در یک برنامه‌ برای قطارهای شهری، نباید امکان حضور دو قطار روی یک ریل در یک زمان وجود داشته باشد( که معنی تصادف بین دو قطار را می‌دهد) و می‌توان از روش وارسی مدل برای اطمینان از عدم وجود چنین ویژگی نامطلوبی استفاده کرد. مثال های دیگر استفاده‌ی این روش در علوم کامپیوتر بررسی صحت عملکرد یک پردازنده یا الگوریتم تقسیمِ وظایف یک سیستم عامل است. این مثال‌ها هیچ کدام یک برنامه‌ی کامپیوتری نیستند( هر چند که ممکن است مجبور باشیم از یک برنامه‌ی کامپیوتری برای پیاده سازی آن‌ها کمک بگیریم که در آن صورت بررسی صحت عملکرد آن برنامه‌ی کامپیوتری داستانی دیگر خواهد داشت)، اما قابل بیان به صورت صوری به جای زبان طبیعی هستند.

روش وارسی مدل برای بیان خواص مورد بررسی از منطق‌های زمانی مختلف استفاده می‌کند. منطق زمانی یک نوع منطق موجهات است. منطق‌های موجهات از گسترش زبان منطق کلاسیک، با اضافه کردن ادات وجهی گوناگون، ساخته می‌شوند. این ادوات غالبا در زبان طبیعی نقش قید را دارند. منطق‌های زمانی دسته‌ای از منطق‌های موجهات هستند که به صوری‌گری ما مفهوم زمان را اضافه می‌کنند، یعنی قیدهایی مانند فعلا، بعدا، و قبلا( که مورد آخری کمتر رایج است). منطقی که در اینجا بیان می‌کنیم \lr{LTL} نام دارد که یکی از منطق‌های زمانی است که برای روش وارسی مدل استفاده می‌شود. البته در مورد قیدهای مذکور، اشاره به این نکته ضروری است که در بیانی که در اینجا از این منطق ارائه داده‌ایم، ادوات جدید به‌طور مستقیم بیانگر این قید‌ها نیستند، هرچند که به کمک ادوات جدید می‌توان ادواتی برای هر یک از این قیود ساخت.
این تعاریف از \cite{buchi} آورده شده‌اند.

ابتدا زبان این منطق را بیان می‌کنیم و سعی می‌کنیم، به طور غیر دقیق، در مورد معنای فرمول‌های این زبان به خواننده یک درک شهودی بدهیم، سپس به سراغ معناشناسی صوری این منطق می‌رویم.

\subsection{زبان \lr{LTL}}
\begin{defn}
	هر عضو مجموعه‌ی $\Phi$ یک فرمول در زبان \lr{LTL} است( و $\Pi$ مجموعه‌ی‌( شمارای نامتناهی) فرمول‌های اتمی است و $\pi \in \Pi$):
	$$
	\Pi \subset \Phi,
	$$
	$$
	\phi \in \Phi \Leftrightarrow
	\phi ::= \pi | \phi \lor \phi |
	\neg \phi |
	\bigcirc \phi |
	\phi \mathcal{U}\phi 
$$	
	
\end{defn}
در این منطق، ما زمان را با اعداد طبیعی نشان می‌دهیم. یعنی برای یک فرمول، زمان از عدد ۰ شروع شده و تا ابد ادامه خواهد داشت و حین گذر زمان ممکن است ارزش فرمول‌ها تغییر کند. مسلما پس از بررسی معناشناسی صوری بهتر می‌شود این مفهوم را به طور شهودی حس کرد، اما به هر حال به خواننده پیشنهاد می‌شود، پیش از رسیدن به آن بخش به ادامه‌ی این بخش که در تلاش است یک درک شهودی از معنای فرمول‌ها بدهد، توجه کند. 

در این زبان ادوات کلاسیک 
$\neg, \lor$
هستند با همان معنایی که در منطق گزاره‌ای کلاسیک داشتند.  
در ادوات جدید 
$\bigcirc \phi$
به معنای برقرار بودن این فرمول دقیقا در لحظه‌ی بعدی( دقیقا یک لحظه) است، مثلا در شکل زیر با در نظر گرفتن اینکه در زمان 0 هستیم، این فرمول در لحظه‌ی ۱ برقرار است.
	\begin{center}
	\begin{tikzpicture}
	\node at (1,0.35) (phi) {$\phi$};
	\node at (1,-0.35) (1) {$1$};
	\node at (0,-0.35) (0) {$0$};
	\node at (2,-0.35) (n) {$n$};
	\node at (3,-0.35) (1) {$n+1$};
	\node at (4,-0.35) (1) {$\infty$};
	\fill (0,0) circle (0.1cm);
	\fill (1,0) circle (0.1cm);
	\fill (2,0) circle (0.1cm);
	\fill (3,0) circle (0.1cm);
	\path (0,0) edge (1,0) 
	(1,0) edge[dotted] (2,0)
	(2,0) edge (3,0)
	(3,0) edge[dotted] (4,0);
	\end{tikzpicture}
\end{center}
$\phi \mathcal{U}\psi$
به این معنی است که فرمول سمت چپی حداقل تا قبل از اینکه فرمول سمت راستی برقرار شود، برقرار است.( مثلا اگر بگوییم "تا وقتی که باران نباریده زمین خشک است" در این صورت "زمین خشک است" به جای فرمول سمت چپ و "باران باریده است" فرمول سمت راست است).
\begin{center}
	\begin{tikzpicture}
	\node at (1,0.35) (phi) {$\phi$};
	\node at (0,0.35) (phi) {$\phi$};
	\node at (2,0.35) (phi) {$\phi$};
	\node at (3,0.35) (phi) {$\psi$};
	\node at (1,-0.35) (1) {$1$};
	\node at (0,-0.35) (0) {$0$};
	\node at (2,-0.35) (n) {$n$};
	\node at (3,-0.35) (1) {$n+1$};
	\node at (4,-0.35) (1) {$\infty$};
	\fill (0,0) circle (0.1cm);
	\fill (1,0) circle (0.1cm);
	\fill (2,0) circle (0.1cm);
	\fill (3,0) circle (0.1cm);
	\path (0,0) edge (1,0) 
	(1,0) edge[dotted] (2,0)
	(2,0) edge (3,0)
	(3,0) edge[dotted] (4,0);
	\end{tikzpicture}
\end{center}

این زبان را می‌توان با ادوات بیشتری از آنچه آورده‌ایم بیان کرد و البته بیان‌های دیگری هم بسته به بحث متداول هستند، اما در اینجا یک شکل ساده از این زبان را آورده‌ایم که به غیر از ادوات منطق گزاره‌ای دو ادات دیگر را در زبان خود دارد. دلیل وجود ادوات متفاوت، می‌تواند راحت‌تر کردن بیان خواص باشد. همان طور که استفاده نکردن از یا و شرطی در منطق گزاره‌ای می‌تواند به سخت کردن بیان جملات در چارچوب این منطق منجر شود، حذف این ادوات وجهی هم بیان خواص را در این منطق مشکل می‌سازد. 

حال که به درکی شهودی از معنای فرمول‌های این زبان رسیده‌ایم، به بیان صوری این مفاهیم می‌پردازیم.

\subsection{معناشناسی \lr{LTL}}

مدل‌های این منطق را به صورت توابع
$M:\mathbb{N}_0 \rightarrow \mathit{P}(\Pi)$ 
تعریف می‌کنیم. به عبارت دیگر، هر مدل یک تابع است که هر عدد طبیعی را به یک مجموعه از فرمول‌های اتمی می‌برد. این در واقع به این معنی است که یک مدل مشخص می‌کند که در هر لحظه کدام یک از فرمول‌های اتمی درست هستند. مثلا، در مدلی به نام $M$ در واقع
$M(5)$
مجموعه‌ی اتم‌هایی است که در لحظه‌ی 5 طبق این مدل درست هستند و اگر اتمی در این مجموعه حاضر نباشد، در لحظه‌ی 5، ارزش غلط دارد.
درستی یک فرمول در یک مدل را با 
$M,i$
نشان می‌دهیم و 
$M,i \models \phi$
به این معنی است که فرمول $\phi$، در لحظه‌ی $i$در مدل $M$ ارزش درست دارد. این مفهوم را، به صورت بازگشتی، به شکل زیر تعریف می‌کنیم:


\begin{flushleft}
$	M,i \models \pi \:\:\: \mathit{iff} \:\:\: \pi \in M(i)$\\
$	M,i \models \neg \phi \:\:\: \mathit{iff} \:\:\: M,i\nvDash \phi$\\
$	M,i \models \phi \lor \psi \:\:\: \mathit{iff} \:\:\: M,i \models \phi \:\:\: \mathit{or} \:\:\: M,i \models \psi$\\
	 $M,i \models \bigcirc \phi  \:\:\:  \mathit{iff} \:\:\: M,i+1 \models \phi$\\
	 $M,i \models \phi \mathcal{U} \psi \:\:\: \mathit{iff} \:\:\: 
	 \exists k \geq i \in \mathbb{N}_0: \forall i\leq j< k: M,j \models \phi \:\:\: \mathit{and} \:\:\: M,k \models \psi$
\end{flushleft}


یک فرمول را ارضاپذیر می‌گوییم اگر و تنها اگر مدلی وجود داشته باشد که فرمول در آن صادق باشد.
اگر یک فرمول در هر مدلی صادق باشد، آن فرمول را \emph{معتبر} می‌گوییم.\\


\section{نظریه تعبیر مجرد}

به‌طور خلاصه، نظریه تعبیر مجرد یک چارچوب برای ساختن یک تقریب از معناشناسی یک زبان‌ برنامه نویسی است. 

معناشناسی یک زبان یک مدل ریاضیاتی مجرد است که چگونگی رفتار برنامه‌ها در این زبان را توصیف می‌کند. 
تقریب نیز یک معناشناسی دیگر است که قرار است بخشی( نه همه) از رفتارهای یک برنامه‌ی کامپیوتری در حال اجرا در یک زبان را توصیف کند. این که تقریب چیست، یک معناشناسی را در چه زمانی می‌توانیم تقریبی برای معناشناسی دیگری بدانیم و از یک تقریب چه چیزهایی را می‌توانیم بفهمیم و مواردی دیگر در مورد ارتباط بین دو مدل ریاضیاتی که درباره‌ی معنای برنامه‌ها در یک زبان برنامه‌نویسی واحد صحبت می‌کنند، همگی موضوع بحث در نظریه‌ي تعبیر مجرد است.
پس تا اینجا مشخص شد که نظریه‌ی تعبیر مجرد در مورد ارتباط بین معناشناسی‌های مختلف صحبت می‌کند. 

برای شروع بحث صوری در مورد این نظریه، از مفهوم دامنه و معناشناسی شروع می‌کنیم.
در واقع، این نوع از مشخص کردن معناشناسی یک زبان برنامه نویسی را معناشناسی دِلالَتی نامیده‌اند. در فصول آینده با یک معناشناسی از این نوع سر و کار خواهیم داشت.
\begin{defn}
	(معناشناسی و دامنه): اگر $\mathbb{P}$ مجموعه‌ی برنامه‌ها در یک زبان برنامه نویسی باشد، به تابع 
	$\mathcal{S}:\mathbb{P} \rightarrow D$
	یک معناشناسی و به مجموعه‌ی $D$ یک دامنه می‌گوییم.
\end{defn}
همان‌طور که از تعریف مشخص است، برای این که بتوانیم معنای برنامه‌های کامپیوتری موجود در یک زبان را تعریف کنیم، به یک مجموعه به اسم دامنه احتیاج داریم. تلاش برای پی بردن به این که در یک معناشناسی باید چه مجموعه‌ای را به عنوان دامنه در نظر گرفت، منجر به تولد یک مبحث به نام نظریه‌ی دامنه شده است.

در فصل‌های بعدی، با یک معناشناسی دلالتی سر و کار خواهیم داشت. 
پس از ارائه‌ی یک زبان برنامه نویسی، یک معناشناسی برای آن زبان معرفی می‌کنیم که معناشناسی رد پیشوندی نام دارد. در این معناشناسی، دامنه یک مجموعه است که شامل موجوداتی به نام رد پیشوندی است. هر رد پیشوندی یک دنباله است که در هر عضو آن اطلاعات موجود در حافظه و مرحله‌ی اجرای برنامه مشخص شده است. 

اما فعلا که در حال صحبت در مورد نظریه‌ی تعبیر مجرد هستیم، معناشناسی خاصی را معرفی نمی‌کنیم و بحث را کلی‌تر پیش می‌بریم. نظریه تعبیر مجرد برای معناشناسی‌ها یک چارچوب مشخص کرده و فقط در مورد معناشناسی‌هایی که در این چارچوب می‌گنجند می‌تواند صحبت کند. یکی از محدودیت‌های این چارچوب این است که دامنه باید یک ترتیب جزئی باشد.
\begin{defn}
	(ترتیب جزئی): یک مجموعه‌ی $D$ را به همراه یک رابطه‌ی $\leq$ روی آن مجموعه ترتیب جزئی می‌گوییم، اگر و تنها اگر خواص زیر را داشته باشند:
	$$\blacktriangleleft \forall a \in D: a\leq a$$
	$$\blacktriangleleft \forall a,b \in D: a \leq b \land b \leq a \rightarrow a=b$$
	$$\blacktriangleleft \forall a,b,c \in D: a \leq b \land b \leq c \rightarrow a \leq c$$
\end{defn}

حال به تعریف بخش بزرگتری از این چارچوب می‌پردازیم. در جبر مجرد مفهومی به اسم تناظر گالوا وجود دارد. این تناظر بین مجموعه‌ای از گروه‌ها و مجموعه‌ای از توسیع میدان‌هایی خاص وجود دارد که به بحث ما مربوط نمی‌شوند. این تناظر یک شکل نظریه ترتیبی هم دارد که در آن به جای مجموعه‌ای از گروه‌ها و میدان‌ها، دو مجموعه‌ی جزئا مرتب داریم. می‌توان گفت در واقع این یک مجرد سازی تناظری است که از جبر آمده. 

به شکل ضعیف‌تر نظریه ترتیبی این تناظر اتصال گالوا می‌گویند که در نظریه تعبیر مجرد به عنوان شرط تقریب تعریف شده است، به این معنی که دامنه‌ی یک معناشناسی باید با دامنه‌ی تقریبش یک اتصال گالوا داشته باشد.

\begin{defn}
	(اتصال گالوا): برای دو ترتیب جزئی
	$(A,\leq)$ 
	و
	$(C,\subseteq)$
	زوج 
	$\langle \alpha , \gamma \rangle$
	شامل دو تابع
	$\gamma:A \rightarrow C$
	و
	$\alpha:C \rightarrow A$،
	 یک اتصال گالوا است اگر و تنها اگر
	$$\forall c \in C :\forall a \in A: \alpha(c)\leq a \leftrightarrow c \subseteq \gamma(a)$$.
\end{defn}









\chapter{صوری‌گری جدید برای روش وارسی مدل}

محوریت کار ما قرار است\cite{calcul} باشد که در آن سعی شده روش وارسی مدل با کمک نظریه تعبیر مجرد، بهبود داده شود. در\cite{clarke} روشی که معرفی شده در‌واقع ویژگی برنامه‌ها را به کمک منطق های \lr{LTL} یا \lr{CTL} بیان می‌کند. خود برنامه‌ها هم با کمک معناشناسی این منطق ها که نوع خاصی از مدل های کریپیکی به اسم سیستم گذار هستند توصیف می شوند. اما در\cite{calcul} کاری که انجام شده به این شکل است که منطق های \lr{LTL} و \lr{CTL} با عبارات منظم\cite{kleene56} جایگزین شده‌اند. علت این کار دو نکته بوده، اولی اینکه استفاده از عبارات منظم به جای منطق های نام برده شده می‌تواند برای برنامه نویسان ساده‌تر باشد و دومی اینکه عبارت منظم از از منطق های نام برده شده قدرت بیان بالاتری دارد.\cite{regisbetter} در ادامه ی کار وارسی مدل با استفاده از موجودات جدید تعریف شده به سه شکل مختلف بیان شده. در هر مرتبه بیان وارسی مدل، به‌گفته‌ی نویسنده، "ساختارمند"تر شده. می‌توان دریافت که فایده ساختارمندتر بودن بیان این است که پیاده‌سازی راحت‌تری دارد.\\
حال به بیان و بررسی تعاریف و خواص آن‌ها می‌پردازیم.

\section{نحو زبان مورد بررسی‬}

زبان بیان برنامه ها زیرمجموعه ای از دستورات زبان \lr{C} است، به شکل زیر:
$$\mathsf{x,y},... \in \mathbb{X}\hspace{4.65cm}$$
$$\mathsf{A} \in \mathbb{A} ::=\mathsf{1\:|\:x\:|\:A_1 - A_2\hspace{2.4cm}}$$  
$$\mathsf{B \in \mathbb{B} ::=A_1<A_2 \:|\: B_1 \: nand\: B_2\hspace{1.2cm}}$$
$$\mathsf{E \in \mathbb{E}::= A \: | \: B\hspace{4cm}}$$
$$\mathsf{S\in \mathbb{S} ::=\hspace{5cm}  }$$
$$\mathsf{x\doteq A; \hspace{2.25cm}}$$
$$\mathsf{|\:\:\:;\hspace{3.75cm}}$$
$$\mathsf{|\:\:\:if\:(B)\:S\:|\:if\:(B)\:S\:else\:S}$$
$$\mathsf{|\:\:\:while\:(B)\:S\: | \: break;\hspace{0.65cm}}$$
$$\mathsf{|\:\:\:\{Sl\}\hspace{3.15cm}}$$
$$\mathsf{Sl \in \mathbb{SL}::=Sl\:\:\:S\:|\:\backepsilon}\hspace{3.05cm}$$
$$\mathsf{P\in \mathbb{P}\:::=Sl}\hspace{4.25cm}$$

\vspace{1cm}
در اینجا زیر مجموعه‌ای از دستورات زبان \lr{C} را داریم. همین‌طور که قابل مشاهده است این زبان تا حد ممکن کوچک شده. علت این کار را بعدا عمیق‌تر حس خواهیم‌کرد. علتْ ساده‌تر شدن کار برای ارائه‌ی معناشناسی و تعبیر مجرد آن است. در اینجا راحتی آن برنامه‌نویسی که قرار است با این زبان برنامه بنویسد مطرح نبوده چون اصلا این زبان برای این کار ساخته نشده. نویسنده‌ی \cite{calcul} در اینجا صرفا می‌خواسته فرآیند را نشان دهد. اگر به فرض برای زبانی مانند پایتون بخواهیم درستی‌یابی با استفاده از روش ارائه شده را درست کینم، می‌توانیم همه‌ی راهی که در \cite{calcul} برای زبان توصیف شده، رفته‌شده را برای پایتون هم برویم و به یک تحلیل‌گر ایستا برای پایتون برسیم.
در مورد قدرت بیان این زبان هم می‌توانیم بگوییم که می‌توانیم باقی اعداد را از روی عدد ۱ و عملگر منها بسازیم. مثلا ابتدا 0 را به کمک ۱-۱ می سازیم و سپس با استفاده از 0 می‌توانیم یکی یکی اعداد منفی را بسازیم و سپس بعد از آن به سراغ اعداد مثبت می‌رویم که با کمک 0 و هر عدد منفی‌ای که ساختیم، ساخته می شوند. باقی اعداد و حتی باقی عملگر‌ها( یعنی به غیر از اعداد طبیعی) نیز از روی آنچه داریم قابل‌ساختن است. در مورد عبارت‌های بولی نیز داستان به همین منوال است. یعنی اینجا صرفا ادات شفر تعریف شده و باقی عملگر‌های بولی را می‌توان با استفاده از همین عملگر ساخت. باقی دستورات نیز دستورات شرط و حلقه هستند. باقی دستورات قرار است مطابق چیزی که از زبان \lr{C} انتظار داریم کار بکنند. در مورد دستور \lr{$\mathsf{break;}$} ذکر این نکته ضروری است که اجرای آن قرار است اجرای برنامه را از دستوری بعد از داخلی‌ترین حلقه‌ای که \lr{$\mathsf{break;}$} داخلش قرار دارد ادامه ‌دهد. در پایان می توان ثابت کرد که این زبان هم قدرت با مدل دیویس\cite{davis} است. 
توجه داریم که هر‌چه در این بخش در‌مورد معنای دستورات این زبان گفتیم، به هیچ وجه صوری نبود و صرفا درک شهودی ای که از معنای اجرای هر‌یک از دستورات داشتیم را بیان کردیم. بیان صوری معنای برنامه‌ها را، که برخلاف درک شهودی‌مان قابل انتقال به کامپیوتر‌ است، در ادامه بیان خواهیم‌کرد. طبیعتا این بیان صوری از روی یک درک شهودی ساخته شده‌است.

\section{معناشناسی زبان مورد بررسی‬}
معناشناسی زبانی را که در بخش پیش آوردیم با کمک مفاهیمی به نام برچسب و رد پیشوندی و عملگر چسباندن روی دو رود پیشوندی مختلف تعریف خواهیم‌کرد و نام این معناشناسی نیز معناشناسی رد پیشوندی است.\\

\subsection{برچسب‌ها}

با‌وجود اینکه خود زبان \lr{C} در قسمتی از زبان خود چیز‌هایی به نام برچسب دارد اما همین‌طور که در بخش پیشین دیدیم، در زبانی که اینجا در حال بحث روی آن هستیم خبری از برچسب‌ها نیست. اما برای تعریف صوری معنای برنامه‌ها، به شکلی که مورد بحث است، به آن‌ها نیاز است. در این بخش ابتدا به توضیحی مختصر در مورد برچسب‌ها در معناشناسی‌ زبان مورد بحث می‌پردازیم. تعاریف صوری دقیق این موجودات در پیوست \cite{calcul} آورده‌شده‌اند. از آوردن مستقیم این تعاریف در اینجا خود‌داری کرده‌ایم. البته در مورد معنای صوری برجسب‌ها هم ذکر این نکته ضروری است که نویسنده‌ی \cite{calcul} حتی به صورت صوری هم برای هر بخش از برنامه این کار را به طور دقیق انجام نداده و انجام این کار به طور دقیق‌تر را احتمالا به کسی که قرار است یک پیاده سازی کامل از این روش داشته باشد سپرده.


در زبانمان \lr{$\mathsf{S}$}ها بخشی از موجودات موجود در زبان هستند. برچسب ها را برای \lr{$\mathsf{S}$}ها تعریف می‌کنیم. برچسب‌ها با کمک توابع \lr{labs, in, brks-of, brk-to, esc, aft, at} تعریف می‌شوند. در‌واقع هر $\mathsf{S}$ به ازای بعضی از این توابع یک برجسب دارد و این‌ها در‌واقع نشان‌دهنده‌ی آن برچسب هستند. بعض دیگر این توابع برای هر $\mathsf{S}$ ممکن است یک مجموعه از برچسب‌ها را تعیین‌ کند و یکی از آن‌ها هم با گرفتن $\mathsf{S}$ یک مقدار بولی را بر‌می‌گرداند. 
\\\\
\lr{at[$\mathsf{S}$]} : برچسب شروع $\mathsf{S}$.\\
\lr{aft[$\mathsf{S}$]} : برچسب پایان $\mathsf{S}$، اگر پایانی داشته باشد.\\
\lr{esc[$\mathsf{S}$]} : یک مقدار بولی را باز‌‌می‌گرداند که بسته به اینکه در $\mathsf{S}$ دستور $\mathsf{break;}$ وجود دارد یا خیر، مقدار درست یا غلط را بر‌می‌گرداند.\\
\lr{brk-to[$\mathsf{S}$]} : برچسبی است که اگر حین $\mathsf{S}$ دستور $\mathsf{break;}$ اجرا شود، برنامه از آن نقطه ادامه پیدا می کند.\\
\lr{brks-of[$\mathsf{S}$]} : مجموعه‌ای از برچسب $\mathsf{break;}$ های $\mathsf{S}$ را بر‌می‌گرداند.\\
\lr{in[$\mathsf{S}$]} : مجموعه‌ای از تمام برچسب‌های درون $\mathsf{S}$ را برمی‌گرداند.\\
\lr{labs[$\mathsf{S}$]} : مجموعه‌ای از تمام بر‌چسب‌هایی که با اجرای $\mathsf{S}$ قابل دسترسی هستند را بر‌می‌گرداند.\\\\\\


\subsection{رد پیشوندی}


پس از تعریف برچسب‌ها به سراغ تعریف رد پیشوندی می‌رویم. پیش از آن باید وضعیت‌ها و محیط‌ها را تعریف کنیم.
\begin{defn}
	(محیط): به ازای مجموعه مقادیر $\mathbb{V}$ و مجموعه متغیرهای $\mathbb{X}$ تابع 
	$\rho : \mathbb{X} \rightarrow \mathbb{V}$ 
	را یک محیط می‌گوییم. مجموعه‌ی همه‌ی محیط‌ها را با $\mathbb{EV}$ نمایش می‌دهیم.
\end{defn}

\begin{defn}
	(وضعیت): به هر زوج مرتبْ به ترتیب متشکل از یک برچسب $l$ و یک محیط $\rho$ یک وضعیت (یا حالت)  
	$\langle l , \rho \rangle$
	می‌گوییم. مجموعه‌ی همه‌ی وضعیت‌ها را با $\mathfrak{S}$ نشان می‌دهیم.
\end{defn}
\begin{defn}
	(رد پیشوندی): به یک دنباله از وضعیت‌ها(با امکان تهی بودن) یک رد پیشوندی می‌گوییم.
\end{defn}

 هر رد پیشوندی یک دنباله است که قرار است توصیفی از چگونگی اجرای برنامه باشد. وضعیت‌ها همان‌طور که از نامشان پیداست قرار‌ است موقعیت لحظه‌ای برنامه را توصیف کنند. $l$ قرار است برچسب برنامه‌ی در حال اجرا باشد و $\rho$ مقدار متغیر‌ها را در آن موقع از اجرای برنامه نشان می‌دهد. دنباله‌های ما می‌توانند متناهی یا نامتناهی باشند. مجموعه‌ی ردهای پیشوندی‌ متناهی را با $\mathfrak{S^+}$ و مجموعه‌ی ردهای پیشوندی نامتناهی را با  $\mathfrak{S^\infty}$ نمایش می‌دهیم. مجموعه‌ی همه‌ی ردهای پیشوندی را هم با $\mathfrak{S^{+\infty}}$ نمایش می‌دهیم. 
با‌توجه به آنچه گفتیم، یک عملگر چسباندن $\Join$ را روی ردهای پیشوندی تعریف می‌کنیم. 
\begin{defn}
(عملگر چسباندن): اگر داشته باشیم 
$\pi_1 , \pi_2 \in \mathfrak{S^{+\infty}}  , \sigma_1 ,\sigma_2 \in \mathfrak{S}$
داریم:\\
\begin{center}
	اگر $\pi_1 \in \mathfrak{S^+} $ داریم  $\hspace{2.60cm}  $                                     
	$\pi_1 \Join \pi_2 = \pi_1$    \\
	اگر $\sigma_1\neq\sigma_2$    $\hspace{1.6cm}  $
	$\pi_1 \Join \pi_2$ تعریف نشده است
	\\اگر $\pi_1 \in \mathfrak{S^\infty} $ داریم   $\hspace{1cm}  $ 
	$\pi_1 \sigma_1 \Join \sigma_1 \pi_2 = \pi_1 \sigma_1 \pi_2 $
	
\end{center}
\end{defn}
همینطور $\epsilon$ هم یک رد پیشوندی است که حاوی هیچ وضعیتی نیست. به عبارت دیگر یک دنباله‌ی تهی است.

\subsection{تعریف صوری معناشناسی رد پیشوندی}
در این بخش قرار است دو تابع $\mathcal{A}$ و $\mathcal{B}$ را به ترتیب روی عبارات حسابی و بولی زبانمان یعنی $\mathsf{A}$ها و $\mathsf{B}$ها تعریف کنیم سپس با کمک آنها $\mathcal{S^*}$ را روی  مجموعه‌ای از اجتماع معنای $\mathsf{S}$ها و $\mathsf{Sl}$ها تعریف می کنیم. پس در نهایت هدف ما تعریف  $\mathcal{S^*}$ است.


\begin{defn}
(معنای عبارات حسابی - تابع $\mathcal{A}$): تابع 
$\mathcal{A}:\mathbb{A}\rightarrow \mathbb{EV} \rightarrow \mathbb{V}$
را به صورت بازگشتی روی ساختار 
$\mathsf{A} \in \mathbb{A}$
به شکل زیر تعریف می‌کنیم:
$$\mathcal{A[\mathsf{1}]\rho = }1     $$
$$\mathcal{A[\mathsf{x}]\rho = } \rho(\mathsf{x})          $$
$$\mathcal{A[\mathsf{A_1-A_2}]\rho = }\mathcal{A[\mathsf{A_1}]\rho }- \mathcal{A[\mathsf{A_2}]\rho }       $$

\end{defn}

\begin{defn}
	(معنای عبارات بولی - تابع $\mathcal{B}$): تابع 
	$\mathcal{B}: \mathbb{B} \rightarrow \mathbb{EV} \rightarrow \mathbb{BOOL}$
	را به صورت بازگشتی روی ساختار 
	$\mathsf{B} \in \mathbb{B}$
	به شکل زیر تعریف می‌کنیم:

\begin{center}
	اگر $\mathcal{A[\mathsf{A_1}]\rho }$ کوچکتر از $\mathcal{A[\mathsf{A_2}]\rho }$ باشد
	$\mathcal{B[\mathsf{A_1<A_2}]\rho = } True   \hspace{2cm}  $\\
	اگر $\mathcal{A[\mathsf{A_1}]\rho }$ بزرگتر از $\mathcal{A[\mathsf{A_2}]\rho }$ باشد
	$\mathcal{B[\mathsf{A_1<A_2}]\rho = } False   \hspace{2cm}  $\\
	$ \mathcal{B[\mathsf{B_1 nand B_2}]\rho = } \neg(\mathcal{B[\mathsf{B_1}]\rho}   \wedge \mathcal{B[\mathsf{B_2}]\rho}) $
\end{center}
\end{defn}

طبعا $\wedge$ و $\neg$ در فرازبان هستند.\\\\
در ادامه به تعریف $\mathcal{S^*}$ می‌پردازیم. این کار را با تعریف $\mathcal{S^*}$ روی هر ساخت $\mathsf{S}$ و $\mathsf{Sl}$ انجام می‌دهیم.
پیش از ادامه‌ی بحث باید این نکته را در‌مورد علامت‌گذاری‌هایمان ذکر کنیم که منظور از $        \mathsf{S} ::= l \mathsf{break;}  $ این است که تاکید کرده‌ایم که $\mathsf{S}$ با برچسب $l$ شروع شده‌است وگرنه همین طور که گفتیم   $l$ جزو زبان نیست.\\\\

\begin{defn}
	(معنای برنامه‌ها - تابع $\mathcal{S}^*$): 
اگر $        \mathsf{S} ::= \mathsf{break;}  $ باشد، ردهای پیشوندی متناظر با اجرای این دستور را به شکل مجموعه‌ی زیر تعریف می کنیم:
$$\mathcal{S^*} [\mathsf{S}] = \{ \langle at[\mathsf{S}] , \rho \rangle | \rho \in \mathbb{EV}       \} \cup     \{ \langle at[\mathsf{S}] , \rho \rangle \langle brk-to[\mathsf{S}] , \rho \rangle | \rho \in \mathbb{EV}       \}             $$   


اگر $        \mathsf{S} ::=  \mathsf{x\doteq A;}  $ باشد، ردهای پیشوندی متناظر با اجرای این دستور را به شکل مجموعه‌ی زیر تعریف می کنیم:
$$\mathcal{S^*} [\mathsf{S}] = \{ \langle at[\mathsf{S}] , \rho \rangle | \rho \in \mathbb{EV}       \} \cup     \{ \langle at[\mathsf{S}] , \rho \rangle \langle aft[\mathsf{S}] , \rho[\mathsf{x}\leftarrow \mathcal{A}[\mathsf{A}]\rho] \rangle | \rho \in \mathbb{EV}       \}             $$   

اگر $         \mathsf{S} ::= \mathsf{if}  \mathsf{ (B) S_t}  $ باشد، ردهای پیشوندی متناظر با اجرای این دستور را به شکل مجموعه‌ی زیر تعریف می کنیم:
$$\mathcal{S^*} [\mathsf{S}] = \{ \langle at[\mathsf{S}] , \rho \rangle | \rho \in \mathbb{EV}       \} \cup     \{ \langle at[\mathsf{S}] , \rho \rangle \langle aft[\mathsf{S}] , \rho \rangle | \mathcal{B}[\mathsf{B}] \rho =False      \} 
$$$$\cup    \{ \langle at[\mathsf{S}] , \rho \rangle \langle at[\mathsf{S_t}] , \rho \rangle 
\pi | \mathcal{B}[\mathsf{B}] \rho =True  \wedge   \langle  at[\mathsf{S_t}]  , \rho \rangle \pi \in \mathcal{S} [\mathsf{S_t}]    \}          $$ 


اگر $         \mathsf{S} ::= \mathsf{if}  \mathsf{ (B) S_t else S_f}  $ باشد، ردهای پیشوندی متناظر با اجرای این دستور را به شکل مجموعه‌ی زیر تعریف می کنیم:
$$\mathcal{S} [\mathsf{S}] = \{ \langle at[\mathsf{S}] , \rho \rangle | \rho \in \mathbb{EV}       \} $$$$\cup     \{ \langle at[\mathsf{S}] , \rho \rangle \langle at[\mathsf{S_f}] , \rho \rangle 
\pi | \mathcal{B}[\mathsf{B}] \rho =False  \wedge   \langle  at[\mathsf{S_f}]  , \rho \rangle \pi \in \mathcal{S} [\mathsf{S_f}]    \}  
$$$$\cup    \{ \langle at[\mathsf{S}] , \rho \rangle \langle at[\mathsf{S_t}] , \rho \rangle 
\pi | \mathcal{B}[\mathsf{B}] \rho =True  \wedge   \langle  at[\mathsf{S_t}]  , \rho \rangle \pi \in \mathcal{S} [\mathsf{S_t}]    \}          $$ \\


اگر 
$         \mathsf{Sl} ::= \backepsilon  $
 باشد، ردهای پیشوندی متناظر با اجرای این دستور را به شکل مجموعه‌ی زیر تعریف می کنیم:

$$\mathcal{S} [\mathsf{Sl}] = \{ \langle at[\mathsf{Sl}] , \rho \rangle | \rho \in \mathbb{EV}       \}        $$ \\

اگر $         \mathsf{Sl} ::= \mathsf{Sl' \:\:\: S}  $ باشد، ردهای پیشوندی متناظر با اجرای این دستور را به شکل مجموعه‌ی زیر تعریف می کنیم:
$$\mathcal{S} [\mathsf{Sl}] = \mathcal{S} [\mathsf{Sl'}] \cup( \mathcal{S} [\mathsf{Sl'}]
\Join \mathcal{S} [\mathsf{S}] )      $$ \\



اگر $         \mathsf{S} ::= \mathsf{while (B)S_b }   $ باشد، ماجرا نسبت به حالات قبل اندکی پیچیده‌تر می‌شود. تابعی به اسم $\mathcal{F} $ را تعریف خواهیم‌کرد که در حقیت دو ورودی دارد. ورودی اول آن یک دستور حلقه است و ورودی دوم آن یک مجموعه. به عبارتی دیگر می‌توانیم بگوییم به ازای هر حلقه یک تابع $\mathcal{F} $  جداگانه تعریف می‌شود که مجموعه‌ای از ردهای پیشوندی را می گیرد و مجموعه‌ای دیگر از همین موجودات را بازمی‌گرداند. کاری که این تابع قرار است انجام دهد این است که انگار یک دور دستورات داخل حلقه را اجرا می کند و دنباله‌هایی جدید را از دنباله‌های قبلی می‌سازد. معنای یک حلقه را کوچکترین نقطه ثابت این تابع در نظر می‌گیریم. در ادامه تعریف $\mathcal{F} $ آمده. با دیدن تعریف می توان به دلیل این کار پی‌برد. آن نقطه‌ای که دیگر $\mathcal{F} $ روی آن اثر نمی‌کند یا حالتی است که در آن دیگر شرط حلقه برقرار نیست و اصولا قرار نیست دستورات داخل حلقه اجرا شوند که طبق تعریف $\mathcal{F} $  می‌توانیم ببینیم که $\mathcal{F} $  در این حالت چیزی به ردهای پیشوندی اضافه نمی‌کند. یا اینکه حلقه به دستور $\mathsf{break;}$ خورده که در آن صورت وضعیتی به ته ردهای پیشوندی اضافه می‌شود که برچسبش خارج از مجموعه برچسب دستورات حلقه است و همین اضافه کردن هر چیزی را به ته ردهای پیشوندی موجود، توسط $\mathcal{F} $  غیرممکن می‌کند. بنابراین نقطه ثابت مفهوم مناسبی است برای اینکه از آن در تعریف صوری معنای حلقه استفاده کنیم. علت اینکه کوچکترین نقطه ثابت را به عنوان معنای حلقه در نظر می‌گیریم هم این است که مطمئن هستیم کوچکترین نقطه ثابت، هر رد پیشوندی ای را در خود داشته باشد به معنای اجرای برنامه مرتبط است. برای درک بهتر این نکته می‌توان به این نکته توجه کرد که با اضافه کردن وضعیت‌هایی کاملا بی‌ربط به اجرای برنامه به ته رد‌های پیشوندی، که صرفا برچسب متفاوتی با آخرین وضعیت هر رد پیشوندی دارند، نقطه ثابت جدیدی ساخته ایم. پس اگر خودمان را محدود به انتخاب کوچکترین نقطه ثابت نکنیم، به توصیفات صوری خوبی از برنامه‌ها دست پیدا نخواهیم‌کرد. در مورد نقطه ثابت تنها این نکته باقی می‌ماند که اصلا از کجا می‌دانیم که چنین نقطه ثابتی وجود دارد که در این صورت باید گفت مجموعه‌هایی که از ردهای پیشوندی تشکیل می‌شوند با عملگر زیرمجموعه بودن یک مشبکه را تشکیل می‌دهند و بنا به قضیه تارسکی\cite{tarski} برای چنین موجودی نقطه ثابت وجود دارد.
تعاریف موجوداتی که درموردشان صحبت کردیم به این شکل است:

$$\mathcal{S} [\mathsf{S}] = lfp^{\subseteq}\: \mathcal{F[\mathsf{S}]}      $$ $$\mathcal{F} [\mathsf{S}] X= \{ \langle at[\mathsf{S}] , \rho \rangle | \rho \in \mathbb{EV}       \} \cup $$
$$  \{ \pi_2 \langle l ,\rho \rangle \langle aft[\mathsf{S}],\rho \rangle |  \pi_2 \langle l ,\rho \rangle \in X \wedge \mathcal{B}[\mathsf{B}]\rho=False \wedge l= at[\mathsf{S}]   \} \cup      $$
$$  \{ \pi_2 \langle l ,\rho \rangle \langle at[\mathsf{S_b}],\rho \rangle \pi_3 |  \pi_2 \langle l ,\rho \rangle \in X \wedge \mathcal{B}[\mathsf{B}]\rho=True \wedge$$$$  \langle at[\mathsf{S_b}],\rho \rangle \pi_3 \in  \mathcal{S} [\mathsf{S_b}]   \wedge   l= at[\mathsf{S}]  \}  $$\\

اگر $         \mathsf{S} ::=;  $ باشد، ردهای پیشوندی متناظر با اجرای این دستور را به شکل مجموعه‌ی زیر تعریف می کنیم:
$$\mathcal{S} [\mathsf{S}] = \{ \langle at[\mathsf{S}] , \rho \rangle | \rho \in \mathbb{EV}       \} \cup     \{ \langle at[\mathsf{S}] , \rho \rangle \langle aft[\mathsf{S}] , \rho \rangle | \rho \in \mathbb{EV}       \}             $$  


اگر $         \mathsf{S} ::=\{\mathsf{Sl}\}  $ باشد، ردهای پیشوندی متناظر با اجرای این دستور را به شکل مجموعه‌ی زیر تعریف می کنیم:
$$\mathcal{S} [\mathsf{S}] = \mathcal{S} [\mathsf{Sl}] $$   \\
\end{defn}
در ادامه چند مثال می‌زنیم تا مفهوم موجودات تعریف شده مشخص‌تر شود.
\begin{exm}
	
\end{exm}
\begin{exm}
	
\end{exm}
\begin{exm}
	
\end{exm}

\section{ویژگی‌های معنایی برنامه‌ها}
تا به اینجای کار یک زبان آورده‌ایم و برای آن معنا تعریف کرده‌ایم. در این فصل می‌خواهیم در مورد ویژگی‌های برنامه‌هایی که در این زبان نوشته می‌شوند با توجه به معنای صوری‌ای که تعریف کرده‌ایم، صحبت کنیم. دقت شود که برای برنامه‌هایی که در یک زبان برنامه‌نویسی نوشته می‌شوند می‌توان به اشکال مختلفی ویژگی تعریف کرد؛ مثلا ویژگی‌های نحوی، مثل اینکه طول برنامه چند خط است یا هر کاراکتر چند بار به کار رفته، یا ویژگی‌های محاسباتی، مثل بررسی‌ سرعت برنامه یا میزان استفاده‌ی آن از حافظه که عموما در نظریه الگوریتم و پیچیدگی محاسبات بررسی می‌شود. منظور ما در اینجا از تعریف ویژگی، متناسب است با معناشناسی‌ای که برای برنامه‌هایمان تعریف کرده‌ایم. معناشناسی‌ای که تعریف کرده‌ایم درواقع سیر محاسباتی برنامه را توصیف می‌کند و ما می‌خواهیم ویژگی‌ها را با توجه به این موصوع تعریف کنیم. در این صورت می‌توانیم صحت عملکرد برنامه‌ها را با توجه به صادق بودن ویژگی‌هایی که در مورد آن‌ها تعریف شده بفهمیم.\\
ابتدا به تعریف ویژگی‌ها می‌پردازیم، سپس به سراغ تعریف یک نوع عبارت منظم می‌رویم که از آن برای بیان ویژگی‌ها استفاده می‌شود.
\subsection{ویژگی‌های معنایی}
همان‌طور که در بخش قبلی دیدیم، معنای هر برنامه با یک مجموعه‌ی 
$\mathcal{S^*} [\mathsf{S}]$
مشخص می‌شود. وقتی می‌خواهیم ویژگی‌هایی را برای موجوداتی که به کمک مجموعه‌ها تعریف شده اند بیان کنیم، اینکه ویژگی‌ها را هم با مجموعه‌ها بیان کنیم کار معقولی به نظر می‌رسد. مثل اینکه بخواهیم ویژگی زوج بودن را در مورد اعداد طبیعی بیان کنیم. می توانیم مجموعه‌ی $\mathbb{E}$ را به عنوان مجموعه‌ی همه‌ی اعداد زوج در نظر بگیریم و اینکه یک عدد زوج هست یا نه را عضویتش در مجموعه‌ی $\mathbb{E}$ تعریف کنیم. پس یعنی در مورد اعداد طبیعی قرار است هر ویژگی به شکل زیرمجموعه‌ای از تمام این اعداد در نظر گرفته شود. یعنی هر عضو 
$\mathit{P}(\mathbb{N})$
بنا به تعریف ما یک ویژگی از اعداد طبیعی است.
در مورد برنامه‌ها نیز قرار است همین رویه را پیش بگیریم. تابع 
$\mathcal{S^*}$
از نوع 
$\mathbb{P} \rightarrow \mathit{P}(\mathfrak{S^+})$
است. یعنی یک برنامه را در ورودی می‌گیرد و یک مجموعه از ردهای پیشوندی را باز می‌گرداند. پس می‌توانیم هر ویژگی را به عنوان زیر مجموعه‌ای از 
$\mathit{P}(\mathfrak{S^+})$
تعریف کنیم، به عبارت دیگر عضوی از
$\mathit{P(P}(\mathfrak{S^+}))$.

\subsection{عبارات منظم}
در اینجا توصیف ویژگی‌ها برای هر برنامه باید یک چارچوب داشته باشد. در صورت قدیمی روش وارسی مدل ما از منطق های زمانی برای بیان ویژگی‌ها به صورت صوری استفاده می‌کردیم و این احتیاج به یک زبان برای صوری کردن کامل کار را، که رسیدن به بیان مسئله‌ی وارسی مدل است، به ما نشان می‌دهد. در اینجا ما با داستان دیگری هم رو به رو هستیم و آن این است که از آنجایی که با مجموعه‌ها سر و کار داریم و مجموعه‌ها چندان موجودات ساختنی‌ای نیستند( برخلاف مدل کریپکی)، بهتر است یک موجود ساختنی مثل یک زبان صوری برای بیان آن‌ها داشته باشیم. در این فصل قصد داریم یک نوع عبارت منظم را برای این منظور تعریف کنیم. پیشتر به نکته‌ی دیگری در مورد استفاده از عبارات منظم، که متداول‌تر بودن بین جامعه‌ی برنامه نویسان است، صحبت کردیم. ابتدا زبان این عبارت منظم را تعریف می‌کنیم، سپس به سراغ معناشناسی آن می‌رویم. 
\subsection{زبان عبارات منظم}
فرق عمده‌ای که زبان عبارات منظم ما با عبارات منظم کلاسیک دارد در کاراکترهاست. کاراکترها در زبان کلاسیک موجوداتی اتمی بودند، اما در اینجا ساختار دارند. در اینجا به جای هر کاراکتر یک زوج متشکل از مجموعه‌ی $\mathsf{L}$ و عبارت بولی $\mathsf{B}$ تشکیل شده‌اند که این زوج را به شکل 
$\mathsf{L : B}$
در زبانمان نمایش می‌دهیم.\\ 
زبان ما به شکل BNF زیر است:
\begin{defn}
$$\mathsf{L} \in \mathit{P}(\mathbb{L})$$        
$$\mathsf{x,y,...} \in \mathbb{X}$$
$$\mathsf{\underline{x},\underline{y},...} \in \mathbb{\underline{X}}$$
$$\mathsf{B} \in \mathbb{B}$$
$$\mathsf{R} \in \mathbb{R}$$
\newpage
$$\mathsf{R} ::=\hspace{0.5cm} \varepsilon\hspace{4.2cm}$$
$$|\:\:\:\mathsf{L : B}\hspace{2.4cm}$$
$$|\:\:\:\mathsf{R_1 R_2}\:\:\:(or\:\mathsf{R_1 \bullet R_2 })$$
$$\:\:\:\:\:\:\:|\:\:\:\mathsf{R_1\:\mid\:R_2}\:\:\:(or\:\mathsf{R_1 + R_2 })$$
$$|\:\:\:\mathsf{R_1^*}\hspace{2.8cm}$$
$$|\:\:\:\mathsf{R_1^+}\hspace{2.7cm}$$
$$|\:\:\:(\mathsf{R_1})\hspace{2.47cm}$$
\end{defn}
همان طور که قابل مشاهده است در اینجا عملگرهای دوتایی چسباندن
$(\bullet)$
 و انتخاب
$(|)$
 را داریم، به همراه عملگرهای یگانی 
$^*$
و
$^+$.
در ادامه خواهیم دید که در فرازبان معنی عملگر یگانی 
$^+$
به وسیله‌ی عملگر یگانی دیگر قابل بیان است، هرچند که در زبانمان هم برای سهولت کار از بیان این عملگر اجتناب نشده. 
توجه شود که پرانتزها هم جزئی از زبان قرار داده شده‌اند.


همین‌طور در اینجا می‌خواهیم از تعدادی عبارات مخفف که در ادامه کارمان را راحت‌تر می‌کنند صحبت کنیم. منظور از زوج 
$\mathsf{? : B}$
همان 
$\mathbb{L}\mathsf{ : B}$
است. عبارت 
$l : \mathsf{B}$
به جای عبارت 
$\{l\} : \mathsf{B}$
به کار می‌رود و منظور از عبارت 
$\neg l \mathsf{: B}$
نیز عبارت 
$\mathbb{L}\setminus\{l\}:\mathsf{B}$
است.

با یک نگاه به دستور این زبان یک نکته‌ی چشمگیر برای ما، با توجه به موجوداتی که در بخش قبل تعریف کردیم، با نگاه به قواعد این زبان می‌تواند وجود یک مجموعه‌ی
$\mathbb{\underline{X}}$
در کنار 
$\mathbb{X}$
که از قبل داشتیم باشد. قرار است به ازای هر 
$\mathsf{x}\in\mathbb{X}$
یک 
$\mathsf{\underline{x}}\in\mathbb{\underline{X}}$
داشته‌باشیم. منظور از 
$\mathsf{\underline{x}}$
مقدار متغیر 
$\mathsf{x}$
در ابتدای هر برنامه است. این یعنی تابع
$\underline{\rho}: \mathbb{\underline{X}} \rightarrow \mathbb{V}$
که 
$\mathbb{V}$
مجموعه‌ی مقادیر متغیرهاست( در بخش قبل به این اشاره نشد اما خود
$\rho$
هایی که در بخش قبل داشتیم هم از نوع
$\mathbb{X} \rightarrow \mathbb{V}$
بود. با توجه به زبانمان و توضیحاتی که در گذشته دادیم، می‌توان در نظر گرفت که در اینجا 
$\mathbb{V}$
همان اعداد صحیح است). همان‌طور که پیش‌تر گفتیم برای اشاره به یک تابع
$\rho$
از کلمه‌ی "محیط" استفاده می‌شود. به همین منوال در ادامه برای اشاره‌ به 
$\underline{\rho}$
از "محیط اولیه" استفاده می‌کنیم. برای اشاره به مجموعه‌ی همه‌ی محیط‌های اولیه هم از نماد 
$\underline{\mathbb{EV}}$
استفاده می‌کنیم.
بقیه‌ی موجودات از جمله برچسب‌ها و عبارات بولی را هم که قبلا داشتیم.\\
در ادامه به بیان صوری معنای زبان بیان شده می‌پردازیم. پس از آن می‌توانیم با بررسی چند مثال، از اینکه معنای هر عبارت منظم چیست درکی شهودی به دست آوریم.\\

\subsection{معناشناسی عبارات منظم}
معنای عبارات منظم را با استفاده از تابع 
$\mathcal{S}^r$
نشان می‌دهیم. این تابع به این شکل تعریف می‌شود که در ورودی یک عبارت منظم 
$\mathsf{R}$
را می‌گیرد، سپس یک مجموعه از زوج مرتب‌های( یا همان‌طور که پیش‌تر نام‌گذاری کردیم "وضعیت‌ها"ی)
$\langle\underline{\rho} , \pi\rangle$
را که 
$\pi \in \mathbb{S^*}$
و 
$\underline{\rho} \in \underline{\mathbb{EV}}$
باز می‌گرداند. بنابراین این تابع از نوع
$\mathbb{R} \rightarrow \mathit{P} (\mathbb{\underline{EV} \times S^*})$
است. همین‌طور دقت شود که تا به حال از 
$\mathbb{S}^*$
صحبتی نکرده بودیم و فقط 
$\mathbb{S}^+$
را معرفی کرده بودیم. 
$\mathbb{S}^*$
نیز برابر است با
$\mathbb{S}^+ \cup \{\epsilon\}$ 
(به لحاظ معنایی همان عملگر
$^*$
است که در زبان عبارات منظمهم هست، مشهور به ستاره‌ی کلینی).

تعریف استقرایی تابع 
$\mathcal{S}^r$
به شکل زیر است:
\begin{defn}
	تابع 
	$\mathcal{S}^r:\mathbb{R}\rightarrow \mathit{P}(\mathbb{\underline{EV}\times S}^*)$
	به صورت استقرایی روی ساختار عبارت منظم $\mathsf{R}$ به صورت زیر تعریف می‌شود:
$$\mathcal{S}^r[\varepsilon]= \{ \langle \underline{\rho} , \epsilon \rangle | \underline{\rho} \in \underline{\mathbb{EV}}\}$$
[یعنی معنای عبارت منظم
$\varepsilon$
مجموعه‌ای شامل زوج مرتب‌هایی از محیط‌های اولیه‌ی مختلف در کنار رد پیشوندی تهی استفاده می‌کند.]
$$\mathcal{S}^r[\mathsf{L:B}] = \{\langle\underline{\rho},\langle l , \rho \rangle \rangle | l \in \mathsf{L} \land \mathcal{B}[\mathsf{B}] \underline{\rho},\rho \}$$
[این یعنی معنای عبارت 
$\mathcal{S}^r[\mathsf{L:B}]$
زوج مرتب‌هایی هستند که عضو اول آن‌ها محیط‌های اولیه مختلف هستند( مانند مورد قبلی و البته در موارد آتی!) و عضو دوم آن‌ها ردهای پیشوندی تک‌عضوی 
$\langle l , \rho \rangle$
هستند که در آن‌ها برچسب 
$l$
باید در
$\mathsf{L}$
که مجموعه‌ای از برچسب‌هاست حضور داشته باشد و عبارت بولی 
$\mathsf{B}$
باید درباره‌ی محیط اولیه
$\underline{\rho}$
و محیط 
$\rho$
برقرار باشد. حتما متوجه این نکته شدید که 
$\mathcal{B}$
در اینجا به جای اینکه از نوع 
$\mathbb{EV} \rightarrow \mathbb{BOOL}$
باشد،همان‌طور که قبلا تعریف کردیم، از نوع
$\underline{\mathbb{EV}} \rightarrow \mathbb{EV} \rightarrow \mathbb{BOOL}$
است.( منظور از 
$\mathbb{BOOL}$
همان مجموعه‌ی 
$\{True,False\}$
است.) در اینجا 
$\mathcal{A}$
و 
$\mathcal{B}$
را در ادامه با نوع‌های متفاوت دوباره تعریف خواهیم کرد، که البته فرق اساسی‌ای با تعریف قبلی ندارد و صرفا گسترشی ساده از آن است.]
$$\mathcal{S}^r[\mathsf{R_1 R_2}]= \mathcal{S}^r[\mathsf{R_1}] \Join \mathcal{S}^r[\mathsf{R_2}]$$
\begin{center}
	به‌طوری که در آن برای هر دو مجموعه‌ی 
	$\mathcal{S}$
	و
	$\mathcal{S'}$
	از رد‌های پیشوندی:
$$\mathcal{S \Join S'}=
\{
\langle \underline{\rho}, \pi \pi' \rangle |
\langle \underline{\rho},\pi \rangle \in \mathcal{S} \land
\langle \underline{\rho},\pi' \rangle \in \mathcal{S'}
\}
$$
\end{center}
[این یعنی اگر یک عبارت منظم داشته باشیم که از چسباندن 
$\mathsf{R_1}$
و 
$\mathsf{R_2}$
به هم ساخته شده باشد، آنگاه معنای این عبارت منظم با چسباندن ردهای پیشوندی موجود در مولفه‌ی دوم زوج مرتب‌هایی که اعضای مجموعه‌ی معنای این دو عبارت منظم هستند و گذاشتن این رد پیشوندی‌های حاصل از چسباندنْ در معنای عبارت منظم جدید تعریف می‌شود. همین‌طور که می‌بینید یک عملگر چسباندن برای دو مجموعه از این زوج‌های 
$\langle \underline{\rho} , \pi \rangle$
تعریف شده و در تعریف 
$\mathcal{S}^r [\mathsf{R_1 R_2}]$
از آن کمک گرفته شده.\\
تا این تکه از تعریف معنای عبارت منظم که رسیده‌ایم، تا حدی به دستیابی به درکی شهودی از اینکه به چه نحوی قرار است عبارات منظم راهی برای توصیف ویژگی‌ در مورد برنامه‌ها باشد نزدیک‌تر شده‌ایم. همان‌طور که در مورد قبل دیدیم هر زوج 
$\mathsf{L:B}$
دقیقا به یک وضعیتْ داخل یک رد پیشوندی اشاره می‌کند. انگار که قرار است این زوج‌ها موازی با وضعیت‌ها در ردهای پیشوندی موجود در معنای یک برنامه جلو روند و منطبق باشند تا وارسی مدل انجام شود. درک این موضوع اولین قدم ماست در دیدن عصاره‌ی روش وارسی مدل در ادبیاتی که از اول این فصلْ عَلَم کرده‌ایم.]
$$\mathcal{S}^r [\mathsf{R_1\:|\:R_2}]= 
\mathcal{S}^r [R_1] \cup
\mathcal{S}^r [R_2]$$
[این موردْ معنای اعمال عملگر انتخاب روی دو عبارت منظم را توصیف می‌کند. معنای اعمال این عملگر به‌سادگی به صورت اجتماع معنای هر دو عبارت منظم تعریف شده.]

$$\mathcal{S}^r [\mathsf{R}]^0 = \mathcal{S}^r[\varepsilon]$$
$$\mathcal{S}^r [\mathsf{R}]^{n+1} = \mathcal{S}^r [\mathsf{R}]^{n} \Join
\mathcal{S}^r [\mathsf{R}]$$
[دو عبارت اخیر برای توصیف معنای عملگرهای $^*$ و $^+$ تعریف شده‌اند. عملگر $\Join$ و معنای 
$\mathcal{S}^r[\varepsilon]$
را هم که قبلا تعریف کرده بودیم و $0$ و $n$ و $n+1$ هم اعداد طبیعی‌اند و $+$ لاجرم همان جمع اعداد طبیعی است.]
$$\mathcal{S}^r[\mathsf{R^*}] =  \bigcup_{n \in \mathbb{N}}
\mathcal{S}^r [\mathsf{R}^n]$$
$$\mathcal{S}^r[\mathsf{R^+}] =  \bigcup_{n \in \mathbb{N}\setminus\{0\}}
\mathcal{S}^r [\mathsf{R}^n]$$
[این دو عبارت هم تعریف معنای خود دو عملگر $^*$ و $^+$ هستند. منظور از $\mathbb{N}$ مجموعه‌ی اعداد طبیعی است. همان‌طور که قبل‌تر هم اشاره شد $^+$ را می‌توان در فرازبان با $^*$ تعریف کرد. اضافه می‌کنیم که خود $^*$ را هم در فرازبان می‌توان با عملگر انتخاب تعریف کرد و در اینجا می‌توان این نکته را هم دید.]

$$\mathcal{S}^r [(\mathsf{B})]=\mathcal{S}^r [\mathsf{B}]$$
[این تکه از تعریف هم صرفا بیان می‌کند که پرانتزها تاثیری در معنای عبارات منظم ندارند که کاملا قابل انتظار است چرا که وجود پرانتز قرار است در صرفا در خواص نحوی زبان اثر بگذارد.]
\end{defn}
تعریف معنای عبارات منظم در اینجا تمام می‌شود اما همان‌گونه که در لا‌به‌لای تعاریف گفتیم، احتیاج داریم که $\mathcal{A}$ و $\mathcal{B}$ را از نو تعریف کنیم:
\begin{defn}
	توابع 
	$\mathcal{A}:\mathbb{A}\rightarrow \mathbb{\underline{EV}} \rightarrow
	\mathbb{EV} \rightarrow \mathbb{V}$
	و 
	$\mathcal{B}:\mathbb{B}\rightarrow \mathbb{\underline{EV}} \rightarrow
	\mathbb{EV} \rightarrow \mathbb{BOOL}$
	به شکل استقرایی به ترتیب روی ساختارهای 
	$\mathsf{A}\in\mathbb{A}$
	و
	$\mathsf{B} \in \mathbb{B}$
	به شکل زیر تعریف می‌شوند:
	
	
$$\mathcal{A}[\mathsf{1}]\underline{\rho},\rho=1$$
$$\mathcal{A}[\mathsf{\underline{\mathsf{x}}}]\underline{\rho},\rho= \underline{\rho}(\mathsf{x})$$
$$\mathcal{A}[\mathsf{\mathsf{x}}]\underline{\rho},\rho= \rho(\mathsf{x})$$
$$\mathcal{A}[\mathsf{A_1 - A_2}]\underline{\rho},\rho= 
\mathcal{A}[\mathsf{A_1}]\underline{\rho},\rho - \mathcal{A}[\mathsf{A_2}]\underline{\rho},\rho$$
$$\mathcal{B}[\mathsf{A_1 < A_2}]\underline{\rho},\rho=
	\mathcal{A}[\mathsf{A_1}]\underline{\rho},\rho\: <\: \mathcal{A}[\mathsf{A_2}]\underline{\rho},\rho$$
$$\mathcal{B}[\mathsf{B_1\:nand\:B_2}]\underline{\rho},\rho=
\mathcal{B}[\mathsf{B_1}]\underline{\rho},\rho \uparrow 
\mathcal{B}[\mathsf{B_2}]\underline{\rho},\rho $$
\end{defn}
به‌راحتی قابل مشاهده است که تعاریف جدید تا حد خوبی به تعاریف قبلی شبیه هستند و فرق عمده صرفا وارد شدن $\underline{\rho}$ است.\\
حال به سراغ چند مثال از عبارات منظم و معنای آن‌ها می‌رویم.
\begin{exm}
	فرض کنید عبارت منظم ما

\end{exm}

\begin{exm}
	
	{این سه تا مثال باید با سه تا مثال بخش معناشناسی برنامه‌ها سینک باشن! همین طور در ادامه بعد از تعریف مدل چکینگ در این فصل هم ۳ تا مثال خواهیم داشت که قراره هر کدومشون از این ۳ تا مثال و ۳ تا مثالی که قبلا تعریف کردیم تشکیل شده باشن.}
\end{exm}

\begin{exm}
	
\end{exm}

تا اینجای کار بیشتر مفاهیمی که برای بیان صورت جدید مسئله‌ی وارسی مدل احتیاج داریم را بیان کرده‌ایم. 
\subsection{واریته‌های مختلف زبان عبارات منظم}
به عنوان قسمت آخر این بخشْ واریته‌های مختلفی از زبان عبارات منظم را بیان می‌کنیم. ، که هر کدام در واقع زیرمجموعه‌هایی از کل عبارات زبانی که توصیف کرده‌ایم را توصیف می‌کنند. بعضی از آن‌ها را در همین فصل برای هدف نهایی این فصل و بعضی دیگر را در فصل بعدی استفاده می‌کنیم.

اولین واریته‌ای که می‌خواهیم بیان کنیم، واریته‌ای است که در اعضای آن اصلا عبارت 
$\mathsf{L : B}$
حضور ندارد و کل عبارت‌های زبان از $\varepsilon$ ها تشکیل شده‌اند.
\begin{defn}
	(عبارت منظم تهی - $\mathbb{R_\varepsilon}$): 
	$$\mathsf{R} \:\:\:\in\:\:\: \mathbb{R}_\varepsilon$$
	$$\mathsf{R}\:\:\: ::= \:\:\: \varepsilon \: | \: \mathsf{R_1 R_2} \: | 
	\: \mathsf{R_1 + R_2} \: | \: \mathsf{R_1^*} \: | \: \mathsf{R_1^+} | (\mathsf{R_1})$$
\end{defn}


 با توجه به بخش قبل متوجه هستیم که معنای همه‌ی این عبارت‌ها برابر 
$\{\langle \underline{\rho} , \epsilon \rangle\}$
خواهد بود.

واریته‌ی بعدی عبارت منظم ناتهی است.
\begin{defn}
	(عبارت منظم ناتهی - $\mathbb{R}^+$):
	$$\mathsf{R} \:\:\:\in\:\:\: \mathbb{R}^+$$
	$$\mathsf{R}\:\:\: ::= \:\:\: \mathsf{L:B} \: |
	\: \mathsf{\varepsilon R_2} \: | \: \mathsf{R_1 \varepsilon} \: |
	\: \mathsf{R_1 R_2} \: | 
	\: \mathsf{R_1 + R_2} \: | \: \mathsf{R_1^+} | (\mathsf{R_1})$$
\end{defn}
دلیل وجود 
$\mathsf{\varepsilon R_2}$
و 
$\mathsf{R_1 \varepsilon}$
در تعریف این است که ممکن است معنای عبارتی با معنای عبارات عضو $\mathbb{R}_\varepsilon$ برابر نباشد( بعنی برابر ${\langle \underline{\rho} , \epsilon \rangle}$ نباشد)، اما در خود عبارتْ $\varepsilon$ حضور داشته باشد. با این تفاصیل می‌توان دید که دو مجموعه‌ی 
$\mathbb{R}_\varepsilon$
و
$\mathbb{R}^+$
یک افراز برای مجموعه‌ی $\mathbb{R}$ هستند، براساس اینکه معنای هر عبارت در $\mathbb{R}$ برابر ${\langle \underline{\rho} , \epsilon \rangle}$ هست یا خیر. بنابراین شاید به نظر برسد که تعریف یکی از آن‌ها به طور ساختاری کافی بود، اما ممکن است درجایی احتیاج داشته باشیم که ساختاری استقرایی روی هر یک از آن‌ها عَلَم کنیم یا اینکه در اثبات حکمی بخواهیم از استقرا روی یکی از این دو ساختار استفاده کنیم.

واریته‌ی آخر عبارات منظم ما نیز عبارات منظم بدون انتخاب است.
\begin{defn}
(عبارت منظم بدون انتخاب - $\mathbb{R}^\nmid$):
$$
\mathsf{R}\:\:\: \in \:\:\: \mathbb{R}^\nmid$$
$$
\mathsf{R}\:\:\: ::= \:\:\: \varepsilon \: | \: \mathsf{L:B} \: | \:
\mathsf{R_1 R_2} \: | \: \mathsf{R_1}^* \: | \: \mathsf{R_1}^+ \: | \:
(\mathsf{R_1})$$
\end{defn}

\section{صورت جدید مسئله‌ی وارسی مدل}
بالاخره به هدف نهایی این فصل رسیدیم. می‌خواهیم صورت جدیدی از مسئله‌ی وارسی مدل را بیان ‌کنیم.\\
پیش از ارائه‌ی تعریف وارسی مدل نیاز داریم تا عملگر بستار پیشوندی را برای یک مجموعه از ردهای پیشوندی معرفی کنیم.
\begin{defn}
	(بستار پیشوندی):
	اگر 
	$\Pi \in \mathit{P}(\mathbb{\underline{EV}\times S^+})$
	آنگاه بستار پیشوندی $\Pi$را به صورت زیر تعریف می‌کنیم:
	$$\mathsf{prefix}(\Pi)=
	\{\langle \underline{\rho},\pi\rangle | \pi \in \mathbb{S^+} \land \exists\
	\pi' \in \mathbb{S^*}: \langle \underline{\rho}, \pi \pi' \rangle \in \Pi \}$$
\end{defn}

برای درک بهتر مفهوم بستار پیشوندی به مثال زیر توجه شود.
\begin{exm}
	اگر 
	$\Pi = \{\langle\underline{\rho},\langle l_1 , \rho_1 \rangle 
	\langle l_2  \rho_2 \rangle\rangle  \langle l_3 , \rho_3 \rangle\rangle 
	,
	\langle\underline{\rho},\langle {l_1}' , {\rho_1}' \rangle  \langle {l_2}'  {\rho_2}' \rangle\rangle
	\}$ 
	باشد( $\Pi$ شامل دو عضو است) آنگاه:
	$$
	\mathsf{prefix}(\Pi)=
	\{
	\langle\underline{\rho},\langle l_1 , \rho_1 \rangle\rangle ,
	\langle\underline{\rho},\langle l_1 , \rho_1 \rangle  \langle l_2  \rho_2 \rangle\rangle,
	\langle\underline{\rho},\langle l_1 , \rho_1 \rangle  \langle l_2  \rho_2 \rangle  \langle l_3 , \rho_3 \rangle\rangle ,$$
	$$
	\langle\underline{\rho},\langle {l_1}' , {\rho_1}' \rangle\rangle,
	\langle\underline{\rho},\langle {l_1}' , {\rho_1}' \rangle  \langle {l_2}'  {\rho_2}' \rangle\rangle
	\}
	$$
	که شامل 5 عضو است.
\end{exm}
حال به ارائه‌ی صورت جدیدمان از روش وارسی مدل می‌رسیم که هدف اصلی این اصل بود و با این تعریف فصل تمام می‌شود.
\begin{defn}
(وارسی مدل):
	اگر 
	$\mathsf{P}\in\mathbb{P} , \mathsf{R} \in \mathbb{R}^+ , \underline{\rho} \in \underline{\mathbb{EV}}$
	آنگاه:
	$$\mathsf{P},\underline{\rho} \models \mathsf{R}
	 \Leftrightarrow
	(\{\underline{\rho}\}\times \mathcal{S}^* [\mathsf{P}]) \subseteq 
	\mathsf{prefix} (\mathcal{S}^r [\mathsf{R} \bullet (?:\mathit{T})^*])
	$$
\end{defn}

این تعریف بیان می‌کند که برنامه‌ی 
$\mathsf{P}$
در صورتی که با محیط اولیه‌ی 
$\underline{\rho}$
اجرا شود، در صورتی خاصیتی که با عبارت منظم 
$\mathsf{R}$
بیان شده را دارد که معنای آن زیرمجموعه‌ی بستار پیشوندی معنای عبارت منظم
$\mathsf{R} \bullet (?:\mathit{T})^*$
باشد.
توجه شود که محیط اولیه‌ای که برای برنامه‌ی مورد بررسی متصور هستیم صرفا به این منظور قرار داده شده که معناشناسی برنامه را بتوانیم با معنای عبارات منظم قابل قیاس کنیم. دلیل حضور محیط اولیه در معنای عبارات منظم نیز در صورت سوم روش وارسی مدل یعنی فصل ۴ مشخص می‌شود و در این صورت از روش وارسی مدل و صورت بعدی آن چندان نقشی ندارد. 

در مورد نقش 
$ (?:\mathit{T})^*$
و
$ \mathsf{prefix} $
این به تصمیم مبدع این روش بوده که این دو در تعریف روش وارسی مدل حضور داشته باشند. حضور این دو طبعا باعث می‌شود به ازای یک عبارت منظم 
$\mathsf{R}$
نسبت به این حالت که صرفا معنی برنامه زیرمجموعه‌ی معنی 
$\mathsf{R}$
باشد، برنامه‌های بیشتری باشند که خاصیت بیان شده با 
$\mathsf{R}$
را ارضا کنند، چون در این صورت مجموعه‌ی سمت راستی در رابطه‌ی زیرمجموعه بودن بزرگتر می‌شود.

\section{در مورد توقف پذیری}

در این بخش مشکلی از کار که به نظر نگارنده رسیده مطرح شده. متاسفانه این مشکل بسیار بزرگ است و به نظر عجیب می‌‌آید اگر نویسنده‌ی \cite{calcul} متوجه آن نبوده باشد! اگر صحبت ما در اینجا درست باشد، این به این معنی خواهد بود که کل کاری که در حال توصیفش هستیم قابل پیاده سازی نیست!

بحث ما در اینجا در مورد توقف پذیری است. در \cite{calcul} در مورد توقف یک برنامه صحبتی به میان نیامده. یعنی حتی گفته نشده که در چه صورتی می‌توانیم بگوییم که یک برنامه متوقف شده است. یک تعریف  صوری معقول که خودمان می‌توانیم برای این معنا بیاوریم این است:

\begin{defn}
(توقف پذیری:) برنامه‌ی $\mathsf{P}$ را به همراه اجرای اولیه $\underline{\rho}$ توقف پذیر می‌گوییم اگر و تنها اگر وجود داشته باشد 
	$\pi \in \mathcal{S}^* \llbracket \mathsf{P} \rrbracket$ 
	که ($\rho$ محیط متناظر با محیط اولیه‌ی $\underline{\rho}$ است.):
	$$\pi = \langle at \llbracket \mathsf{P} \rrbracket , \rho \rangle \pi'$$
	و اینکه $\langle aft\llbracket \mathsf{P} \rrbracket , \rho' \rangle$ در $\pi$ حضور داشته باشد. این اتفاق را با 
	$\mathsf{P,\underline{\rho}\downarrow}$
	نشان می‌دهیم. همین تعریف را برای لیست دستورات $\mathsf{Sl}$ یا دستور $\mathsf{S}$ هم صرفا با جایگذاری این‌ها با برنامه‌ی $\mathsf{P}$ داریم.
\end{defn} 
در این تعریف توقف پذیری صرفا برای یک محیط اولیه تعریف شده.
در اینجا توقف پذیری به متناهی بودن ردهای پیشوندی موجود در برنامه ربط داده نشده. با توجه به معناشناسی‌ای که داریم، تعریف توقف پذیری به معنای وجود رد پیشوندی متناهی با محیط اولیه‌ی مورد بررسی در معنای برنامه که اصلا جور در نمی‌آید، چون معناشناسی ما خاصیت پیشوندی بودن را دارد و مطمئن هستیم در معنای هر برنامه‌ای حتما یک رد پیشوندی متناهی با محیط اولیه‌ی مورد بررسی وجود دارد.

اگر هم بخواهیم تعریف توقف پذیری را وجودنداشتن ردهای پیشوندی نامتناهی با محیط اولیه‌ی مورد بررسی در معنای برنامه در نظر بگیریم در ابتدا به نظر می‌آید که به تعریف قوی‌تری نسبت به آنچه ارائه دادیم رسیده‌ایم. ما در اینجا سعی داریم تعریفی را ارائه کنیم که برای حرف‌هایی که در \cite{calcul} زده شده تا حد امکان مشکل درست نکند، که اگر دیدیم با این وجود مشکل وجود دارد مطمئن باشیم که اشتباه در \cite{calcul} است و نه حرف ما. پس سعی از ارائه‌ی این تعریف که به نظر از تعریف ارائه شده با کار ناسازگارتر می‌آید اجتناب می‌کنیم( در ادامه به بیان ناسازگاری پراخته شده) اما در قضیه‌ی بعدی می‌بینیم که تعریفی که ارائه کردیم با همان که بگوییم در برنامه رد پیشوندی نامتناهی وجود ندارد معادل است.
\begin{thm}
	برای برنامه‌ی $\mathsf{P}$ و محیط اولیه‌ی $\underline{\rho}$ داریم $\mathsf{P} , \underline{\rho} \downarrow $ اگر و تنها اگر با فرض اینکه $\rho$ محیط متناظر با محیط اولیه‌ی $\underline{\rho}$ است و
	$$\forall \pi \in \mathbb{S^{+\infty}} : 
	\langle at \llbracket \mathsf{P} \rrbracket , \rho \rangle \pi \in \mathcal{S^*} \llbracket \mathsf{P} \rrbracket \rightarrow
	\langle at \llbracket \mathsf{P} \rrbracket , \rho \rangle \pi \in \mathbb{R^+}$$
	 
\end{thm}
\begin{proof}
	$(\Rightarrow)$
	برای این قسمت باید ثابت کنیم که در معنای هر برنامه‌ای رد پیشوندی‌ای وجود دارد که با
	$\langle at\llbracket \mathsf{P} \rrbracket , \rho \rangle$
	شروع شده و به ازای یک محیط $\rho'$ به 
	$\langle aft\llbracket \mathsf{P} \rrbracket , \rho' \rangle$
	ختم شده. 
	در این اثبات از تعریف برچسب‌ها که در ضمیمه‌ی \cite{calcul} آمده استفاده شده.
	داریم 
	$\mathsf{P=Sl}$
	و 
	$aft \llbracket \mathsf{P} \rrbracket = aft \llbracket \mathsf{Sl} \rrbracket $
	
	حکم را با استقرا روی ساختار $\mathsf{Sl}$ ثابت می‌کنیم. 
	$$\blacktriangleright \mathsf{Sl = \backepsilon}:$$
	داریم:
	$$\mathcal{S^*} \llbracket \backepsilon \rrbracket = \{\langle at \llbracket \backepsilon \rrbracket , \dot{\rho} \rangle | \dot{\rho} \in \mathbb{EV}\}$$
	و طبق تعریف برچسب‌ها داریم:
	$$at \llbracket \backepsilon \rrbracket = aft \llbracket \backepsilon \rrbracket$$
	پس حکم برقرار است.
	$$\blacktriangleright \mathsf{Sl =Sl'\;S}:$$
	اینکه در معنای $\mathsf{Sl}$ دنباله‌ای شامل
	 $\langle aft \llbracket \mathsf{Sl} \rrbracket , \rho '\rangle$ وجود داشته باشد، به با توجه به تعاریفی که داریم به این وابسته است که در معنای $\mathsf{S}$ دنباله‌ای شامل $\langle aft \llbracket \mathsf{S} \rrbracket , \rho' \rangle$ وجود داشته باشد. برای اینکه این را ثابت کنیم هم باید همین حکم را روی ساختار $\mathsf{S}$ ثابت کنیم که در واقع بخش اصلی اثبات این سمت قضیه است. 
	$$\blacktriangleright\blacktriangleright \mathsf{S\;=\; x\doteq A;}:$$
	در این حالت با توجه به تعریف معنای $\mathsf{S}$ که قبل‌تر ارائه شد، دنباله‌ی 
	$$\langle at \llbracket \mathsf{S} \rrbracket ,\rho \rangle \langle aft\llbracket\mathsf{S}\rrbracket , \rho[\mathsf{x}\leftarrow \mathcal{A}\llbracket\mathsf{A}\rrbracket\rho] \rangle $$
	در معنای دستور به ازای هر $\rho$ وجود دارد که خب در هر صورت این شامل محیط متناظر با $\underline{\rho}$ هم می‌شود .
	$$\blacktriangleright\blacktriangleright \mathsf{S=\; ; \;\;}:$$	 
با توجه به معنای این دستوردنباله‌ی زیر در معنای این دستور وجود دارد.
$$\langle at \llbracket \mathsf{S} \rrbracket ,\rho \rangle \langle aft\llbracket\mathsf{S}\rrbracket , \rho \rangle $$

$$\blacktriangleright\blacktriangleright \mathsf{S=\; if\;(B)\;S_t}:$$
	در صورتی که 
	$\mathcal{B}\llbracket\mathsf{B}\rrbracket\rho=\mathit{T}$
	دنباله‌ی 
	$$\langle at \llbracket \mathsf{S} \rrbracket , \rho \rangle
	\langle at \llbracket \mathsf{S_t} \rrbracket , \rho \rangle \pi$$
	در مجموعه‌ی معنای این دستور حضور دارد در حالیکه 
	$\langle at \llbracket \mathsf{S_t} \rrbracket , \rho \rangle \pi$
	داخل معنای $\mathsf{S_t}$ است و طبق فرض استقرا می‌دانیم که برچسب آخرین موقعیت $\pi$ برابر است با $aft \llbracket \mathsf{S_t} \rrbracket$ که طبق تعاریف مربوط به برچسب‌ها 
	$aft \llbracket \mathsf{S_t} \rrbracket=aft \llbracket \mathsf{S} \rrbracket$.
	در صورتی که معنای عبارت بولی غلط باشد هم دنباله‌ی زیر در معنای دستور طبق تعریف موجود است.
	$$\langle at \llbracket \mathsf{S} \rrbracket ,\rho \rangle \langle aft\llbracket\mathsf{S}\rrbracket , \rho \rangle $$
	
$$\blacktriangleright\blacktriangleright \mathsf{S=\; if\;(B)\;S_t\;else\;S_f}:$$
مانند حالت قبل است منتها با این تفاوت که در صورتی که معنای عبارت بولی غلط باشد دنباله‌ی زیر در معنای دستور حضور دارد:
	$$\langle at \llbracket \mathsf{S} \rrbracket , \rho \rangle
\langle at \llbracket \mathsf{S_f} \rrbracket , \rho \rangle \pi$$
و تساوی 
$aft \llbracket \mathsf{S_t} \rrbracket=aft \llbracket \mathsf{S} \rrbracket=aft \llbracket \mathsf{S_f} \rrbracket$
هم طبق تعریف برچسب‌ها برقرار است.

$$\blacktriangleright\blacktriangleright \mathsf{S=\; while\;(B)\;S_t}:$$
در اثبات این سمت قضیه این حالت پیچیده ترین حالت است و در واقع تنها حالتی است که در اثبات آن به فرض قضیه احتیاج داریم! همان طور که پیشتر گفتیم معنای حلقه با استفاده از یک تابع تعریف می‌شود. معنای حلقه کوچکترین نقطه ثابت این تابع است، در حالیکه انگار این تابع وقتی روی یک مجموعه از ردهای پیشوندی اعمال شود، تاثیرات یک بار اجرای دستورات درون حلقه را روی ردهای پیشوندی درون مجموعه اعمال می‌کند.

طبق تعریف $\mathcal{F}$ مطمئن هستیم که رد پیشوندی‌ای که با محیط $\underline{\rho}$ شروع شود در مجموعه‌ی معنای $\mathsf{S}$ وجود دارد، چونکه به ازای هر محیط $\dot{\rho}$(نقطه به این خاطر است که با $\rho$ خاص موجود در فرض اشتباه گرفته نشود) حالت 
$\langle at \llbracket \mathsf{S} \rrbracket, \dot{\rho} \rangle$
 در هر اعمال تابع $\mathcal{F}$ روی هر مجموعه‌ی دلخواه وجود دارد. وقتی معنای $\mathsf{S}$ را به عنوان کوچک‌ترین نقطه ثابت $\mathcal{F}$ در نظر گرفته‌ایم پس مطمئن هستیم که آن مجموعه‌ای که کوچکترین نقطه ثابت است شامل رد پیشوندی 
 $\langle at \llbracket \mathsf{S} \rrbracket, \rho \rangle$
 است. این رد پیشوندی با اجرای $\mathcal{F}$ تحت تاثیر قرار می‌گیرد. اگر معنای $\mathsf{B}$ در یکی از اعمال های $\mathcal{F}$ غلط باشد، رد پیشوندی
$ \langle at \llbracket \mathsf{S} \rrbracket, \rho \rangle \pi \langle aft \llbracket \mathsf{S} \rrbracket , \rho' \rangle$ 
در معنای برنامه قرار خواهد گرفت و می‌توانیم بگوییم اجرای دستور با این محیط اولیه توقف پذیر است. می‌دانیم که طبق تعریف تابع به انتهای این رد پیشوندی چیزی اضافه نمی‌شود. از طرف دیگر هم با این محیط اولیه، با توجه به تعریف رد پیشوندی دیگری وجود ندارد که طولانی‌تر از رد پیشوندی مورد اشاره باشد. 

در حالت دیگر اگر فرض کنیم هیچ گاه به حالتی نمی‌رسیم که در آن معنای $\mathsf{B}$ غلط باشد هم با فرض مسئله به تناقض می‌خوریم، چون در آن صورت تابع $\mathcal{F}$ مدام به طول دنباله‌‌هایی که با محیط $\rho$ شروع می‌شوند می‌افزاید و این یک دنباله‌ی نامتناهی را خواهد ساخت. در صورتی که معنای $\mathsf{B}$ هیچ گاه صحیح نباشد، حداقل حالت 
$\langle at \llbracket \mathsf{S_t} \rrbracket , \rho'' \rangle$
به ته دنباله‌های پیشین اضافه خواهد شد و از این جهت مطمئن هستیم که دنباله‌ی نامتناهی گفته شده در معنای دستور حضور خواهد داشت. 

پس با این تفاصیل، این مورد هم ثابت می‌شود.

\break$$\blacktriangleright\blacktriangleright \mathsf{S=\; break;}:$$
در تعریف تابع $aft$ روی برچسب‌ها در \cite{calcul} این تعریف برای این دستور مشخص نیست! در \cite{cousotbook} که در مورد برچسب‌ها بحث شده، نویسنده‌ی \cite{calcul} گفته که در مورد آن بخش از تعاریف توابع مربوط به برچسب‌ها که تعریف نشده‌اند برداشت آزاد است و ما در اینجا سعی داریم معقول ترین برداشتی که نسبت به درکمان از این کار می‌توانیم داشته باشیم را بیان کنیم. مهم‌ترین چیزی که در مورد برچسب‌ها در مورد این دستور قرار است برقرار باشد این است که اگر این دستور بخشی از $\mathsf{S_t}$ در حلقه‌ی زیر باشد
$$\mathsf{S'=\; while\; (B)\; S_t}$$
در این صورت 
$aft \llbracket \mathsf{S'} \rrbracket = brk-to \llbracket \mathsf{S_t} \rrbracket $
را طبق تعریف داریم. انتظار می‌رود که\break 
$aft \llbracket \mathsf{break;} \rrbracket = aft \llbracket \mathsf{S'} \rrbracket$ 
باشد. اینکه دستورات برنامه پس از اجرای $\mathsf{break;}$ از خارج(یا به عبارت بهتر بعد از) حقله‌ی $\mathsf{S'}$ پی گرفته شود انتظار معقولی است از سیستمی که در حال توصیف رد اجرای برنامه‌های کامپیوتری است. البته در نظر گرفته شود که فرض کرده‌ایم که $\mathsf{S'}$ داخلی ترین حلقه‌ای است که $\mathsf{break;}$ درون آن جای دارد.

 از پس این فرض‌های ما 
$aft \llbracket \mathsf{break;} \rrbracket = break-to \llbracket \mathsf{S_t} \rrbracket$ 
نتیجه می‌شود و طبق تعریف معنای دستورات 
$\mathsf{break;}$
رد پیشوندی زیر در معنای این دستور وجود دارد
$$\langle at \llbracket \mathsf{break;} \rrbracket , \rho \rangle
\langle aft \llbracket \mathsf{break;} \rrbracket , \rho \rangle$$
 که نشانه‌ی توقف است.
 

$$\blacktriangleright\blacktriangleright \mathsf{S=\; \{ Sl'' \}}:$$
	در این صورت توقف پذیری $\mathsf{Sl''}$ را از فرض استقرای استقرایی که روی لیست دستورات زده بودیم داریم پس $\mathsf{\{Sl''\}}$ هم توقف پذیر است.
	
	
در اینجا اثبات این طرف قضیه به پایان می‌رسد.

($\Leftarrow$)
دوباره باید روی ساختار برنامه‌ها استقرا بزنیم و دوباره چون هر برنامه مساوی با یک لیست از دستورات است استقرا را ابتدا روی ساختار لیست دستورات و در دل آن روی ساختار دستورات استقرا می‌زنیم.

در این اثبات به غیر از یک حالت ساختار دستور، که دستور حلقه است، هر آنچه در مورد اثبات طرف راست قضیه گفتیم، به ما حکم را بدون نیاز به فرض نشان می‌دهد. بنابراین فقط در مورد اثبات همین یک مورد بحث می‌کنیم. 
$$\blacktriangleright\blacktriangleright \mathsf{S=\; while\; (B)\; S_t}:$$

اگر فرض کنیم این دستور به ازای محیط $\rho$ در حالت اول متوقف شده، در واقع فرض کرده‌ایم در معنای این دستور رد پیشوندی 
$\langle at \llbracket \mathsf{S} \rrbracket , \rho \rangle \pi \langle aft \llbracket \mathsf{S} \rrbracket , \rho' \rangle \pi'$ 
وجود دارد.
باید ثابت کنیم به ازای $\pi'$ دلخواه اگر رد پیشوندی
$\langle at \llbracket \mathsf{S} \rrbracket , \rho \rangle \pi \langle aft \llbracket \mathsf{S} \rrbracket , \rho' \rangle \pi'$ 
داخل 
$\mathcal{S^*}\llbracket \mathsf{S} \rrbracket$
وجود داشته باشد آنگاه
$\pi'=\epsilon$
برقرار است.

اگر برچسب
$aft \llbracket \mathsf{S} \rrbracket$
در یک حالت در رد پیشوندی‌ای که گفتیم حضور داشته باشد، یعنی در یک دور اجرای حلقه عبارت بولی معنی غلط می‌داده که حالتی شامل این برچسب به یک رد پیشوندی چسبانده شده و این رد پیشوندی ساخته شده. از طرفی دیگر هم می‌دانیم که وقتی عبارت بولی حاضر در ساختار حلقه غلط شده، دیگر به ردهای پیشوندی داخل معنای حلقه چیزی اضافه نمی‌شود. بنابراین سناریو‌ای جز $\pi'=\epsilon$ باقی نمی‌ماند. 
\end{proof}

پس با توجه به آنچه گفتیم می‌توانیم با خیال راحت توقف پذیری یک برنامه با یک محیط اولیه را معادل متناهی بودن همه‌ی ردهای پیشوندی‌ای بدانیم که با محیط متناظر با آن محیط اولیه شروع شده‌اند.
اگر در صورت ارائه شده از وارسی مدل عبارت منظم $\mathsf{R}$ را با عبارت منظم $\varepsilon$ جایگزین کنیم داریم:
$$\mathsf{P},\underline{\rho} \models \mathsf{R}
\Leftrightarrow
(\{\underline{\rho}\}\times \mathcal{S}^* \llbracket \mathsf{P}\rrbracket) \subseteq 
\mathsf{prefix} (\mathcal{S}^r \llbracket \varepsilon \bullet (?:\mathit{T})^*\rrbracket)
=\mathsf{prefix} (\mathcal{S}^r \llbracket (?:\mathit{T})^*\rrbracket)$$
طبق تعریف معنای عبارات منظم، هر رد پیشوندی متناهی‌ای داخل مجموعه‌ی سمت راستی رابطه‌ی زیرمجموعه بودن قرار می‌گیرد. این یعنی اگر الگوریتمی برای بررسی 
$\mathsf{P} , \underline{\rho} \models \mathsf{R} $
داشته باشیم، این الگوریتم می‌تواند تشخیص دهد آیا برنامه‌ی $\mathsf{P}$ با محیط اولیه‌ی $\underline{\rho}$ متوقف می‌شود یا خیر! این یعنی الگوریتمی برای مسئله‌ی توقف پذیری، مسئله‌ای که تصمیم ناپذیر است! بنابراین چنین الگوریتمی نباید وجود داشته باشد که یعنی پیاده سازی‌ای برای شیوه‌ای که در حال بیانش هستیم وجود ندارد! ادامه‌ی کار روی همین تعریف پیش می‌رود و دو صورت دیگر هم که قرار است ساختارمندتر باشند در نهایت با این صورت معادل‌اند، هرچند که پس از رسیدن به بیان دو صورت دیگر خواهیم دید که همین صحبت‌هایی که در مورد این صورت می‌کنیم در مورد صورت‌های دیگر هم بدون در نظر گرفتن معادل بودن این ۳ صورت برقرار است.




\chapter{صوری‌سازی جدید برای روش وارسی مدل}

در این فصل، صورت جدیدی برای روش وارسی مدل ارائه ‌می‌شود. در این صورت، برای بیان ویژگی‌های مورد بررسی به‌جای منطق زمانی از \gls*{lregex} استفاده می‌شود. با صحبت در مورد ویژگیهای برنامه‌ها شروع می‌کنیم، سپس به معرفی نحو و معناشناسی عبارات منظم، به عنوان یک منطق برای بیان ویژگی‌ها، می‌پردازیم و پس از آن، صورت جدید روش وارسی مدل را ارائه می‌کنیم. در انتها، به بحث در مورد \gls*{decidability} این روش می‌پردازیم.


\section{ویژگی‌های برنامه‌ها}
تا به اینجای کار، نحو و معناشناسی یک زبان برنامه نویسی را تعریف کرده‌ایم. در این بخش، در مورد ویژگی‌های برنامه‌هایی که در این زبان نوشته می‌شوند، با توجه به معناشناسی رد پیشوندی که در فصل پیش تعریف کردیم، صحبت می‌کنیم. برای برنامه‌هایی که در یک زبان برنامه‌نویسی نوشته می‌شوند، می‌توان به اشکال مختلفی ویژگی تعریف کرد. مثلا ویژگی‌های کاملا نحوی، مثل اینکه طول برنامه چند خط است یا هر کاراکتر چند بار به کار رفته است، یا ویژگی‌های محاسباتی، مثل بررسی‌ سرعت برنامه یا میزان استفاده‌ی آن از حافظه که عموما در \gls*{algtheory} و \gls*{complexity} بررسی می‌شود. در اینجا، منظور ما از تعریف ویژگی، متناسب است با گردش کار برنامه‌ها. 
به کمک عبارات منظم که در این فصل تعریف خواهد شد; ویژگی‌های مورد نظر قابل بیان هستند.
\subsection{ویژگی‌ها}
همان‌طور که در فصل قبل دیدیم، معنای هر برنامه‌ی $\mathsf{P}$ با یک مجموعه‌ی 
$\mathcal{S^*}\llbracket {\mathsf{P}}\rrbracket$
مشخص می‌شود. درواقع،  
$\mathcal{S^*}$
تابعی
از  
$\mathbb{P}$
(مجموعه‌ی تمام برنامه‌ها) به
$\mathit{P}(\mathfrak{S^*})$
است. بنابراین، یک برنامه را در ورودی می‌گیرد و یک مجموعه از ردهای پیشوندی را باز می‌گرداند. پس هر ویژگی زیر مجموعه‌ای از 
$\mathit{P}(\mathfrak{S^*})$
است، به عبارت دیگر عضوی از
$\mathit{P(P}(\mathfrak{S^*}))$. 
تعریف ویژگی‌ها به کمک رابطه‌ی زیرمجموعه بودن از سنت‌های نظریه‌ی \gls*{abstint}\cite{cousot1} است.

\subsection{عبارات منظم}
برای بیان ویژگی‌های برنامه‌ها به یک چارچوب نیاز داریم. در صورت سنتی روش وارسی مدل از منطق های زمانی برای بیان ویژگی‌ها به صورت صوری استفاده می‌کرد و این احتیاج به یک منطق را برای ارائه‌ی یک \gls*{formalization} کامل که رسیدن به صورت جدید وارسی مدل است، به ما نشان می‌دهد. طبق \cite{calcul} یک دلیل برای استفاده از عبارات منظم، که متداول‌تر بودن آن بین جامعه‌ی برنامه نویسان است. منظور این است که برنامه نویسان از شاخه‌های مختلف این تخصص همگی با عبارات منظم آشنا هستند، درحالیکه برنامه نویسانی که با منطق‌های زمانی آشنایی داشته باشند، بسیار کم هستند و جایگزینی اولی با دومی می‌تواند به رواج استفاده از روش وارسی مدل برای درستی‌یابی برنامه‌های کامپیوتری کمک کند. ابتدا، نحو این منطق را تعریف می‌کنیم، سپس به سراغ معناشناسی آن می‌رویم. عبارات منظمی که در اینجا تعریف می‌کنیم، با عبارات منظم سنتی که در علوم کامپیوتر رایج هستند، در جزئیات متفاوت هستند. 
\subsection{نحو عبارات منظم}
فرق عمده‌ای که نحو عبارات منظم ما با عبارات منظم سنتی دارد، در \gls{atom}هاست. اتم‌ها در عبارات منظم سنتی بدون ساختار بودند، اما در اینجا، ساختار دارند. در اینجا، هر اتم یک زوج مرتب است، متشکل از مجموعه‌ی $\mathsf{L}$ شامل برچسب‌ها و عبارت بولی $\mathsf{B}$. این زوج مرتب را به شکل 
$\mathsf{L : B}$
نمایش می‌دهیم.\\ 
با یادآوری اینکه 
$\mathbb{X,B,R}$
به ترتیب مجموعه‌های عبارت‌های منظم، عبارت‌های بولی و متغیرها هستند، نحو عبارات منظم به شکل زیر است:
\begin{defn}
	(نحو عبارات منظم): اگر فرض کنیم
$$\mathsf{L} \in \mathit{P}(\mathbb{L}),$$        
$$\mathsf{x,y,...} \in \mathbb{X},$$
$$\mathsf{\underline{x},\underline{y},...} \in \mathbb{\underline{X}},$$
$$\mathsf{B} \in \mathbb{B},$$
$$\mathsf{R} \in \mathbb{R},$$
$\mathbb{R}$ توسط گرامر زیر تولید می‌شود.
$$\mathsf{R} ::=\hspace{0.5cm} \varepsilon\hspace{4.2cm}$$
$$|\:\:\:\mathsf{L : B}\hspace{2.4cm}$$
$$|\:\:\:\mathsf{R_1 R_2}\:\:\:(or\:\mathsf{R_1 \bullet R_2 })$$
$$\:\:\:\:\:\:\:|\:\:\:\mathsf{R_1\:\mid\:R_2}\:\:\:(or\:\mathsf{R_1 + R_2 })$$
$$|\:\:\:\mathsf{R_1^*}\hspace{2.8cm}$$
$$|\:\:\:\mathsf{R_1^+}\hspace{2.7cm}$$
$$|\:\:\:(\mathsf{R_1})\hspace{2.47cm}$$
\end{defn}
همان طور که قابل مشاهده است، در اینجا، \glspl*{binop} چسباندن
$(\bullet)$
 و \gls{choice}
$(|)$
 را به همراه \glspl*{unop}
$^*$
و
$^+$
داریم.
در ادامه، با توجه به تعریف \gls{regexsem}، خواهیم دید که معنای عملگر یگانی 
$^+$
به وسیله‌ی عملگر یگانی $^*$ قابل بیان است.
توجه شود که پرانتزها هم جزئی از نحو قرار داده شده‌اند.


نمادگذاری: منظور از عبارت منظم 
$\mathsf{? : B}$
همان عبارت منظم
$\mathbb{L}\mathsf{ : B}$
است. عبارت منظم 
$l : \mathsf{B}$
به جای عبارت منظم 
$\{l\} : \mathsf{B}$
به کار می‌رود و منظور از عبارت منظم
$\neg l \mathsf{: B}$
نیز عبارت منظم 
$\mathbb{L}\setminus\{l\}:\mathsf{B}$
است. 

همچنین یک نکته‌ی قابل توجه، با توجه به تعاریف فصل قبل، وجود یک مجموعه به نام
$\mathbb{\underline{X}}$
در کنار 
$\mathbb{X}$
که از قبل داشتیم، است. به ازای هر 
$\mathsf{x}\in\mathbb{X}$
یک 
$\mathsf{\underline{x}}\in\mathbb{\underline{X}}$
داریم، به عبارت دیگر 
$\mathbb{\underline{X}}=\{\underline{\mathsf{x}}|\mathsf{x}\in \mathbb{X}\}$.
$\mathsf{\underline{x}}$
متغیر 
$\mathsf{x}$
در ابتدای هر برنامه است.  همان‌طور که تعریف می‌کنیم:
\begin{defn}
	(\gls{initenv}): به هر تابع  
	$\underline{\rho}:\underline{\mathbb{X}}\rightarrow\mathbb{V}$
	یک محیط اولیه می‌گوییم. مجموعه‌ی همه‌ی محیط‌های اولیه را با 
	$\underline{\mathbb{EV}}$
	نمایش می‌دهیم.
\end{defn}
 
\subsection{معناشناسی عبارات منظم}
معنای هر عبارت منظم را با استفاده از تابع 
$\mathcal{S}^r$
به نام \gls{regexsemantics} تعریف می‌کنیم. این تابع در ورودی‌اش یک عبارت منظم 
$\mathsf{R}$
را می‌گیرد و یک مجموعه از زوج مرتب‌های
$\langle\underline{\rho} , \pi\rangle$
را که 
$\pi \in \mathfrak{S^*}$
و 
$\underline{\rho} \in \underline{\mathbb{EV}}$
باز می‌گرداند. بنابراین این تابع از \gls*{type}
$\mathbb{R} \rightarrow \mathit{P} (\mathbb{\underline{EV}} \times \mathfrak{S}^*)$
است.
 منظور از 
 $\mathfrak{S}^*$
 نیز 
 $\mathfrak{S}^+ \cup \{\epsilon\}$
 است.
 
تعریف \gls*{inductive} تابع 
$\mathcal{S}^r$
را در ادامه آورده‌ایم.
\begin{defn}
	(معناشناسی عبارات منظم):
	تابع 
	$\mathcal{S}^r:\mathbb{R}\rightarrow \mathit{P}(\mathbb{\underline{EV}\times \mathfrak{S}}^*)$
	به صورت استقرایی روی نحو عبارات منظم به صورت زیر تعریف می‌شود:
$$\mathcal{S}^r\llbracket\varepsilon\rrbracket= \{ \langle \underline{\rho} , \epsilon \rangle | \underline{\rho} \in \underline{\mathbb{EV}}\}.$$
معنای عبارت منظم
$\varepsilon$
مجموعه‌ای شامل زوج مرتب‌هایی از محیط‌های اولیه‌ی مختلف در کنار رد پیشوندی تهی است.
$$\mathcal{S}^r\llbracket\mathsf{L:B}\rrbracket = \{\langle\underline{\rho},\langle l , \rho \rangle \rangle | l \in \mathsf{L} \land \mathcal{B}\llbracket\mathsf{B}\rrbracket \underline{\rho},\rho \}.$$
معنای عبارت منظم
$\mathsf{L:B}$
مجموعه‌ای است شامل زوج مرتب‌هایی که عضو اول آن‌ها محیط‌های اولیه‌ی مختلف و عضو دوم آن‌ها ردهای پیشوندی تک‌عضوی 
$\langle l , \rho \rangle$
هستند. در این ردهای پیشوندی، برچسب 
$l$
باید در
$\mathsf{L}$
که مجموعه‌ای از برچسب‌هاست حضور داشته باشد. همین طور باید عبارت بولی 
$\mathsf{B}$
درباره‌ی محیط اولیه
$\underline{\rho}$
و محیط 
$\rho$
برقرار باشد. با توجه به حضور محیط‌های اولیه، در اینجا
$\mathcal{B}$
به جای اینکه از نوع 
$\mathbb{EV} \rightarrow \mathbb{BOOL}$
باشد، از نوع
$\underline{\mathbb{EV}} \rightarrow \mathbb{EV} \rightarrow \mathbb{BOOL}$
است( منظور از 
$\mathbb{BOOL}$
همان مجموعه‌ی 
$\{\mathit{T,F}\}$
است). بعد از تمام شدن تعریف معناشناسی عبارات منظم، دوباره
$\mathcal{A}$
و 
$\mathcal{B}$
را با در نظر گرفتن محیط اولیه تعریف خواهیم کرد.
$$\mathcal{S}^r\llbracket\mathsf{R_1 R_2}\rrbracket= \mathcal{S}^r\llbracket\mathsf{R_1}\rrbracket \Join \mathcal{S}^r\llbracket\mathsf{R_2}\rrbracket.$$
\begin{center}
	به ‌طوریکه، با فرض اینکه دو مجموعه‌ی 
	$\mathcal{S}$
	و
	$\mathcal{S'}$
	هر یک معنای یک عبارت منظم باشند:
$$\mathcal{S \Join S'}=
\{
\langle \underline{\rho}, \pi \pi' \rangle |
\langle \underline{\rho},\pi \rangle \in \mathcal{S} \land
\langle \underline{\rho},\pi' \rangle \in \mathcal{S'}
\}.
$$
\end{center}
اگر یک عبارت منظم داشته باشیم که از چسباندن دو عبارت منظم
$\mathsf{R_1}$
و 
$\mathsf{R_2}$
به هم ساخته شده باشد، آنگاه معنای این عبارت منظم از زوج مرتب‌هایی تشکیل شده است که عضو اول آن‌ها محیط‌های اولیه هستند و عضو دوم آن‌ها از چسباندن ردهای پیشوندی موجود در عضو دوم اعضای مجموعه‌ی معنای این دو عبارت منظم تشکیل شده است. این عملگر چسباندن که در اینجا برای معنای عبارت‌های منظم تعریف شده است، با عملگر چسباندن برای معنای برنامه‌ها که در فصل پیش ارائه شد، متفاوت است. مورد دوم هم با نماد $\Join$ روی ردهای پیشوندی تعریف ‌شد، اما در تعریفی که در این فصل از چسباندن ارائه داده‌ایم، این عملگر روی ردهای پیشوندی تعریف نشده است، بلکه روی زوج‌های مرتبی که در معنای هر عبارت منظمی حضور دارند، تعریف می‌شود.

تا این قسمت از تعریف معناشناسی عبارات منظم که رسیده‌ایم، تا حدی به درکی شهودی از اینکه به چه شکلی عبارات منظم راهی برای توصیف ویژگی‌ در مورد برنامه‌ها هستند، دست پیدا کرده‌ایم. همان‌طور که در حالت قبل دیدیم، \gls*{consistency} هر زوج مرتب
$\mathsf{L:B}$
با یک وضعیت داخل یک رد پیشوندی بررسی می‌‌شود. این زوج‌ مرتب‌ها موازی با وضعیت‌ها در ردهای پیشوندی معنای برنامه پیش می‌روند و سازگاری آن‌ها با یکدیگر بررسی می‌شود تا وارسی مدل انجام شود.
$$\mathcal{S}^r \llbracket\mathsf{R_1\:|\:R_2}\rrbracket= 
\mathcal{S}^r \llbracket R_1\rrbracket \cup
\mathcal{S}^r \llbracket R_2\rrbracket.$$
این حالت، معنای \gls*{apply} عملگر انتخاب روی دو عبارت منظم را توصیف می‌کند. معنای اعمال این عملگر به صورت اجتماع معنای هر دو عبارت منظم تعریف شده است.

$$\mathcal{S}^r \llbracket\mathsf{R}\rrbracket^0 = \mathcal{S}^r\llbracket\varepsilon\rrbracket,$$
$$\mathcal{S}^r \llbracket\mathsf{R}\rrbracket^{n+1} = \mathcal{S}^r \llbracket\mathsf{R}\rrbracket^{n} \Join
\mathcal{S}^r \llbracket\mathsf{R}\rrbracket.$$
دو عبارت اخیر، برای تعریف معنای عملگرهای $^*$ و $^+$ تعریف شده‌اند. عملگر $\Join$ و معنای 
$\mathcal{S}^r\llbracket\varepsilon\rrbracket$
را هم که قبلا تعریف کرده بودیم و $0$ و $n$ هم اعداد طبیعی‌اند.
$$\mathcal{S}^r\llbracket\mathsf{R^*}\rrbracket =  \bigcup_{n \in \mathbb{N}}
\mathcal{S}^r \llbracket\mathsf{R}\rrbracket^n,$$
$$\mathcal{S}^r\llbracket\mathsf{R^+}\rrbracket =  \bigcup_{n \in \mathbb{N}\setminus\{0\}}
\mathcal{S}^r \llbracket\mathsf{R}\rrbracket^n.$$
این دو عبارت تعریف معنای دو عملگر $^*$ و $^+$ هستند. منظور از $\mathbb{N}$ مجموعه‌ی اعداد طبیعی است. همان‌طور که قبل‌تر هم اشاره شد، معنای عملگر $^+$ با عملگر $^*$ قابل بیان است. اضافه می‌کنیم که معنای عملگر $^*$ را در \gls*{meta}( و نه در منطق تعریف شده در اینجا) می‌توان با عملگر انتخاب تعریف کرد.

$$\mathcal{S}^r \llbracket(\mathsf{R})\rrbracket=\mathcal{S}^r \llbracket\mathsf{R}\rrbracket.$$
این قسمت از تعریف هم صرفا بیان می‌کند که پرانتزها تاثیری در معنای عبارات منظم ندارند که کاملا قابل انتظار است، چرا که وجود پرانتز قرار است صرفا در خواص نحو منطق نقش داشته باشد.

\end{defn}

\begin{defn}
	(\gls*{satisfy}): می‌گوییم، در محیط اولیه‌ی $\underline{\rho}$ رد پیشوندی $\pi$ عبارت منظم $\mathsf{R}$ را ارضا می‌کند، اگر و تنها اگر 
	$\langle \underline{\rho},\pi \rangle \in \mathcal{S}^r \llbracket \mathsf{R} \rrbracket$. 
\end{defn}

تعریف معناشناسی عبارات منظم در اینجا تمام می‌شود، اما همان‌گونه که در لا‌به‌لای تعاریف گفته ‌شد، احتیاج داریم که $\mathcal{A}$ و $\mathcal{B}$ را از نو تعریف کنیم:
\begin{defn}
	توابع 
	$\mathcal{A}:\mathbb{A}\rightarrow \mathbb{\underline{EV}} \rightarrow
	(\mathbb{EV} \rightarrow \mathbb{V})$
	و 
	$\mathcal{B}:\mathbb{B}\rightarrow (\mathbb{\underline{EV}} \rightarrow
	(\mathbb{EV} \rightarrow \mathbb{BOOL}))$
	به شکل استقرایی به ترتیب روی ساختار عبارت‌های حسابی 
	$\mathsf{A}\in\mathbb{A}$
	و عبارت‌های بولی
	$\mathsf{B} \in \mathbb{B}$
	به شکل زیر تعریف می‌شوند:
	
	
$$\mathcal{A}\llbracket\mathsf{1}\rrbracket\underline{\rho},\rho=1,$$
$$\mathcal{A}\llbracket\mathsf{\underline{\mathsf{x}}}\rrbracket\underline{\rho},\rho= \underline{\rho}(\mathsf{x}),$$
$$\mathcal{A}\llbracket\mathsf{\mathsf{x}}\rrbracket\underline{\rho},\rho= \rho(\mathsf{x}),$$
$$\mathcal{A}\llbracket\mathsf{A_1 - A_2}\rrbracket\underline{\rho},\rho= 
\mathcal{A}\llbracket\mathsf{A_1}\rrbracket\underline{\rho},\rho - \mathcal{A}\llbracket\mathsf{A_2}\rrbracket\underline{\rho},\rho,$$
$$\mathcal{B}\llbracket\mathsf{A_1 < A_2}\rrbracket\underline{\rho},\rho=
	\mathcal{A}\llbracket\mathsf{A_1}\rrbracket\underline{\rho},\rho\: <\: \mathcal{A}\llbracket\mathsf{A_2}\rrbracket\underline{\rho},\rho,$$
$$\mathcal{B}\llbracket\mathsf{B_1\:nand\:B_2}\rrbracket\underline{\rho},\rho=
\mathcal{B}\llbracket\mathsf{B_1}\rrbracket\underline{\rho},\rho \uparrow 
\mathcal{B}\llbracket\mathsf{B_2}\rrbracket\underline{\rho},\rho .$$
\end{defn}
قابل مشاهده است که تعاریف جدید تا حد خوبی به تعاریف قبلی شبیه هستند و فرق عمده صرفا وارد شدن $\underline{\rho}$ است.\\

حال که معناشناسی عبارات منظم را داریم، به طور مختصر به مقایسه‌ی عبارات منظمی که در این بحث تعریف کرده‌ایم و عبارات منظم کلاسیک در باقی نوشته‌ها و موضوعات می‌پردازیم. \gls{kleenealgebra} یک  \gls*{algstruc} است که تعمیمی است از عبارات منظم معرفی شده در \cite{kleene56}. سر و کله‌ی عبارات منظم با تعاریف غیر هم‌ارز در قسمت‌های مختلفی از علوم کامپیوتر پیدا می‌شود. هدف از ارائه‌ی جبر کلین ارائه‌ی یک تعریف عام به منظور در بر داشتن تعاریف \gls*{unequiv} بوده است. البته، برای جبر کلین هم تعاریف مختلف غیر هم‌ارز معرفی شده اند که در \cite{DexterKleene} این تعاریف و ارتباطشان با یکدیگر بررسی شده‌ است. همین طور، این مقاله خود با یک تعریف از جبر کلین شروع کرده است که به شکل زیر است. 
\begin{defn}
	(جبر کلین): به مجموعه‌ی $\mathbb{K}$ به همراه دو عملگرهای دوتایی $\bullet$ و $+$ و عملگر یگانی $^*$ و ثابت‌های $0$ و $1$ یک جبر کلین می‌گویند، اگر و تنها اگر 
	$(\mathbb{K}, +, \bullet, 0 , 1)$
	یک \gls*{semiring}ی \gls*{idempotent} باشد، به طوریکه عملگر $+$ \gls*{associative}، \gls*{commutative} و خودتوان با \gls*{identity}‌ی $0$ باشد; عملگر $\bullet$ شرکت پذیر با عنصر یکه‌ی $1$ باشد; $\bullet$ روی $+$ از هر دو طرف پخش پذیر باشد و $0$ برای $\bullet$ قانون جذب داشته باشد.
\end{defn}

طبق این تعریف، اگر عبارات منظمی که در اینجا تعریف کرده‌ایم، یک جبر کلین می‌بود، می‌بایستی که برای هر عبارت منظم $\mathsf{R}$ می‌داشتیم
$\mathcal{S}^r \llbracket \mathsf{\varepsilon R} \rrbracket = \mathcal{S}^r \llbracket \mathsf{\varepsilon} \rrbracket$،
چون $\varepsilon$ نقش صفر عملگر چسباندن( که \gls*{multop} جبر کلین است) را دارد و یک جبر کلین برای صفر \gls*{absroplaw} را به عنوان یک \gls*{law} دارد. اما در مورد عبارات منظمی که در اینجا تعریف کرده‌ایم، داریم 
$\mathcal{S}^r \llbracket \mathsf{\varepsilon R} \rrbracket = \mathcal{S}^r \llbracket \mathsf{R} \rrbracket$.
پس می‌توانیم همینجا نتیجه بگیریم که عبارات منظمی که داریم، با یک تعریف استاندارد و فراگیر از عبارات منظم ناسازگار است. پس باید در نظر داشته باشیم که ممکن است بسیاری از خواص مورد مطالعه قرار گرفته در مورد عبارات منظم کلاسیک، در مورد عبارات منظمی که در اینجا تعریف کرده‌ایم برقرار نباشد.
بیشتر از این به این بحث نمی‌پردازیم که بحث دامنه‌دار و منحرف کننده‌ایست. یک قضیه را در مورد عبارات منظم ارائه می‌دهیم که \gls*{distrib} عملگر انتخاب به عملگر چسباندن است و بعد به ادامه‌ی راه اصلیمان می‌پردازیم.

\begin{thm}
	برای عبارت‌های منظم
	 $\mathsf{R , R}_1,\mathsf{R}_2$
	 داریم:
	  $$\mathcal{S}^r \llbracket \mathsf{R}\bullet (\mathsf{R}_1 + \mathsf{R}_2) \rrbracket = \mathcal{S}^r \llbracket (\mathsf{R \bullet R}_1) + (\mathsf{R \bullet R}_2) \rrbracket. $$
\end{thm}

\begin{proof}
	$$\mathcal{S}^r \llbracket \mathsf{R}\bullet (\mathsf{R}_1 + \mathsf{R}_2) \rrbracket =\mathcal{S}^r \llbracket \mathsf{R} \rrbracket \Join \mathcal{S}^r \llbracket (\mathsf{R}_1+\mathsf{R}_2) \rrbracket$$
	$$=\mathcal{S}^r \llbracket \mathsf{R} \rrbracket \Join \mathcal{S}^r \llbracket \mathsf{R}_1+\mathsf{R}_2 \rrbracket=\mathcal{S}^r \llbracket \mathsf{R} \rrbracket \Join (\mathcal{S}^r \llbracket \mathsf{R}_1 \rrbracket \cup \mathcal{S}^r\llbracket \mathsf{R}_2 \rrbracket)$$
	$$= \{ \langle \underline{\rho}, \pi \pi' \rangle| \langle \underline{\rho}, \pi \rangle \in \mathcal{S}^r \llbracket \mathsf{R} \rrbracket \land (\langle \underline{\rho},\pi'\rangle \in \mathcal{S}^r \llbracket \mathsf{R}_1 \rrbracket\lor \langle \underline{\rho} , \pi' \rangle \in \mathcal{S}^r \llbracket \mathsf{R}_2 \rrbracket) \}$$
	$$=\{ \langle \underline{\rho}, \pi \pi' \rangle| (\langle \underline{\rho}, \pi \rangle \in \mathcal{S}^r \llbracket \mathsf{R} \rrbracket \land \langle \underline{\rho},\pi'\rangle \in \mathcal{S}^r \llbracket \mathsf{R}_1 \rrbracket) \lor (\langle \underline{\rho}, \pi \rangle \in \mathcal{S}^r \llbracket \mathsf{R} \rrbracket \land \langle \underline{\rho} , \pi' \rangle \in \mathcal{S}^r \llbracket \mathsf{R}_2 \rrbracket  \}$$
	$$=\{ \langle \underline{\rho}, \pi \pi' \rangle|\langle \underline{\rho}, \pi \rangle \in \mathcal{S}^r \llbracket \mathsf{R} \rrbracket \land \langle \underline{\rho},\pi'\rangle \in \mathcal{S}^r \llbracket \mathsf{R}_1 \rrbracket\} \cup
	\{ \langle \underline{\rho}, \pi \pi' \rangle|\langle \underline{\rho}, \pi \rangle \in \mathcal{S}^r \llbracket \mathsf{R} \rrbracket \land \langle \underline{\rho},\pi'\rangle \in \mathcal{S}^r \llbracket \mathsf{R}_2 \rrbracket\}$$
	$$=(\mathcal{S}^r \llbracket \mathsf{R} \rrbracket \Join \mathcal{S}^r \llbracket \mathsf{R}_1\rrbracket ) \cup (\mathcal{S}^r \llbracket \mathsf{R} \rrbracket \Join \mathcal{S}^r \llbracket \mathsf{R}_2 \rrbracket)=
	\mathcal{S}^r \llbracket \mathsf{R\bullet R}_1 \rrbracket \cup \mathcal{S}^r \llbracket \mathsf{R\bullet R}_2 \rrbracket$$
	$$=\mathcal{S}^r \llbracket (\mathsf{R \bullet R}_1) + (\mathsf{R \bullet R}_2) \rrbracket $$
\end{proof}
تا اینجای کار، بیشتر مفاهیمی که برای بیان صورت جدید روش وارسی مدل احتیاج داریم را بیان کرده‌ایم، هرچند که هنوز برخی از تعاریف باقی مانده‌اند.
\subsection{گونه‌های مختلف نحو عبارات منظم}
به عنوان قسمت آخر این بخش، گونه‌های مختلفی از نحو عبارات منظم را بیان می‌کنیم که هر کدام در واقع زیرمجموعه‌ای از منطقی که توصیف کرده‌ایم را تشکیل می‌دهند. بعضی از آن‌ها را در همین فصل و برخی دیگر را در فصل بعدی استفاده می‌کنیم.

می‌توانیم به هر یک از این عبارات منظم به چشم یک منطق متفاوت نگاه کنیم. هرچند که معناشناسی عبارات منظم به عنوان معناشناسی هر یک از این منطق‌ها باقی می‌ماند. با توجه به اینکه هر یک از این منطق‌ها زیرمجموعه‌ی عبارات منظم هستند، می‌توان گفت معناشناسی هر یک از آن‌ها همان تابع $\mathcal{S}^r$ با \gls*{domain}ی محدود شده است.

اولین گونه‌ای که می‌خواهیم بیان کنیم، گونه‌ای است که در اعضای آن اصلا اتم حضور ندارد و کل عبارت‌های منظم عضو این منطق از $\varepsilon$ها تشکیل شده‌اند.
\begin{defn}
	(\gls{emptyregex} - $\mathbb{R_\varepsilon}$): 
	$\mathbb{R}_\varepsilon$ 
	توسط گرامر زیر تولید می‌شود.
	$$\mathsf{R}\:\:\: ::= \:\:\: \varepsilon \: | \: \mathsf{R_1 R_2} \: | 
	\: \mathsf{R_1 + R_2} \: | \: \mathsf{R_1^*} \: | \: \mathsf{R_1^+} | (\mathsf{R_1}).$$
\end{defn}


 با توجه به آنچه گفتیم، معنای همه‌ی این عبارت‌های منظم عضو این منطق برابر 
$\{\langle \underline{\rho} , \epsilon \rangle\}$
خواهد بود.

گونه‌ی بعدی \gls{nonemptyregex} است.
\begin{defn}
	(عبارات منظم ناتهی - $\mathbb{R}^+$):
$\mathbb{R}^+$ توسط گرامر زیر تولید می‌شود.
	$$\mathsf{R}\:\:\: ::= \:\:\: \mathsf{L:B} \: |
	\: \mathsf{\varepsilon R_2} \: | \: \mathsf{R_1 \varepsilon} \: |
	\: \mathsf{R_1 R_2} \: | 
	\: \mathsf{R_1 + R_2} \: | \: \mathsf{R_1^+} | (\mathsf{R_1}).$$
\end{defn}
دلیل وجود 
$\mathsf{\varepsilon R_2}$
و 
$\mathsf{R_1 \varepsilon}$
در تعریف این است که ممکن است معنای عبارت منظمی در این منطق با ${\langle \underline{\rho} , \epsilon \rangle}$ برابر نباشد، اما در آن عبات منظم، $\varepsilon$ حضور داشته باشد. با این تفاصیل، می‌توان دید که دو مجموعه‌ی 
$\mathbb{R}_\varepsilon$
و
$\mathbb{R}^+$
یک \gls*{partition} برای مجموعه‌ی $\mathbb{R}$ هستند، چونکه معنای هر عبارت منظم در $\mathbb{R}$ یا با ${\langle \underline{\rho} , \epsilon \rangle}$ برابر هست یا نیست. بنابراین شاید به نظر برسد که تعریف یکی از این نحوها کافی باشد. اما این‌طور نیست، چون ممکن است که درجایی احتیاج داشته باشیم که تعریفی استقرایی روی هر یک از این دو ساختار ارائه دهیم، یا اینکه در اثبات حکمی بخواهیم از استقرا روی یکی از این دو ساختار استفاده کنیم.

گونه‌ی آخر عبارات منظم ما نیز \gls{choicefreeregex} است.
\begin{defn}
(عبارات منظم بدون انتخاب - $\mathbb{R}^\nmid$):
$\mathbb{R}^\nmid$ توسط گرامر زیر تولید می‌شود.
$$
\mathsf{R}\:\:\: ::= \:\:\: \varepsilon \: | \: \mathsf{L:B} \: | \:
\mathsf{R_1 R_2} \: | \: \mathsf{R_1}^* \: | \: \mathsf{R_1}^+ \: | \:
(\mathsf{R_1}).$$
\end{defn}

\section{صورت جدید مسئله‌ی وارسی مدل}
بالاخره، به هدف نهایی این فصل رسیدیم. می‌خواهیم صورت جدیدی از مسئله‌ی وارسی مدل را بیان ‌کنیم.\\
پیش از ارائه‌ی تعریف وارسی مدل، نیاز داریم که عملگر \gls{prefixclosure} را برای یک مجموعه از ردهای پیشوندی معرفی کنیم.
\begin{defn}
	(بستار پیشوندی):
	اگر 
	$\Pi \in \mathit{P}(\mathbb{\underline{EV}\times \mathfrak{S}^+})$،
	آنگاه بستار پیشوندی $\Pi$ را به صورت زیر تعریف می‌کنیم:
	$$\mathsf{prefix}(\Pi)=
	\{\langle \underline{\rho},\pi\rangle | \pi \in \mathfrak{S^+} \land \exists\
	\pi' \in \mathfrak{S^*}: \langle \underline{\rho}, \pi \pi' \rangle \in \Pi \}.$$
\end{defn}

برای درک بهتر مفهوم بستار پیشوندی به مثال زیر توجه شود.
\begin{exm}
	اگر 
	$\Pi = \{\langle\underline{\rho},\langle l_1 , \rho_1 \rangle 
	\langle l_2  \rho_2 \rangle\rangle  \langle l_3 , \rho_3 \rangle\rangle 
	,
	\langle\underline{\rho},\langle {l_1}' , {\rho_1}' \rangle  \langle {l_2}'  {\rho_2}' \rangle\rangle
	\}$
	باشد( $\Pi$ شامل دو عضو است)، آنگاه:
	$$
	\mathsf{prefix}(\Pi)=
	\{
	\langle\underline{\rho},\langle l_1 , \rho_1 \rangle\rangle ,
	\langle\underline{\rho},\langle l_1 , \rho_1 \rangle  \langle l_2  \rho_2 \rangle\rangle,
	\langle\underline{\rho},\langle l_1 , \rho_1 \rangle  \langle l_2  \rho_2 \rangle  \langle l_3 , \rho_3 \rangle\rangle ,$$
	$$
	\langle\underline{\rho},\langle {l_1}' , {\rho_1}' \rangle\rangle,
	\langle\underline{\rho},\langle {l_1}' , {\rho_1}' \rangle  \langle {l_2}'  {\rho_2}' \rangle\rangle
	\}.
	$$

\end{exm}
حال به ارائه‌ی صورت جدید روش وارسی مدل می‌رسیم.
\begin{defn}
(وارسی مدل):
	اگر 
	$\mathsf{P}\in\mathbb{P} , \mathsf{R} \in \mathbb{R}^+ , \underline{\rho} \in \underline{\mathbb{EV}}$
	آنگاه:
	$$\mathsf{P},\underline{\rho} \models \mathsf{R}
	 \Leftrightarrow
	(\{\underline{\rho}\}\times \mathcal{S}^* \llbracket\mathsf{P}\rrbracket) \subseteq 
	\mathsf{prefix} (\mathcal{S}^r \llbracket\mathsf{R} \bullet (?:\mathit{T})^*\rrbracket).
	$$
\end{defn}

این تعریف بیان می‌کند که اگر برنامه‌ی 
$\mathsf{P}$
با محیط اولیه‌ی 
$\underline{\rho}$
اجرا شود، این برنامه در صورتی ویژگی‌ای که با یک عبارت منظم 
$\mathsf{R}$
بیان شده است را دارد که معنای آن زیرمجموعه‌ی بستار پیشوندی معنای عبارت منظم
$\mathsf{R} \bullet (?:\mathit{T})^*$
باشد.
توجه شود که محیط اولیه‌ای که برای برنامه‌ی تحت بررسی در نظر گرفته‌ایم، صرفا به این منظور در تعریف قرار داده شده است که بتوانیم،  معنای برنامه‌ها را با معنای عبارت‌های منظم قابل قیاس کنیم. دلیل حضور محیط اولیه در معنای عبارت‌های منظم نیز در صورت سوم روش وارسی مدل مشخص می‌شود، یعنی جایی که وارسی مدل روی ساختار برنامه‌ها تعریف شده است و ردهای پیشوندی موجود در معنای هر عبارت دستوری یا لیست عبارت‌های دستوری با محیط متفاوتی شروع می‌شوند و اطلاعات محیط اولیه‌ی برنامه در این ردها حضور ندارد (با اینکه ممکن است به آن نیاز داشته باشیم). در این صورت از روش وارسی مدل و صورت بعدی، محیط‌های اولیه صرفا حضور دارند و در عمل نقش مهمی ندارند. 

در مورد نقش عبارت منظم
$ (?:\mathit{T})^*$
و عملگر
$ \mathsf{prefix} $
نیز می‌توان گفت، بنا به تصمیم مبدع این روش، اگر یک رد پیشوندی، با عبارت منظم ناسازگاری نداشته باشد و از آن کوتاه‌تر باشد، این رد پیشوندی با ویژگی بیان شده را دارد. این نکته صرفا در مورد حضور 
$\mathsf{prefix}$ 
است. حضور 
$(?:\mathit{T})^*$
نیز باعث می‌شود اگر طول عبارت منظم مورد بررسی کمتر از رد پیشوندی باشد و ناسازگاری‌ای مشاهده نشده باشد، رد پیشوندی دارای ویژگی $\mathsf{R}$ در نظر گرفته شود.
در اینجا، بیان صورت جدید وارسی مدل به اتمام می‌رسد.
\section{در مورد توقف پذیری}

در این بخش نکته‌ای در مورد کار که به نظر نگارنده رسیده، مطرح شده است. اگر صحبت ما در اینجا درست باشد، این به این معنی خواهد بود که روشی که در حال توصیفش هستیم،\gls*{undecidable} است، به عبارت دیگر قابلیت پیاده‌سازی را ندارد!

بحث ما در اینجا در مورد \gls*{halting} است. در \cite{calcul} در مورد تعریف \gls*{halt} یک برنامه صحبتی به میان نیامده است. یک تعریف  صوری که خودمان می‌توانیم برای این مفهوم، بیاوریم این است:

\begin{defn}
(توقف پذیری:) برنامه‌ی $\mathsf{P}$ را به همراه محیط اولیه
 $\underline{\rho}$
  \gls{terminable}
  می‌گوییم، اگر و تنها اگر وجود داشته باشد 
	$\pi \in \mathcal{S}^* \llbracket \mathsf{P} \rrbracket$، 
	که ($\rho$ محیط متناظر با محیط اولیه‌ی $\underline{\rho}$ است):
	$$\pi = \langle at \llbracket \mathsf{P} \rrbracket , \rho \rangle \pi'$$
	و اینکه $\langle aft\llbracket \mathsf{P} \rrbracket , \rho' \rangle$ در $\pi$ حضور داشته باشد. در این صورت می‌نویسیم
	$\mathsf{P,\underline{\rho}\downarrow}$.
همین تعریف را برای لیست عبارت‌های دستوری $\mathsf{Sl}$ یا عبارت‌ دستوری $\mathsf{S}$ صرفا با جایگذاری به‌جای برنامه‌ی $\mathsf{P}$ نیز داریم.
\end{defn} 
در این تعریف، توقف پذیری صرفا برای یک محیط اولیه تعریف شده است.
در اینجا، توقف پذیری به متناهی بودن ردهای پیشوندی موجود در برنامه ارتباطی ندارد. با توجه به معناشناسی‌ رد پیشوندی، تعریف کردن توقف پذیری به معنای وجود رد پیشوندی متناهی در معنای برنامه به هیچ عنوان مناسب نیست، چون معناشناسی خاصیت پیشوندی بودن را دارد و مطمئن هستیم، در معنای هر برنامه‌ای حتما یک رد پیشوندی متناهی با محیط اولیه‌ی مورد بررسی وجود دارد.

اگر هم بخواهیم، تعریف توقف پذیری را عدم حضور ردهای پیشوندی نامتناهی در معنای برنامه در نظر بگیریم، به تعریف قوی‌تری نسبت به آنچه ارائه دادیم، رسیده‌ایم. در اینجا، ما سعی می‌کنیم، ضعیف‌ترین تعریف معقول را ارائه دهیم. در این صورت، از اثباتی که بر اساس تعریف ارائه شده، برای تصمیم ناپذیری می‌آوریم، مطمئن‌تر خواهیم شد. با این حال، در قضیه‌ی بعدی می‌بینیم که تعریفی که ارائه کرده‌ایم با اینکه بگوییم، در معنای برنامه رد پیشوندی نامتناهی وجود ندارد معادل است.
\begin{thm}
	برای برنامه‌ی $\mathsf{P}$ و محیط اولیه‌ی $\underline{\rho}$ داریم $\mathsf{P} , \underline{\rho} \downarrow $، اگر و تنها اگر $\rho$ محیط متناظر با محیط اولیه‌ی $\underline{\rho}$ باشد و
	$$\forall \pi \in \mathfrak{S^{+}} : 
	\langle at \llbracket \mathsf{P} \rrbracket , \rho \rangle \pi \in \mathcal{S^*} \llbracket \mathsf{P} \rrbracket \rightarrow
	\langle at \llbracket \mathsf{P} \rrbracket , \rho \rangle \pi \in \mathbb{R^+}.$$
	 
\end{thm}
\begin{proof}
	$(\Rightarrow)$
	برای این حالت، باید ثابت کنیم که در معنای هر برنامه‌ای رد پیشوندی‌ای وجود دارد که با
	$\langle at\llbracket \mathsf{P} \rrbracket , \rho \rangle$
	شروع شده است و به ازای یک محیط $\rho'$ به 
	$\langle aft\llbracket \mathsf{P} \rrbracket , \rho' \rangle$
	ختم شده است. 
	در این اثبات، از تعریف برچسب‌ها که در ضمیمه‌ی \cite{calcul} آمده است، استفاده کرده‌ایم.
	داریم:
	$\mathsf{P=Sl}$
	و 
	$aft \llbracket \mathsf{P} \rrbracket = aft \llbracket \mathsf{Sl} \rrbracket $.
	
	حکم را با \gls*{induction} روی ساختار $\mathsf{Sl}$ ثابت می‌کنیم. 
	$$\blacktriangleright \mathsf{Sl = \backepsilon}:$$
	داریم:
	$$\mathcal{S^*} \llbracket \backepsilon \rrbracket = \{\langle at \llbracket \backepsilon \rrbracket , \dot{\rho} \rangle | \dot{\rho} \in \mathbb{EV}\}.$$
	همین طور طبق تعریف برچسب‌ها داریم:
	$$at \llbracket \backepsilon \rrbracket = aft \llbracket \backepsilon \rrbracket$$
	پس حکم برقرار است.
	$$\blacktriangleright \mathsf{Sl =Sl'\;S}:$$
	اینکه در معنای $\mathsf{Sl}$، رد پیشوندی‌ای شامل
	 $\langle aft \llbracket \mathsf{Sl} \rrbracket , \rho '\rangle$ وجود داشته باشد، با توجه به تعاریفی که داریم، به این وابسته است که در معنای
	  $\mathsf{S}$،
	   دنباله‌ای شامل $\langle aft \llbracket \mathsf{S} \rrbracket , \rho' \rangle$ وجود داشته باشد. برای اینکه این گزاره را ثابت کنیم، باید آن را روی ساختار $\mathsf{S}$ ثابت کنیم که در واقع بخش اصلی اثبات این سمت قضیه است. 
	$$\blacktriangleright\blacktriangleright \mathsf{S\;=\; x\doteq A;}:$$
	در این حالت، با توجه به تعریف معنای $\mathsf{S}$، دنباله‌ی 
	$$\langle at \llbracket \mathsf{S} \rrbracket ,\rho \rangle \langle aft\llbracket\mathsf{S}\rrbracket , \rho[\mathsf{x}\leftarrow \mathcal{A}\llbracket\mathsf{A}\rrbracket\rho] \rangle $$
	در معنای عبارت‌ دستوری به ازای هر $\rho$ وجود دارد. پس دنباله‌ای به شکل بالا با محیطی متناظر با
	$\underline{\rho}$
	هم در معنای عبارت‌ دستوری حضور دارد.
	$$\blacktriangleright\blacktriangleright \mathsf{S=\; ; \;\;}:$$	 
با توجه به معنای این عبارت‌ دستوری، دنباله‌ی زیر در معنای این عبارت‌ دستوری وجود دارد.
$$\langle at \llbracket \mathsf{S} \rrbracket ,\rho \rangle \langle aft\llbracket\mathsf{S}\rrbracket , \rho \rangle $$

$$\blacktriangleright\blacktriangleright \mathsf{S=\; if\;(B)\;S_t}:$$
	در صورتی که 
	$\mathcal{B}\llbracket\mathsf{B}\rrbracket\rho=\mathit{T}$
	برقرار باشد، دنباله‌ی 
	$$\langle at \llbracket \mathsf{S} \rrbracket , \rho \rangle
	\langle at \llbracket \mathsf{S_t} \rrbracket , \rho \rangle \pi$$
	در مجموعه‌ی معنای این عبارت‌ دستوری حضور دارد، در حالیکه، 
	$\langle at \llbracket \mathsf{S_t} \rrbracket , \rho \rangle \pi$
	داخل معنای $\mathsf{S_t}$ نیز هست. طبق \gls*{indhyp} می‌دانیم که برچسب آخرین وضعیت $\pi$ برابر است با $aft \llbracket \mathsf{S_t} \rrbracket$ که طبق تعاریف مربوط به برچسب‌ها داریم
	$aft \llbracket \mathsf{S_t} \rrbracket=aft \llbracket \mathsf{S} \rrbracket$.\\
	در صورتی که معنای عبارت بولی غلط باشد نیز، طبق تعریف، دنباله‌ی زیر در معنای عبارت‌ دستوری موجود است.
	$$\langle at \llbracket \mathsf{S} \rrbracket ,\rho \rangle \langle aft\llbracket\mathsf{S}\rrbracket , \rho \rangle $$
	
$$\blacktriangleright\blacktriangleright \mathsf{S=\; if\;(B)\;S_t\;else\;S_f}:$$
مانند حالت قبل است، منتها با این تفاوت که در صورتی که معنای عبارت بولی غلط باشد، دنباله‌ی زیر در معنای عبارت‌ دستوری حضور دارد:
	$$\langle at \llbracket \mathsf{S} \rrbracket , \rho \rangle
\langle at \llbracket \mathsf{S_f} \rrbracket , \rho \rangle \pi$$
همین ‌طور، تساوی 
$aft \llbracket \mathsf{S_t} \rrbracket=aft \llbracket \mathsf{S} \rrbracket=aft \llbracket \mathsf{S_f} \rrbracket$
هم طبق تعریف برچسب‌ها برقرار است.

$$\blacktriangleright\blacktriangleright \mathsf{S=\; while\;(B)\;S_t}:$$
در اثبات این سمت قضیه، این حالت پیچیده ترین حالت است و در واقع تنها حالتی است که در اثبات آن به فرض قضیه احتیاج داریم! همان طور که پیش‌تر گفتیم، معنای عبارت دستوری حلقه به کمک یک تابع تعریف می‌شود. معنای حلقه کوچکترین نقطه‌ثابت این تابع است، در حالیکه، این تابع وقتی روی یک مجموعه از ردهای پیشوندی اعمال شود، تاثیرات یک بار اجرای عبارت‌(های) دستوری درون حلقه را روی ردهای پیشوندی درون مجموعه اعمال می‌کند.

طبق تعریف $\mathcal{F}$، مطمئن هستیم که رد پیشوندی‌ای که با محیط $\underline{\rho}$ شروع شود، در مجموعه‌ی معنای $\mathsf{S}$ وجود دارد، چون به ازای هر محیط $\dot{\rho}$(نقطه به این خاطر است که با $\rho$ خاص موجود در فرض اشتباه گرفته نشود) حالت 
$\langle at \llbracket \mathsf{S} \rrbracket, \dot{\rho} \rangle$
 در هر اعمال تابع $\mathcal{F}$ روی هر مجموعه‌ی دلخواه وجود دارد. چون معنای $\mathsf{S}$ را به عنوان کوچک‌ترین نقطه‌ثابت $\mathcal{F}$ در نظر گرفته‌ایم، مطمئن هستیم که مجموعه‌ای که کوچکترین نقطه‌ثابت تابع $\mathcal{F}$ است، شامل رد پیشوندی 
 $\langle at \llbracket \mathsf{S} \rrbracket, \rho \rangle$
 می‌شود. این رد پیشوندی، با اجرای $\mathcal{F}$ تحت تاثیر قرار می‌گیرد. اگر معنای $\mathsf{B}$ در یکی از اعمال های $\mathcal{F}$ غلط باشد، رد پیشوندی
$ \langle at \llbracket \mathsf{S} \rrbracket, \rho \rangle \pi \langle aft \llbracket \mathsf{S} \rrbracket , \rho' \rangle$ 
در معنای برنامه قرار خواهد گرفت. بنابراین، اجرای عبارت‌ دستوری با این محیط اولیه توقف پذیر است. می‌دانیم که طبق تعریف تابع، چیزی به انتهای این رد پیشوندی اضافه نمی‌شود. از طرف دیگر، با این محیط اولیه، با توجه به تعریف، رد پیشوندی دیگری وجود ندارد که طولانی‌تر از رد پیشوندی مورد اشاره باشد. 

در حالت دیگر، اگر فرض کنیم هیچ گاه به حالتی نمی‌رسیم که در آن معنای $\mathsf{B}$ غلط باشد هم با فرض مسئله به تناقض می‌خوریم، چون در آن صورت تابع $\mathcal{F}$ مدام به دنباله‌‌هایی که با محیط $\rho$ شروع می‌شوند، می‌افزاید و این یک دنباله‌ی نامتناهی را خواهد ساخت. در صورتی که معنای $\mathsf{B}$ هیچ گاه صحیح نباشد، حداقل حالت 
$\langle at \llbracket \mathsf{S_t} \rrbracket , \rho'' \rangle$
به آخر دنباله‌های پیشین اضافه خواهد شد و از این جهت مطمئن هستیم که دنباله‌ی نامتناهی گفته شده در معنای عبارت‌ دستوری حضور خواهد داشت. پس با این تفاصیل، این حالت هم ثابت می‌شود.

$$\blacktriangleright\blacktriangleright \mathsf{S=\; break;}:$$

مهم‌ترین چیزی که در مورد برچسب‌ها در این عبارت‌ دستوری\footnote{
در تعریف تابع $aft$ روی برچسب‌ها در \cite{calcul} تعریف برای این عبارت‌ دستوری مشخص نیست! در \cite{cousotbook} هم در مورد این برچسب‌ها بحث شده است. در آنجا، گفته شده است که در مورد بخش هایی از تعریف این توابع، که تعریف مشخصی ارائه نشده است، برداشت آزاد است! ما در اینجا سعی داریم، معقول ترین برداشتی که نسبت به درکمان از این کار می‌توانیم داشته باشیم را بیان کنیم.
} باید برقرار باشد، این است که اگر این عبارت‌ دستوری بخشی از $\mathsf{S_t}$ در حلقه‌ی زیر باشد،
$$\mathsf{S'=\; while\; (B)\; S_t}$$
در این صورت، تساوی 
$aft \llbracket \mathsf{S'} \rrbracket = brk-to \llbracket \mathsf{S_t} \rrbracket $
را طبق تعریف خواهیم داشت. انتظار می‌رود که\break 
$aft \llbracket \mathsf{break;} \rrbracket = aft \llbracket \mathsf{S'} \rrbracket$ 
باشد. اینکه اجرای برنامه پس از اجرای $\mathsf{break;}$ بعد از عبارت دستوری $\mathsf{S'}$ پی گرفته شود، انتظار معقولی است از صوری‌سازی‌ای که در حال توصیف گردش کار کامپیوتری است. البته در نظر گرفته شود که فرض کرده‌ایم که $\mathsf{S'}$ داخلی ترین حلقه‌ای است که $\mathsf{break;}$ درون آن جای دارد.

 از پس این فرض‌های ما 
$aft \llbracket \mathsf{break;} \rrbracket = break-to \llbracket \mathsf{S_t} \rrbracket$ 
نتیجه می‌شود و طبق معنای عبارت دستوری 
$\mathsf{break;}$،
رد پیشوندی زیر در معنای این عبارت‌ دستوری وجود دارد،
$$\langle at \llbracket \mathsf{break;} \rrbracket , \rho \rangle
\langle aft \llbracket \mathsf{break;} \rrbracket , \rho \rangle$$
 که نشانه‌ی توقف است.
 

$$\blacktriangleright\blacktriangleright \mathsf{S=\; \{ Sl'' \}}:$$
	در این صورت، توقف پذیری $\mathsf{Sl''}$ را از فرض استقرای، استقرای قبلی( روی لیست عبارت‌های دستوری) داریم، پس $\mathsf{\{Sl''\}}$ هم توقف پذیر است.
	
	
در اینجا، اثبات این طرف قضیه به پایان می‌رسد.

($\Leftarrow$)
دوباره باید از استقرا روی ساختار برنامه‌ها استفاده کنیم و دوباره چون هر برنامه‌ای مساوی با یک لیست از عبارت‌های دستوری است، ابتدا از استقرا روی ساختار لیست عبارت‌های و در دل آن روی ساختار خود عبارت‌های دستوری استفاده کنیم.

در این اثبات به غیر از یک حالت ساختار یک عبارت‌ دستوری، که عبارت‌ دستوری حلقه است، هر آنچه در مورد اثبات طرف راست قضیه گفتیم، به ما حکم را بدون نیاز به فرض نشان می‌دهد. بنابراین، فقط در مورد اثبات این یک حالت بحث می‌کنیم: 
$$\blacktriangleright\blacktriangleright \mathsf{S=\; while\; (B)\; S_t}:$$

اگر فرض کنیم این عبارت‌ دستوری به ازای محیط $\rho$ در حالت اول متوقف شده است، در واقع فرض کرده‌ایم که در معنای این عبارت‌ دستوری رد پیشوندی 
$\langle at \llbracket \mathsf{S} \rrbracket , \rho \rangle \pi \langle aft \llbracket \mathsf{S} \rrbracket , \rho' \rangle \pi'$ 
وجود دارد.
باید ثابت کنیم، به ازای $\pi'$ دلخواه، اگر رد پیشوندی
$\langle at \llbracket \mathsf{S} \rrbracket , \rho \rangle \pi \langle aft \llbracket \mathsf{S} \rrbracket , \rho' \rangle \pi'$ 
داخل 
$\mathcal{S^*}\llbracket \mathsf{S} \rrbracket$
وجود داشته باشد، آنگاه
$\pi'=\epsilon$
برقرار است.

اگر برچسب
$aft \llbracket \mathsf{S} \rrbracket$
در یک وضعیت از رد پیشوندی‌ای که گفتیم، حضور داشته باشد، یعنی در یک دور از اجرای حلقه، عبارت بولی لزوما معنای غلط داشته است. این باعث شده است که یک وضعیت شامل برچسب $aft$ به یک رد پیشوندی چسبانده بشود و درنتیجه، این رد پیشوندی ساخته شده است. از طرفی دیگر هم می‌دانیم که وقتی معنای عبارت بولی حاضر در ساختار حلقه غلط شده باشد، دیگر به ردهای پیشوندی داخل معنای حلقه چیزی اضافه نمی‌شود. بنابراین حالت ممکنی جز $\pi'=\epsilon$ باقی نمی‌ماند. 
\end{proof}

پس با توجه به آنچه گفتیم، می‌توانیم با خیال راحت توقف پذیری یک برنامه با یک محیط اولیه را معادل متناهی بودن همه‌ی ردهای پیشوندی‌ای بدانیم که با محیط متناظر با آن محیط اولیه شروع شده‌اند.
اگر در صورت ارائه شده از وارسی مدل عبارت منظم $\mathsf{R}$ را با عبارت منظم $\varepsilon$ جایگزین کنیم، داریم:
$$\mathsf{P},\underline{\rho} \models \mathsf{R}
\Leftrightarrow
(\{\underline{\rho}\}\times \mathcal{S}^* \llbracket \mathsf{P}\rrbracket) \subseteq 
\mathsf{prefix} (\mathcal{S}^r \llbracket \varepsilon \bullet (?:\mathit{T})^*\rrbracket)
=\mathsf{prefix} (\mathcal{S}^r \llbracket (?:\mathit{T})^*\rrbracket).$$
طبق تعریف معنای عبارات منظم، هر رد پیشوندی متناهی‌ای داخل مجموعه‌ی سمت راستی رابطه‌ی زیرمجموعه بودن قرار می‌گیرد. این یعنی اگر الگوریتمی برای بررسی 
$\mathsf{P} , \underline{\rho} \models \mathsf{R} $
داشته باشیم، این \gls*{algorithm} می‌تواند تشخیص دهد که آیا برنامه‌ی $\mathsf{P}$ با محیط اولیه‌ی $\underline{\rho}$ متوقف می‌شود یا خیر که الگوریتمی است برای مسئله‌ی توقف پذیری، مسئله‌ای که تصمیم ناپذیر است. بنابراین، چنین الگوریتمی نباید وجود داشته باشد که یعنی پیاده سازی‌ای برای روشی که در حال بیانش هستیم وجود ندارد. ادامه‌ی کار روی همین تعریف پیش می‌رود و دو صورت دیگر هم که قرار است ساختارمندتر باشند، در نهایت با این صورت معادل‌اند و در نتیجه تصمیم ناپذیرند.














\chapter{وارسی مدل منظم}
در این فصل، به صورتی ساختارمندتر برای روش وارسی مدل می‌رسیم. ساختاری که در این فصل به صورت روش وارسی مدل اضافه می‌شود، ساختار عبارات منظم است و نام این صورت هم \gls{regmodelc} است. از این رو، پیش از اینکه به بیان وارسی مدل منظم بپردازیم، نیاز داریم که ابتدا، به بررسی و تعریف برخی خواص عبارات منظم بپردازیم که در ادامه، برای صورت وارسی مدل منظم مورد نیاز هستند.

\section{برخی مفاهیم}
در این بخش، ابتدا مفهوم هم‌ارز بودن را برای عبارت‌های منظم تعریف می‌کنیم، سپس به سراغ تعریف دو تابع 
$\mathsf{dnf}$
و 
$\mathsf{fstnxt}$
می‌رویم. 
\subsection{هم‌ارزی عبارت‌های منظم}
‌\gls{equivalence} بین دو عبارت منظم را با برابر بودن معنای آن دو تعریف می‌کنیم.
\begin{defn}
	(هم‌ارزی عبارات منظم):
	دو عبارت منظم
	$\mathsf{R}_1$
	و
	$\mathsf{R_2}$
	 را هم‌ارز می‌گوییم، اگر و تنها اگر شرط زیر برقرار باشد:
	 $$\mathcal{S}^r \llbracket \mathsf{R_1} \rrbracket = \mathcal{S}^r \llbracket \mathsf{R_2} \rrbracket$$ 
	 و آن را با 
	 $\mathsf{R_1} \Bumpeq \mathsf{R_2}$
	 نمایش می‌دهیم.
\end{defn}

\begin{thm}
	هم‌ارزی $\Bumpeq$ تعریف شده روی مجموعه‌ی عبارت‌های منظم یک رابطه‌ی هم‌ارزی است.
\end{thm}
\begin{proof}
	برای هر عبارت منظم $\mathsf{R}$ داریم:
	$$\mathcal{S}^r \llbracket \mathsf{R} \rrbracket = 
	\mathcal{S}^r \llbracket \mathsf{R} \rrbracket \Rightarrow
	\mathsf{R} \Bumpeq \mathsf{R}$$
	پس $\Bumpeq$ \gls*{reflexive} است.
	
	اگر 
	$\mathsf{R_1 , R_2} \in \mathbb{R}$
	آنگاه داریم:
	$$\mathcal{S}^r \llbracket \mathsf{R_1} \rrbracket =
	\mathcal{S}^r \llbracket \mathsf{R_2} \rrbracket \rightarrow
	\mathcal{S}^r \llbracket \mathsf{R_2} \rrbracket =
	\mathcal{S}^r \llbracket \mathsf{R_1} \rrbracket \Rightarrow
	\mathsf{R_1} \Bumpeq \mathsf{R_2} \rightarrow
	\mathsf{R_2} \Bumpeq \mathsf{R_1}$$
	پس $\Bumpeq$ \gls*{symmetric} است.
	
	اگر
	 $\mathsf{R_1,R_2,R_3} \in \mathbb{R}$
	 	داریم:
	 $$\mathcal{S}^r \llbracket \mathsf{R_1} \rrbracket =
	 \mathcal{S}^r \llbracket \mathsf{R_2} \rrbracket \land
	 \mathcal{S}^r \llbracket \mathsf{R_2} \rrbracket =
	 \mathcal{S}^r \llbracket \mathsf{R_3} \rrbracket \rightarrow
	 \mathcal{S}^r \llbracket \mathsf{R_1} \rrbracket =
	 \mathcal{S}^r \llbracket \mathsf{R_3} \rrbracket$$
	 $$\Rightarrow
	 \mathsf{R_1} \Bumpeq \mathsf{R_2} \land
	 \mathsf{R_2} \Bumpeq \mathsf{R_3} \rightarrow
	 \mathsf{R_1} \Bumpeq \mathsf{R_3}$$
	 پس $\Bumpeq$ \gls*{transitive} است.
	 
	 پس $\Bumpeq$ یک رابطه‌ی هم‌ارزی است.
\end{proof}


\subsection{فرم نرمال فصلی}

	 در این فصل، برای معرفی صورت جدید وارسی مدل نیاز به مفهوم فرم \gls{dnf} که در تعریف زیر آمده است.
	 \begin{defn}
	 	(فرم نرمال فصلی): عبارت منظم 
	 	$\mathsf{R} \in \mathbb{R}$
	 	را یک فرم نرمال فصلی می‌گوییم، اگر و تنها اگر عبارت‌های منظم بدون انتخاب
	 	$\mathsf{R_1 , R_2, ..., R_n} \in \mathbb{R}^{\nmid}$
	 	چنان وجود داشته باشند که 
	 	$\mathsf{R= R_1 + R_2 + ... R_n}$.
	 \end{defn}
 در تعریف بالا، = با 
 $\Bumpeq$
فرق می‌کند که منظور از آن همان تساوی نحوی است.
 
 در ادامه می‌خواهیم، یک تابع به اسم $\mathsf{dnf}$ تعریف کنیم که یک عبارت منظم $\mathsf{R}$ را می‌گیرد و عبارت منظم هم‌ارز به فرم نرمال فصلی تحویل می‌دهد.
 
  
  \begin{defn}
  	(تابع $\mathsf{dnf}$): تابع $\mathsf{dnf}$ روی عبارات منظم به شکل زیر تعریف می‌شود:
  	
  	$$\blacktriangleleft\mathsf{dnf}(\varepsilon)=\varepsilon$$
  	$$\blacktriangleleft\mathsf{dnf}(\mathsf{L:B})=\mathsf{L:B}$$
  	$$\blacktriangleleft\mathsf{dnf}(\mathsf{R_1 R_2})= \mathsf{\Sigma_{i=1}^{n_1} \Sigma_{j=1}^{n_2} R_1^i R_2^j }$$
  	$$\mathsf{where\;R_1^1 + R_1^2 + ... + R_1^{n_1} = dnf(R_1)\;and\; R_2^1 + R_2^2 + ... + R_2^{n_2}= dnf(R_2)}$$
   $$\blacktriangleleft\mathsf{dnf (R_1+R_2)=dnf(R_1)+dnf(R_2)}$$
   $$\blacktriangleleft\mathsf{dnf (R^*)}= \mathsf{((R_1)^* (R_2)^* ... (R_n)^*)^*}$$
   $$\mathsf{where\;dnf(R)=R^1+R^2+...+R^n}$$
   $$\blacktriangleleft\mathsf{dnf(R^+)=dnf(RR^*)}$$
   $$\blacktriangleleft\mathsf{ dnf((R)) = ( dnf(R) ) }$$  
   	
  \end{defn}

\begin{thm}
	اگر $\mathsf{R} \in \mathbb{R}$ آنگاه $\mathsf{dnf(R)}$ یک ترکیب نرمال فصلی است.
\end{thm}
\begin{proof}
از استقرا روی ساختار $\mathsf{R}$ استفاده می‌کنیم.
$$\blacktriangleright \mathsf{R}=\varepsilon:$$
$$\mathsf{dnf(\varepsilon)}=\varepsilon,$$
که یک فرم نرمال فصلی است.


$$\blacktriangleright \mathsf{R=L:B}:$$
$$\mathsf{dnf(\mathsf{L:B})}=\mathsf{L:B}$$
که یک فرم نرمال فصلی است.


$$\blacktriangleright \mathsf{R=R_1R_2}:$$
فرض استقرا این است که 
$\mathsf{dnf(R_1)=R_1^1+R_1^2+...+R_1^n}$
و
$\mathsf{dnf(R_2)=R_2^1+R_2^2+...+R_2^n}$.
درحالیکه، $\mathsf{dnf(R_1)}$ و $\mathsf{dnf(R_2)}$ ترکیب نرمال فصلی هستند، یعنی هر $\mathsf{R_1^i}$ و هر $\mathsf{R_2^j}$ عضو $\mathbb{R^{\nmid}}$ است.
طبق تعریف، داریم:
$$\mathsf{dnf}(\mathsf{R_1 R_2})=\mathsf{\Sigma_{i=1}^{n_1}\Sigma_{j=1}^{n_2} R_1^i R_2^j}$$
که طرف راست گزاره‌ی بالا یک ترکیب نرمال فصلی است، چون هر  
$\mathsf{R_1^i R_2^j}$
یک عضو از $\mathbb{R}^\nmid$ است.

$$\blacktriangleright \mathsf{R=R_1+R_2}:$$
فرض استقرا این است که $\mathsf{dnf(R_1)}$ و $\mathsf{dnf(R_2)}$ ترکیب نرمال فصلی هستند. طبق تعریف،
$\mathsf{dnf(R_1+R_2)}$
برابر با
$\mathsf{dnf(R_1)+dnf(R_2)}$
است. بنابراین، این عبارت منظم یک ترکیب نرمال فصلی است.

$$\blacktriangleright \mathsf{R=R_1^*}:$$
طبق فرض استقرا، داریم که $\mathsf{dnf(R_1)}$ یک ترکیب نرمال فصلی است. همین طور طبق تعریف $\mathsf{dnf}$ داریم:
$$\mathsf{dnf(R_1^*)= ((R_1^1)^* (R_1^2)^* ... (R_1^n)^*)}$$
که:
$$\mathsf{dnf(R_1)=R_1^1+R_1^2+...+R_1^n}$$
واضح است که $\mathsf{((R_1^1)^* (R_1^2)^* ... (R_1^n)^*)}$ یک فرم نرمال فصلی است، چون می‌دانیم که در هیچ کدام از $\mathsf{R_1^i}$ ها عملگر $+$ وجود ندارد و عملگر $ ^*$ و عملگر چسباندن هم تغییری در این وضع ایجاد نمی‌کنند.

$$\blacktriangleright \mathsf{R=R_1^+}:$$
طبق تعریف، داریم:
$$\mathsf{dnf(R_1^+)=dnf(R_1 R_1^*)}$$
$$\mathsf{dnf(R_1^*)= ((R_1^1)^* (R_1^2)^* ... (R_1^n)^*)}$$
توجه شود 
$\mathsf{R'=dnf(R_1^*)}$
 عضو 
 $\mathbb{R^\nmid}$
است. بعلاوه، از فرض استقرا داریم:
$$\mathsf{R_1= R_1^1 + ... + R_1^n}$$
پس با توجه به تعریف $\mathsf{dnf}$ برای عملگر چسباندن و پخش پذیری چسباندن نسبت به انتخاب، خواهیم داشت:
$$\mathsf{dnf(R_1^+) = \Sigma_{i=1}^n R_1^i R'}$$

$$\blacktriangleright \mathsf{R=(R_1)}:$$
طبق تعریف، داریم:
$$\mathsf{dnf((R_1))=(dnf(R_1))}$$
طبق فرض استقرا 
$\mathsf{dnf(R_1)}$
یک ترکیب نرمال فصلی است، بنابراین 
$\mathsf{(dnf(R_1))=R'} \in \mathbb{R^\nmid}$ 
هم یک ترکیب نرمال فصلی خواهد بود.

\end{proof}

گزاره‌ی دیگری که برای اثبات مانده است، برقرار بودن 
$\mathsf{R \Bumpeq dnf(R) }$
است. برای اثبات آن باید ابتدا قضیه‌ی زیر را اثبات کنیم.
\begin{thm}
	برای هر دو عبارت منظم 
	$\mathsf{R_1 , R_2} \in \mathbb{R}$
	داریم:
	$$\mathsf{(R_1 + R_2)^* \Bumpeq (R_1^* R_2^*)^*}.$$
\end{thm} 
\begin{proof}
	ثابت می‌کنیم:
	$$1.\mathcal{S}^r \llbracket \mathsf{(R_1 + R_2)}^* \rrbracket \subseteq \mathcal{S}^r \llbracket \mathsf{(R_1^* R_2^*)}^* \rrbracket,$$
	$$2.\mathcal{S}^r \llbracket \mathsf{(R_1^* R_2^*)}^* \rrbracket \subseteq \mathcal{S}^r \llbracket  \mathsf{(R_1 + R_2)}^*\rrbracket.$$ 
	
	اثبات $1$:
	$$\mathcal{S}^r \llbracket \mathsf{(R_1 + R_2)}^* \rrbracket = \bigcup_{n\in\mathbb{N}}
	\mathcal{S}^r \llbracket\mathsf{(R_1 + R_2)}^n \rrbracket
	$$
	$$\Rightarrow \exists m\; (\pi=\pi_1\pi_2...\pi_m \;\land\;\forall i \; (1\leq i \leq m \rightarrow \langle\underline{\rho},\pi_i \rangle \in \mathcal{S}^r \llbracket \mathsf{R_1+R_2}\rrbracket))$$
	$$\Rightarrow\forall i \; (1\leq i \leq m \rightarrow \langle\underline{\rho},\pi_i  \rangle \in \mathcal{S}^r \llbracket \mathsf{R_1}\rrbracket \cup \mathcal{S}^r \llbracket \mathsf{R_2}\rrbracket)$$
	$$\Rightarrow\forall i \; (1\leq i \leq m \rightarrow \langle\underline{\rho},\pi_i  \rangle \in \mathcal{S}^r \llbracket \mathsf{R_1}\rrbracket \lor \langle\underline{\rho},\pi_i  \rangle \mathcal{S}^r \llbracket \mathsf{R_2}\rrbracket)$$
	$$\Rightarrow\forall i \; (1\leq i \leq m \rightarrow \langle\underline{\rho},\pi_i  \rangle \in \mathcal{S}^r \llbracket \mathsf{R_1^* R_2^*}\rrbracket)$$
	$$\Rightarrow\langle\underline{\rho},\pi  \rangle \in \mathcal{S}^r \llbracket \mathsf{(R_1^* R_2^*)^m}\rrbracket \subseteq \mathcal{S}^r \llbracket \mathsf{(R_1^* R_2^*)^*}\rrbracket.$$
	
	اثبات $2$:
	$$\langle\underline{\rho},\pi  \rangle \in \mathcal{S}^r \llbracket \mathsf{(R_1^* R_2^*)^*}\rrbracket \Rightarrow \exists m \; (\langle\underline{\rho},\pi  \rangle \in \mathcal{S}^r \llbracket \mathsf{(R_1^* R_2^*)^m}\rrbracket)$$
	$$\Rightarrow \pi=\pi_1\pi_2...\pi_m \;\land\;\forall i \; (1\leq i \leq m \rightarrow \langle\underline{\rho},\pi_i \rangle \in \mathcal{S}^r \llbracket \mathsf{R_1^* R_2^*}\rrbracket)$$
	$$\Rightarrow\forall i \; (1\leq i \leq m \rightarrow \pi_i=\pi_i^1 \pi_i^2 \land  \langle\underline{\rho},\pi_i^1  \rangle \in \mathcal{S}^r \llbracket \mathsf{R_1}^*\rrbracket \land \langle\underline{\rho},\pi_i^2  \rangle \in \mathcal{S}^r \llbracket \mathsf{R_2}^*\rrbracket)$$
	$$\Rightarrow\forall i \; (1\leq i \leq m \rightarrow \pi_i=\pi_i^1 \pi_i^2 \land  \langle\underline{\rho},\pi_i^1  \rangle \in \mathcal{S}^r \llbracket \mathsf{(R_1+R_2)}^*\rrbracket \land \langle\underline{\rho},\pi_i^2  \rangle \in \mathcal{S}^r \llbracket \mathsf{(R_1+R_2)}^*\rrbracket)$$
	$$\Rightarrow\forall i \; (1\leq i \leq m \rightarrow \langle\underline{\rho},\pi_i  \rangle \in \mathcal{S}^r \llbracket \mathsf{(R_1 + R_2)^*} \bullet \mathsf{(R_1 + R_2)^*}\rrbracket)$$
	$$\Rightarrow\forall i \; (1\leq i \leq m \rightarrow \langle\underline{\rho},\pi_i  \rangle \in \mathcal{S}^r \llbracket \mathsf{(R_1 + R_2)^*}\rrbracket)$$	
	$$\Rightarrow \langle\underline{\rho},\pi  \rangle \in \mathcal{S}^r \llbracket \mathsf{((R_1 + R_2)^*)^n}\rrbracket\Rightarrow \langle\underline{\rho},\pi  \rangle \in \mathcal{S}^r \llbracket \mathsf{(R_1 + R_2)^*}\rrbracket.$$
\end{proof}


به عنوان نتیجه‌ای از قضیه‌ی بالا، می‌توانیم با استفاده از یک برهان ساده به کمک استقرا روی اعداد طبیعی، حکم بالا را به جای ۲ برای تعداد دلخواه متنهاهی‌ای از عبارات منظم اثبات کنیم. در ادامه در واقع از این حکم در اثبات استفاده شده است.

\begin{thm}
	برای هر
	$\mathsf{R} \in \mathbb{R}$ 
	داریم:
	$$\mathsf{dnf(R) \Bumpeq R}$$
\end{thm}

\begin{proof}
	این اثبات با استقرا روی ساختار $\mathsf{R}$ انجام می‌شود.
	توجه شود که در هر حالت از استقرا اگر عبارات منظم 
	$\mathsf{R}_1$
	و
	$\mathsf {R}_2$
	در ساختار $\mathsf{R}$ حضور داشته باشند، در مورد آن‌ها فرض گرفته‌ایم که 
	$\mathsf{dnf(R_1)=R_1^1+R_1^2+...+R_1^n}$
	و
	$\mathsf{dnf(R_2)=R_2^1+R_2^2+...+R_2^m}$.
		
	$$\blacktriangleright \mathsf{R=\varepsilon:}$$
	$$\mathsf{dnf}(\varepsilon)=\varepsilon \Rightarrow 
		\mathcal{S}^r \llbracket \mathsf{dnf(\varepsilon)} \rrbracket =
		\mathcal{S}^r \llbracket \varepsilon \rrbracket.$$
	
	$$\blacktriangleright \mathsf{R=L:B\;:}$$
	$$\mathsf{dnf(L:B)=L:B} \Rightarrow 
	\mathcal{S}^r \llbracket \mathsf{dnf(L:B)} \rrbracket =
	\mathcal{S}^r \llbracket \mathsf{L:B} \rrbracket.$$
	
	$$\blacktriangleright \mathsf{R=R_1 R_2\;:}$$
	برای اثبات این حالت، باید دو گزاره‌ی زیر را ثابت کنیم:
	$$1.\mathcal{S}^r \llbracket \mathsf{R_1 R_2} \rrbracket \subseteq
	  \mathcal{S}^r \llbracket \mathsf{dnf(R_1 R_2)} \rrbracket,$$
	$$2.\mathcal{S}^r \llbracket \mathsf{R_1 R_2} \rrbracket \supseteq
	\mathcal{S}^r \llbracket \mathsf{dnf(R_1 R_2)} \rrbracket.$$
	
	
	اثبات $1$:
	
	فرض می‌کنیم که 
	$\langle \underline{\rho} , \pi \rangle$ 
	یک عضو دلخواه از 
	$\mathcal{S}^r \llbracket \mathsf{dnf(R_1 R_2)} \rrbracket$
	باشد.
	چون 
	$\mathsf{dnf(R_1R_2)=} \mathsf{\Sigma_{i=1}^{n_1} \Sigma_{j=1}^{n_2} R_1^i R_2^j }$،
	 پس داریم:
	$$\exists k_1,k_2\;
	(\pi \in \mathcal{S}^r \llbracket \mathsf{R_1^{k_1} R_2^{k_2}} \rrbracket)
	$$
	$$\Rightarrow
	\exists \pi_1, \pi_2 \; ( \; \pi=\pi_1 \pi_2 , 
	\langle \underline{\rho} , \pi_1 \rangle \in \mathcal{S}^r \llbracket \mathsf{R_1^{k_1}} \rrbracket,
	\langle \underline{\rho} , \pi_2 \rangle \in \mathcal{S}^r \llbracket \mathsf{R_2^{k_2}} \rrbracket)
	$$
	در این صورت، داریم:
	$$\mathcal{S}^r \llbracket \mathsf{R_1^{k_1}} \rrbracket \subseteq
	\mathcal{S}^r \llbracket \mathsf{R_1} \rrbracket,\;\;
	\mathcal{S}^r \llbracket \mathsf{R_2^{k_2}} \rrbracket \subseteq
	\mathcal{S}^r \llbracket \mathsf{R_2} \rrbracket$$
	بنابراین:
	$$
	\langle \underline{\rho} , \pi_1 \pi_2 \rangle =
	\langle \underline{\rho} , \pi \rangle \in 
	\mathcal{S}^r \llbracket \mathsf{R_1 R_2} \rrbracket.$$ 
	
	اثبات $2$: فرض کنید
	
	$$\langle \underline{\rho} ,\pi \rangle \in \mathcal{S}^r \llbracket \mathsf{R_1 R_2} \rrbracket$$
	آنگاه
	$$\exists \pi_1, \pi_2 ( \pi = \pi_1 \pi_2\;s.t.\;
	\langle \underline{\rho} , \pi_1 \rangle \in \mathcal{S}^r \llbracket \mathsf{R_1} \rrbracket
	,\langle \underline{\rho} , \pi_2 \rangle \in \mathcal{S}^r \llbracket \mathsf{R_2} \rrbracket).$$
	طبق فرض استقرا، داریم:
	$$\mathcal{S}^r \llbracket \mathsf{R_1} \rrbracket=
	\mathcal{S}^r \llbracket \mathsf{dnf(R_1)} \rrbracket, \;\;
	\mathcal{S}^r \llbracket \mathsf{R_2} \rrbracket=
	\mathcal{S}^r \llbracket \mathsf{dnf(R_2)} \rrbracket
	$$
	$$\Rightarrow
	\exists k_1,k_2 (\langle \underline{\rho} , \pi_1 \rangle \in 
	\mathcal{S}^r \llbracket \mathsf{R_1^{k_1}} \rrbracket,
	\langle \underline{\rho} , \pi_2 \rangle \in 
	\mathcal{S}^r \llbracket \mathsf{R_2^{k_2}} \rrbracket)$$
	$$\Rightarrow
	\langle \underline{\rho} , \pi \rangle \in
	\mathcal{S}^r \llbracket \mathsf{R_1^{k_1} R_2^{k_2}} \rrbracket
	\subseteq \mathcal{S}^r \llbracket \mathsf{dnf(R_1 R_2)} \rrbracket.
	$$ 
	
	ادامه‌ی اثبات اصلی:
	$$\blacktriangleright \mathsf{R=R_1 + R_2:}$$	
	$$\mathcal{S}^r \llbracket \mathsf{dnf(R_1+R_2)} \rrbracket=
	\mathcal{S}^r \llbracket \mathsf{dnf(R_1)+dnf(R_2)} \rrbracket$$
	$$=\mathcal{S}^r \llbracket \mathsf{dnf(R_1)} \rrbracket \cup
	\mathcal{S}^r \llbracket \mathsf{dnf(R_2)} \rrbracket.$$
	طبق فرض استقرا عبارت بالایی با 
	$\mathcal{S}^r \llbracket \mathsf{R_1} \rrbracket \cup
	\mathcal{S}^r \llbracket \mathsf{R_2} \rrbracket$
	برابر است، پس داریم:
	$$\mathcal{S}^r \llbracket \mathsf{R_1} \rrbracket \cup
	\mathcal{S}^r \llbracket \mathsf{R_2} \rrbracket$$
	$$=\mathcal{S}^r \llbracket \mathsf{R_1+R_2} \rrbracket.$$
	
	$$\blacktriangleright \mathsf{R=R_1^*:}$$
	$$\mathcal{S}^r \llbracket \mathsf{dfn(R_1^*)} \rrbracket=\mathcal{S}^r \llbracket \mathsf{((R_1^1)^* (R_1^2)^* ... (R_1^n))^*} \rrbracket$$
	(طبق نتیجه‌ای که از قضیه‌ی قبل گرفتیم)
	$$=\mathcal{S}^r \llbracket \mathsf{(R_1^1+R_1^2+...+R_1^n)^*} \rrbracket=\mathcal{S}^r \llbracket \mathsf{(R_1)^*} \rrbracket$$
	$$=\mathcal{S}^r \llbracket \mathsf{R_1^*} \rrbracket.$$
	
	$$\blacktriangleright \mathsf{R=R_1^+:}$$
	$$\mathcal{S}^r \llbracket \mathsf{dfn(R_1^+)} \rrbracket=\mathcal{S}^r \llbracket \mathsf{dfn(R_1 R_1^*)} \rrbracket.$$
	در اینجا، عملگر چسباندن را داریم. در حالت‌های قبلی، این را نشان دادیم که چه‌ طور در حضور عملگر چسباندن، حکم برقرار می‌شود. همان اثبات را در مورد همین گزاره هم می‌توانیم، بگوییم و پس از آن، به این شکل ادامه دهیم:
	$$\mathcal{S}^r \llbracket \mathsf{dfn(R_1 R_1^*)} \rrbracket=
	\mathcal{S}^r \llbracket \mathsf{R_1 R_1^*} \rrbracket=
	\mathcal{S}^r \llbracket \mathsf{R_1^+} \rrbracket$$
	
	$$\blacktriangleright \mathsf{R=(R_1):}$$
	$$\mathcal{S}^r \llbracket \mathsf{dnf((R_1))} \rrbracket=\mathcal{S}^r \llbracket \mathsf{dnf(R_1)} \rrbracket=\mathcal{S}^r \llbracket \mathsf{R_1} \rrbracket=\mathcal{S}^r \llbracket \mathsf{(R_1)} \rrbracket.$$
\end{proof}


\subsection{سر و دم عبارات منظم}
در این بخش، تابعی را روی عبارت‌های منظم تعریف می‌کنیم که یک عبارت منظم را می‌گیرد و یک زوج مرتب از عبارت‌های منظم را تحویل می‌دهد، سپس به بیان یک قضیه در مورد این تابع می‌پردازیم. این تابع را با
$\mathsf{fstnxt}$
نشان می‌دهیم و نام آن \gls{fstnxt} است. قرار است این تابع یک عبارت منظم را بگیرد و آن را به این شکل تجزیه کند که اولین اتم موجود در عبارت منظم که انگار سر عبارت منظم است، از باقی آن که دم آن عبارت منظم می‌شود، جدا شود. تابع روی عبارات منظم تهی و عبارات منظمی که عملگر $+$ را دارند تعریف نمی‌شود.
\begin{defn}
	(تابع سر و دم): تابع سر و دم
	$\mathsf{fstnxt}:\mathbb{( \mathbb{R^+ \cap R^\nmid})} \rightarrow \mathbb{R \times R}$ 
	را به شکل زیر تعریف می‌‌کنیم:\\
	\lr{
	$
	\blacktriangleleft\mathsf{fstnxt (L:B) = \langle L:B , \varepsilon \rangle }
	$\\
	$
	\blacktriangleleft\mathsf{fstnxt (R_1 R_2) = fstnxt (R_2)}\;\;$
	$\mathsf{where\; R_1 \in \mathbb{R}_\varepsilon}$\\
	$\blacktriangleleft\mathsf{fstnxt(R_1 R_2)= \llparenthesis  R_1^n \in \mathbb{R}_\varepsilon \; \mathbb{?} \; \langle R_1^f , R_2 \rangle : 
		\langle R_1^f , R_1^n \bullet R_{2} \rangle \rrparenthesis} $\\
	$\mathsf{where\; R_1 \notin \mathbb{R}_\varepsilon and\;
		fstnxt(R_1)= \langle R_1^f , R_1^n \rangle}$\\
	$\blacktriangleleft\mathsf{fstnxt(R^+)= \llparenthesis R^n \in \mathbb{R_\varepsilon}\;?\;\langle R^f , R^* \rangle : \langle R^f , R^n \bullet R^* \rangle \rrparenthesis}\;\;$
	$\mathsf{where\; fstnxt(R)=\langle R^f , R^n \rangle}$\\
	$\blacktriangleleft\mathsf{fstnxt((R))=fstnxt(R)}$\\}
\end{defn}

از این تعریف در صورتی از وارسی مدل که در این فصل ارائه می‌شود، استفاده می‌شود. یک قضیه در آخر این بخش آمده است که مهم‌ترین نتیجه در مورد تابع سر و دم است. 
برای اثبات آن قضیه ابتدا یک نحو
\footnote{
ساختاری که تابع سر و دم روی آن تعریف شده است، با ساختاری که در صورت قضیه‌ی بعدی برای مجموعه‌ی 
$\mathbb{R}^+ \cap \mathbb{R}^\nmid$ 
آورده‌ایم،
 ریخت متفاوتی دارد و البته لزومی هم ندارد که یکی باشند. ساختاری که در صورت قضیه‌ی بعدی ارائه کرده‌ایم، در \cite{calcul} نیامده است و خودمان با هدف اثبات قضیه‌ی 3.10، این ساختار را در اینجا ارائه‌ کرده‌ایم.
}
 برای $\mathbb{R^+ \cap R^\nmid}$ می‌آوریم.
\begin{thm}
	$\mathbb{R'=R^+ \cap R^\nmid}$ توسط گرامر زیر تولید می‌شود.
	$$\mathsf{R}\:\:\: ::= \:\:\: \mathsf{L:B} \: |
	\: \mathsf{\varepsilon R_2} \: | \: \mathsf{R_1 \varepsilon} \: |
	\: \mathsf{R_1 R_2} \: | 
	\: \mathsf{R_1^+} | (\mathsf{R_1})$$
\end{thm}
\begin{proof}
	برای 
	$\mathbb{R' \subseteq R^+ \cap R^\nmid}$
	از استقرا روی ساختار بالا استفاده می‌کنیم:
	$$\blacktriangleright \mathsf{R=L:B}\;:$$
	طبق ساختار  $\mathbb{R^+}$ و $\mathbb{R^\nmid}$ داریم
	$\mathsf{L:B} \in \mathbb{R^+ \cap R^\nmid}$.

	$$\blacktriangleright \mathsf{R=\varepsilon R_2}\;:$$
	که در آن 
	$\mathsf{R_2} \in \mathbb{R^+ \cap R^\nmid}$.
	طبق ساختار $\mathbb{R^+}$ ،	$\varepsilon\mathsf{R_2} \in \mathbb{R^+}$ و طبق ساختار $\mathbb{R^\nmid}$ داریم چون 
	$\mathsf{R_2} \in \mathbb{R^\nmid}$
	و
	$\varepsilon \in \mathbb{R^\nmid}$،
	پس 
	$\mathsf{\varepsilon R_2} \in \mathbb{R^\nmid}$.
	پس
	$\mathsf{\varepsilon R_2} \in \mathbb{R^+ \cap R^\nmid}$.
	
	$$\blacktriangleright \mathsf{R=R_1 \varepsilon}\;:$$
	مشابه حالت قبل ثابت می‌شود.
	
	$$\blacktriangleright \mathsf{R=R_1 R_2}\;:$$
	طبق فرض استقرا، داریم 
	$\mathsf{R_1,R_2} \in \mathbb{R^\nmid \cap R^+}$ 
	و در هر دو ساختار عملگر چسباندن را داریم. پس این حالت هم اثبات می‌شود.
	$$\blacktriangleright \mathsf{R=R_1^+}\;:$$
	مثل حالت قبل، چون عملگر $ ^+$ در هر دو ساختار موجود است، به کمک فرض استقرا اثبات می‌شود.
	$$\blacktriangleright \mathsf{R=(R_1)}\;:$$
	مثل حالت قبل اثبات می‌شود.
	
	در اینجا، اثبات 
	$\mathbb{R' \subseteq R^+ \cap R^\nmid}$ 
	کامل می‌شود. حال به سراغ اثبات 
	$\mathbb{R^+ \cap R^\nmid \subseteq R'}$
	می‌رویم.
	
	برای اثبات این بخش، با فرض اینکه 
	$\mathsf{R} \in \mathbb{R^+ \cap R^\nmid}$،
	از استقرا روی ساختار اعضای $\mathbb{R^\nmid}$ استفاده می‌کنیم( این اثبات می‌توانست با استقرا روی ساختار اعضای $\mathbb{R^+}$ هم انجام شود).
	$$\blacktriangleright \mathsf{R=\varepsilon}\;:$$
	چون $\varepsilon \notin \mathbb{R^+ \cap R^\nmid}$، پس این حالت باطل است و در مورد آن نیازی به ارائه‌ی اثبات نیست.
	
	
	$$\blacktriangleright \mathsf{R=L:B}\;:$$
	در این صورت، طبق ساختار $\mathbb{R'}$ داریم  
	$\mathsf{R} \in \mathbb{R'}$.
	
	$$\blacktriangleright \mathsf{R=R_1 R_2}\;:$$
	اینجا هم با توجه به اینکه طبق فرض استقرا 
	$\mathsf{R_1 , R_2} \in \mathbb{R'}$،
	مثل حالت قبل چون $+$ در ساختار $\mathbb{R'}$ حضور دارد، حکم ثابت می‌شود.
	
	$$\blacktriangleright \mathsf{R=R_1^*}\;:$$
	چون به ازای هیچ عبارت منظم $\mathsf{R_1}$ای 
	$\mathsf{R_1}^*$ داخل $\mathbb{R^+}$ نمی‌افتد، پس بررسی این حالت هم نیاز نیست.
	
	
	$$\blacktriangleright \mathsf{R=R_1^+}\;:$$
	مثل عملگر $+$، با توجه به فرض استقرا و اینکه $ ^+$ در ساختار $\mathbb{R'}$ حضور دارد، این حالت هم اثبات می‌شود.
	
	
	$$\blacktriangleright \mathsf{R=(R_1)}\;:$$
	شبیه به حالت قبلی است.
	
	
	
	
	
	
	
\end{proof}





\begin{thm}
	برای هر عبارت منظم 
	$\mathsf{R} \in \mathbb{R^+ \cap R^\nmid}$،
	اگر 
	$\mathsf{fstnxt(R)=\langle L:B , R'\rangle}$
		آنگاه 
	$\mathsf{R'}\in \mathbb{R}^\nmid$
	و 
	$\mathsf{R \Bumpeq L:B\bullet R'}$.
\end{thm}
\begin{proof}
	قضیه را به استقرا روی ساختار اعضای 
	$\mathbb{R^+ \cap R^\nmid}$
	اثبات می‌کنیم.
	$$\blacktriangleright\mathsf{R=L:B\;:}$$
	در این حالت، طبق تعریف تابع سر و دم، داریم
	$\mathsf{fstnxt(R)=\langle L:B , \varepsilon \rangle}$.
	از طرف دیگر \break
	$\mathcal{S}^r\llbracket \mathsf{L:B \bullet \varepsilon} \rrbracket = 
	\mathcal{S}^r \llbracket \mathsf{L:B} \rrbracket$
	و $\varepsilon$ عضو $\mathbb{R^\nmid}$ است.
	
	$$\blacktriangleright\mathsf{R=\varepsilon R_2\;:}$$
	طبق تعریف تابع سر و دم، داریم
	$\mathsf{fstnxt(\varepsilon R_2) = fstnxt(R_2)}$.
		فرض استقرا این است که اگر 
		$\mathsf{fstnxt(R_2)=\langle L:B, R'_2 \rangle}$
	آنگاه 
	$\mathsf{R'_2}\in \mathbb{R^\nmid}$
	و
	$\mathsf{L:B \bullet R'_2} \Bumpeq  \mathsf{R_2}$.
	پس 
	$\mathsf{fstnxt(R)=\langle L:B , R'_2 \rangle}$.
	طبق فرض استقرا، 
	$\mathsf{R'_2} \in \mathbb{R^\nmid}$
	و از طرفی داریم:
	$$\mathsf{R=\varepsilon R_2 \Bumpeq R_2 \Bumpeq L:B \bullet R'_2 }$$
			
	
	
	
	
	
	$$\blacktriangleright\mathsf{R=R_1 \varepsilon\;:}$$
	در این حالت امکان ندارد $\mathsf{R_1=\varepsilon}$ باشد، چون در آن صورت خواهیم داشت، 
	$\mathsf{R=\varepsilon \varepsilon} \in \mathbb{R_\varepsilon}$
	که تناقض است، چون $\varepsilon \varepsilon$ در دامنه‌ی تابع سر و دم نیست.
	طبق تعریف سر و دم اگر داشته باشیم 
	$\mathsf {fstnxt(R_1)=\langle L:B , R'_1 \rangle}$
در این صورت، بنا به اینکه 
$\mathsf{R'_1} \in \mathbb{R_\varepsilon}$
برقرار باشد یا نباشد، دو حالت را داریم:
	$$\blacktriangleright\blacktriangleright \mathsf{R'_1} \in \mathbb{R_\varepsilon}: $$
	در این صورت، 
	$\mathsf{fstnxt(R)=\langle L:B , \varepsilon \rangle}$
برقرار است. داریم
 $\mathsf{R} \in \mathbb{( \mathbb{R^+ \cap R^\nmid})}$
  پس 
  چون
  $\mathsf{R_2}$
  زیر رشته‌ی $\mathsf{R}$ است، پس
  $\mathsf{R_2 \in} \mathbb{R^\nmid}$
  برقرار است. در این صورت، $\mathsf{R=L:B\;\varepsilon}$ نیز برقرار خواهد بود. بنابراین،
  $\mathsf{L:B\;\varepsilon \Bumpeq L:B \bullet \varepsilon}$
  هم بدیهی خواهد بود.
  $$\blacktriangleright\blacktriangleright \mathsf{R'_1} \notin \mathbb{R_\varepsilon}: $$
در این صورت، گزاره‌ی 
$\mathsf{fstnxt(R)=\langle L:B , R'_1 \bullet \varepsilon \rangle}$
طبق تعریف سر و دم برقرار است. چون 
$\mathsf{R} \in \mathbb{R^+ \cap R^\nmid}$،
پس زیر رشته‌های آن نیز عملگر $+$ را نخواهند داشت، پس 
$\mathsf{R} \in \mathbb{R^\nmid}$.
در این صورت نیز، واضح است که:
$$\mathsf{L:B \bullet R'_1 \bullet \varepsilon \Bumpeq R_1 \varepsilon = R }$$


	
	$$\blacktriangleright\mathsf{R=R_1 R_2\;:}$$
	اگر یکی از $\mathsf{R_1}$ و $\mathsf{R_2}$ مساوی $\varepsilon$ باشند، حالاتی که بررسی کرده‌ایم اتفاق می‌افتند. اگر هر دو عبارت منظم برابر با $\varepsilon$ باشند نیز به تناقض می‌خوریم، چون در این صورت دیگر در $\mathbb{R^+}$ این عبارت منظم را نداریم. پس تنها یک حالت می‌ماند و آن این است که هیچ یک از این دو عبارت منظم تهی نباشند. در این صورت، اگر فرض کنیم 
	$\mathsf{fstnxt(R_1)=\langle L:B , R'_1 \rangle}$،
	مسئله بنا به اینکه 
	$\mathsf{R'_1} \in \mathbb{R_\varepsilon}$
	برقرار هست یا خیر، به دو حالت افراز می‌شود.
	$$\blacktriangleright\blacktriangleright \mathsf{R'_1} \notin \mathbb{R_\varepsilon}: $$
	در این صورت 
	$\mathsf{fstnxt(R)=\langle L:B , R'_1 \bullet R_2}\rangle$.
	مانند استدلالی که در حالت قبلی آوردیم، خواهیم داشت 
	$\mathsf{R'_1 \bullet R_2 \in } \mathbb{R^\nmid}$.
	علاوه‌بر این، طبق فرض استقرا داریم 
	$\mathsf{ R_1 \Bumpeq L:B \bullet R'_1 }$،
	پس:
	$$\mathsf{L:B \bullet R'_1 \bullet R_2 \Bumpeq R_1 \bullet R_2 = R}$$
	(عملگر چسباندن \gls*{assoc} است). پس این حالت اثبات می‌شود.
	$$\blacktriangleright\blacktriangleright \mathsf{R'_1} \in \mathbb{R_\varepsilon}: $$
	در این صورت،
	$\mathsf{fstnxt(R)=\langle L:B , R_2 \rangle}$.
	مثل حالت‌های قبل ثابت می‌شود که $\mathsf{R_2} \in \mathbb{R^\nmid}$ و داریم:
	$$\mathsf{R = L:B \bullet R_2 \Rightarrow L:B \bullet R_2 \Bumpeq R}$$ 
	
	$$\blacktriangleright\mathsf{R=R_1^+\;:}$$
	با فرض اینکه 
	$\mathsf{fstnxt(R_1)=\langle L:B , R'_1 \rangle}$،
		بنا به اینکه 
	$\mathsf{R'_1} \in \mathbb{R_\varepsilon}$
	برقرار باشد یا نباشد، دو حالت خواهیم داشت:
	$$\blacktriangleright\blacktriangleright\mathsf{R'_1} \in \mathbb{R_\varepsilon}:$$
	در این صورت طبق تعریف تابع سر و دم،
	$\mathsf{fstnxt(R)=\langle L:B , R^*_1 \rangle}$
	را داریم. 
 چون  داخل 
	$\mathsf{R'_1}$
	و 
	$\mathsf{L:B}$
	عملگر $+$ وجود ندارد 	$\mathsf{R_1^*}$
	عضو 
	$\mathbb{R^\nmid}$ 
	خواهد بود. جای دیگری از این عبارت منظم وجود ندارد که در آن بتوان وجود این عملگر را متصور شد. همین طور، داریم:
	$$\mathsf{fstnxt(R_1) = \langle L:B , R'_1 \rangle \rightarrow
		R'_1 \in \mathbb{R^\nmid}\; \land \; L:B \bullet R'_1 \Bumpeq R_1}$$
		(گزاره‌ی بالا فرض استقراست.)
	$$\mathsf{R_1^* \Bumpeq (L:B \bullet R'_1)^* \Bumpeq (L:B \bullet \varepsilon)^* \Bumpeq (L:B)^*}$$
	(هم‌ارزی وسطی به خاطر این است که $\mathsf{R'_1}$ عضو $\mathbb{R^\nmid}$ است. اگر یکی از دو هم‌ارزی‌ دیگر هم برقرار نباشند کلا عملگر $ ^*$ خوش تعریف نخواهد بود، پس این دو هم‌ارزی باید برقرار باشند.)
	$$\Rightarrow \mathsf{L:B\bullet R_1^* \Bumpeq L:B \bullet (L:B)^* \Bumpeq 
	(L:B)^+}$$

	$$\blacktriangleright\blacktriangleright\mathsf{R'_1} \notin \mathbb{R_\varepsilon}:$$
با توجه به تعریف تابع سر و دم و فرض استقرا، در این حالت داریم 
	$\mathsf{fstnxt(R)=\langle L:B , R'_1 \bullet R_1^* \rangle}$.
	به همان دلیل حالت قبلی، می‌دانیم که $\mathsf{R_1^*}$ عضو $\mathbb{R^\nmid}$ است و طبق فرض استقرا داریم $\mathsf{R'_1} \in \mathbb{R^\nmid}$. بنابراین داریم
		$\mathsf{R'_1 \bullet R_1^*} \in \mathbb{R^\nmid}$.
		
		با استفاده از فرض استقرا داریم:
		$$\mathsf{L:B \bullet R'_1 \Bumpeq R_1}$$
		$$\Rightarrow \mathsf{L:B \bullet R'_1 \bullet R_1^* \Bumpeq R_1 R_1^* \Bumpeq R_1^+}$$
	
	$$\blacktriangleright\mathsf{R=(R_1)\;:}$$
	از فرض استقرا نتیجه می‌شود.
\end{proof}

این بخش، در این قسمت به پایان می‌رسد. حال ابزارهای کافی برای بیان صورت وارسی مدل منظم را در اختیار داریم.

\section{وارسی مدل منظم}
همان‌طور که پیش‌تر گفتیم، می‌خواهیم در این فصل یک صورت معادل با صورتی که در فصل پیش برای روش وارسی مدل آورده شده بود، ارائه کنیم. تا اینجای این فصل، صرفا به معرفی چند مفهوم که برای بیان  صورت جدید به آن‌ها احتیاج داریم، پرداخته‌ایم. در این یخش ابتدا این صورت جدید را بیان می‌‌کنیم و سپس اثبات می‌‌کنیم که صورت جدید با صورت قبلی معادل است. همان‌طور که پیش‌تر هم اشاره شد، تفاوت این دو صورت در این است که در این صورت ساختار عبارات منظم اثر دارد، در حالیکه، صورت قبلی ساختاری نداشت.
\subsection{صورت وارسی مدل منظم}

در نهایت، برای تعریف صورت وارسی مدل منظم به یک تابع به نام $\mathcal{M}$ خواهیم رسید که در ورودی‌اش، یک زوج متشکل از یک محیط اولیه و یک عبارت منظم را در کنار یک برنامه می‌گیرد و در خروجی، همه‌ی ردهای پیشوندی موجود در معنای برنامه که عبارت منظم را ارضا می‌کنند، داخل یک مجموعه بر می‌گرداند.
اما در این بین، مفهوم سازگاری یک رد پیشوندی با یک عبارت منظم چگونه مشخص می‌شود؟ این نکته‌ای است که تا به حال در مورد آن بحث نکرده‌ایم و حالا می‌خواهیم، تعریف تابع $\mathcal{M}^t$ را با این هدف به بحث وارد کنیم. البته، این تابع یک ویژگی بیشتر هم دارد. ویژگی دیگر این است که اگر عبارت منظم با رد پیشوندی سازگار نباشد، تابع به ما می‌گوید که کجای عبارت منظم ناسازگاری وجود داشته است. همین ‌طور اگر رد پیشوندی با عبارت منظم سازگار باشد، این تابع به ما نشان می‌دهد که عبارت منظم تا کجا بررسی شده است. فهمیدن این موضوع با نگاه به تعریف ساده‌تر است و البته، نباید فراموش کرد که سازگاری‌ای که داریم، هماهنگ با صورت قبلی با داشتن عملگر 
$\mathsf{prefix}$
و الحاق عبارت منظم
$\mathsf{(?:\mathit{T})^*}$
است. 
\begin{defn}
	(وارسی‌گر رد پیشوندی): تابع وارسی‌گر رد پیشوندی \break
	$\mathcal{M}^t:\mathbb{(\underline{EV} \times R^\nmid) \rightarrow  \mathfrak{S}^{+\infty}}
	\rightarrow (\mathbb{B \times R^\nmid} ) $
	با ضابطه‌ی زیر تعریف می‌شود:
	$$\blacktriangleleft\mathcal{M}^t \langle \underline{\rho} , \varepsilon \rangle \pi = 
	\langle \mathit{T} , \varepsilon \rangle$$
	(برای هر عضو دیگر $\mathbb{R_\varepsilon}$ هم حالت بالایی برقرار است. دو حالت پایینی برای عبارت‌های منظم عضو $\mathbb{R^+ \cap R^\nmid}$ هستند.)
	$$\blacktriangleleft\mathcal{M}^t \langle \underline{\rho} , \mathsf{R} \rangle \epsilon = 
	\langle \mathit{T} , \mathsf{R} \rangle$$
	$$\blacktriangleleft\mathcal{M}^t \langle \underline{\rho} , \mathsf{R} \rangle \pi = 
	\llparenthesis \langle \underline{\rho}, \langle l_1, \rho_1 \rangle \rangle \in
	\mathcal{S}^r \llbracket \mathsf{L:B} \rrbracket \; ? \; 
	\mathcal{M}^t \langle \underline{\rho}, \mathsf{R'} \rangle \pi' \; : \; 
	\langle \mathit{F}, \mathsf{R} \rangle \rrparenthesis$$
	$$\mathsf{where}\; \pi=\langle l_1,\rho_1\rangle \pi' \;\mathsf{and}\; 
	\langle \mathsf{L:B,R'}\rangle = \mathsf{fstnxt(R)}$$
	
\end{defn}

در مسیر رسیدن به تعریف $\mathcal{M}$، به معرفی یک تابع دیگر هم می‌پردازیم. این تابع را با
 $\mathcal{M}^\nmid$ 
نشان می‌دهیم. در واقع، همان کاری را که $\mathcal{M}$ قرار است به ازای همه‌ی عبارت‌های منظم انجام دهد، این تابع روی عبارت‌های منظمی که $+$ ندارند انجام می‌دهد.
\begin{defn}
	(وارسی مدل منظم محدود به $\mathbb{R^\nmid}$): تابع وارسی مدل منظم محدود به $\mathbb{R^\nmid}$ \break
	$\mathcal{M}^\nmid:\mathbb{(\underline{EV} \times R^\nmid)} \rightarrow  \mathit{P}({\mathfrak{S}^{+\infty})}
	\rightarrow \mathit{P}(\mathfrak{S}^{+\infty} \times \mathbb{R}^\nmid ) $
	با ضابطه‌ی زیر تعریف می‌شود:
	
	 $$\mathcal{M}^\nmid \langle \underline{\rho} , \mathsf{R} \rangle \Pi= \{ \langle \pi , \mathsf{R'} \rangle | \pi \in \Pi \land  
	  \mathcal{M}^t \langle \underline{\rho} , \mathsf{R} \rangle \pi =
	 \langle \mathit{T} , \mathsf{R'} \rangle \}$$
\end{defn}


حالا خود $\mathcal{M}$ را تعریف می‌کنیم. تعریف این تابع چیزی نیست جز \gls*{union} گرفتن از خروجی تابع $\mathcal{M}^\nmid$، به ازای عبارت‌های منظمی که در فرم نرمال فصلی عبارت منظم ورودی تابع حضور دارند. البته، بخشی از اطلاعاتمان از هر رد پیشوندی در هر زوجی که در خروجی $\mathcal{M}^\nmid$ وجود دارد، حذف می‌شود. به عبارت دیگر، صرفا ردهای پیشوندی را در مجموعه‌ای که خروجی $\mathcal{M}$ است، داریم. اطلاعات برای هر رد پیشوندی یک عبارت منظم است که بخشی از $\mathsf{R}$ است که سازگاری‌اش با رد پیشوندی بررسی نشده است. برای ردهای پیشوندی‌ای که در خروجی $\mathcal{M}$ حضور دارند و طولشان بیشتر یا مساوی عبارت منظم مورد بررسی است، این عبارت منظم برابر با تهی است.

\begin{defn}
(وارسی مدل منظم):

تابع وارسی مدل منظم 
$\mathcal{M}:\mathbb{(\underline{EV} \times R)} \rightarrow  \mathit{P}({\mathfrak{S}^{+\infty})}
\rightarrow (\mathbb{\underline{EV}} \times \mathit{P} (\mathfrak{S}^{+\infty} ) )$
با ضابطه‌ی زیر تعریف می‌شود:
$$\mathcal{M} \langle \underline{\rho} , \mathsf{R} \rangle \Pi= \bigcup_{i=1}^n \{  \langle \underline{\rho},\pi\rangle | \exists \mathsf{R'} \in \mathbb{R}:\; \langle \pi , \mathsf{R'} \rangle \in 
\mathcal{M}^\nmid \langle \underline{\rho}, \mathsf{R_i} \rangle \Pi \}$$
$$\mathsf{where \; dnf (R) = R_1 + R_2 + ... + R_n}$$ 



در این صورت، اگر ویژگی $\mathsf{R} \in \mathbb{R}$ در محیط اولیه‌ی $\underline{\rho}$ برای برنامه‌ی $\mathsf{P}$ برقرار باشد، می نویسیم
$$\mathsf{P} , \underline{\rho} \models_r \mathsf{R}$$
و برقرار بودن این رابطه با شرط زیر تعریف می‌شود:
$$\mathsf{P} , \underline{\rho} \models_r \mathsf{R} \iff \{\underline{\rho}\} \times \mathcal{S}^* \llbracket \mathsf{P} \rrbracket \subseteq \mathcal{M} \langle \underline{\rho} , \mathsf{R} \rangle \mathcal{S}^* \llbracket \mathsf{P} \rrbracket$$

\end{defn}

با این تعریف، در واقع زمانی می‌توانیم بگوییم، برنامه‌ی $\mathsf{P}$ ویژگی $\mathsf{R}$ را دارد \break که 
$\mathcal{M}\langle \underline{\rho} , \mathsf{R} \rangle \mathcal{S}^* \llbracket \mathsf{P} \rrbracket = 
\{\underline{\rho}\} \times \mathcal{S}^* \llbracket \mathsf{P} \rrbracket$.
در واقع، انگار تابع 
$\mathcal{M} \langle \underline{\rho} , \mathsf{R} \rangle$
مثل یک صافی روی مجموعه‌ی 
$\mathcal{S}^* \llbracket \mathsf{P} \rrbracket$
عمل می‌کند، پس خروجی این مجموعه زیر مجموعه‌ی $\{\underline{\rho}\} \times \mathcal{S}^* \llbracket \mathsf{P} \rrbracket$ است.
\begin{thm}
	برای هر برنامه‌ی $\mathsf{P}$، محیط اولیه‌ی $\underline{\rho}$ و عبارت منظم $\mathsf{R}$ داریم:
	$$\mathcal{M}\langle \underline{\rho},\mathsf{R} \rangle \mathcal{S}^* \llbracket \mathsf{P} \rrbracket \subseteq \{\underline{\rho}\} \times \mathcal{S}^* \llbracket \mathsf{P}\rrbracket.$$
\end{thm}

\begin{proof}
	اگر زوج
	$\langle \underline{\rho}, \pi \rangle$ 
	عضو 
	$\mathcal{M}\langle \underline{\rho},\mathsf{R} \rangle \mathcal{S}^* \llbracket \mathsf{P} \rrbracket $
	باشد، با فرض 
	$\mathsf{dnf(R)=R_1+R_2+...+R_n}$
	طبق تعریف 
	$\mathcal{M}$،
	 وجود دارند $\mathsf{R'} \in \mathbb{R}$ و عدد $i$ بین ۱ و $n$ که: 
	$$\langle \pi , \mathsf{R'} \rangle \in \mathcal{M}^\nmid \langle \underline{\rho},\mathsf{R_i} \rangle \mathcal{S}^* \llbracket \mathsf{P} \rrbracket$$
	که طبق تعریف $\mathcal{M^\nmid}$ یعنی: 
	$$\pi \in \mathcal{S}^* \llbracket \mathsf{P} \rrbracket \land \mathcal{M}^t \langle \underline{\rho},\mathsf{R_i}\rangle \pi = \langle \mathit{T},\mathsf{R'} \rangle$$
که از 
$\pi \in \mathcal{S}^* \llbracket \mathsf{P} \rrbracket$
می‌توان نتیجه گرفت 
$\langle \underline{\rho}, \pi \rangle \in \{\underline{\rho}\} \times \mathcal{S}^* \llbracket \mathsf{P} \rrbracket$.

\end{proof}


مجموعه‌ی معنای یک برنامه را می‌توان به مجموعه‌ای از \glspl*{class} افراز کرد که در هر یک از این دسته‌ها ردهای پیشوندی‌ای حضور دارند که وضعیت اول آن‌ها یکسان است. قاعدتا، در هر یک از این دسته‌ها باید وضعیت‌های بعدی هم، در صورت وجود، به طور موازی با یکدیگر یکسان باشند، یعنی مثلا در یک مجموعه از افراز توصیف شده، همه‌ی ردهای پیشوندی‌ای که عضو دوم دارند، عضو دومشان با هم برابر است. این گزاره در مورد عضو سوم و چهارم و غیره هم برقرار است. در هر دسته از این افراز، یک رد پیشوندی \gls{maximal} وجود خواهد داشت که توصیف تمام و کمال برنامه در اجرا با وضعیت اول مختص آن دسته است.
اگر ردهای پیشوندی با محیط اولیه‌ی یکسان به اشکال مختلفی ادامه پیدا کنند، باید زبان برنامه نویسی مورد بررسی غیرقطعی باشد. در صورتیکه، در زبان و معناشناسی این زبانی که در فصل اول تعریف کردیم، مولفه‌ای از غیرقطعی بودن حضور ندارد.

حال برای هر برنامه‌ی $\mathsf{P}$ که ویژگی $\mathsf{R}$ در مورد آن در حال بررسی است، می‌توانیم همین افراز را روی مجموعه‌ی 
$\mathcal{M} \langle \underline{\rho}, \mathsf{R} \rangle \mathcal{S}^* \llbracket \mathsf{P} \rrbracket$
در نظر بگیریم. می‌توانیم هر دسته از این افراز را متناظر با دسته‌ای در افرازی که روی 
$\mathcal{S}^* \llbracket \mathsf{P}\rrbracket$ 
توصیف کردیم بدانیم، دقیقا موقعی که وضعیت اولیه در ردهای پیشوندی موجود در دو دسته یکسان باشند.

اگر در هر دسته از این افراز روی 
$\mathcal{S}^* \llbracket \mathsf{P} \rrbracket$،
رد پیشوندی ماکسیمال این دسته در دسته‌ی متناظر در
$\mathcal{M} \langle \underline{\rho}, \mathsf{R} \rangle \mathcal{S}^* \llbracket \mathsf{P} \rrbracket$
 وجود داشته باشد، در این صورت،
$\mathcal{M} \langle \underline{\rho}, \mathsf{R} \rangle \mathcal{S}^* \llbracket \mathsf{P} \rrbracket =\{\underline{\rho}\} \times \mathcal{S}^* \llbracket \mathsf{P} \rrbracket$.
اگر دسته‌ای از افراز روی 
$\mathcal{M} \langle \underline{\rho}, \mathsf{R} \rangle \mathcal{S}^* \llbracket \mathsf{P} \rrbracket$
دارای عضو ماکسیمال متفاوتی نسبت به همان دسته روی 
$\mathcal{S}^* \llbracket \mathsf{P} \rrbracket$
باشد، این عضو ماکسیمال کوتاه‌تر از عضو ماکسیمال همان دسته در 
$\mathcal{S}^* \llbracket \mathsf{P} \rrbracket$
است و این یعنی ناسازگاری‌ای با عبارت منظم مورد بررسی وجود داشته است. محل وقوع ناسازگاری را تابع  
$\mathcal{M}^\nmid$
می‌تواند به ما بگوید. محل ناسازگاری عبارت منظمی است که زوج عضو ماکسیمال دسته‌ی مورد نظر در خروجی یکی از اعمال های 
$\mathcal{M}^\nmid$
روی بخش‌های مختلف 
$\mathsf{dnf(R)}$
است. 




\subsection{درستی و تمامیت}
حال به اثبات معادل بودن صورت جدید با صورت قبلی می‌پردازیم. در \cite{calcul} این اثبات
\footnote{نگارنده‌ی این پایان نامه، به درستی دو اثبات موجود در \cite{calcul} بسیار بد بین است! در اثبات تمامیت، برهان به شکل عجیبی بی‌ربط است و در اثبات قضیه درستی، ایرادات فنی ریزی در جزئیات وجود دارد که با تعاریف در تناقض است. از این رو برهان‌هایی که در اینجا آورده‌ایم، جدید هستند.}
 که یک قضیه‌ی دوطرفه است، تحت دو قضیه به نام‌های \gls{soundness} و \gls{completeness} آمده است. درستی به این معناست که اگر یک بررسی در صورت جدید انجام شود، نتیجه‌ای یکسان با انجام بررسی برای همان برنامه و همان عبارت منظم در صورت قبلی دارد. تمامیت نیز یعنی هر بررسی‌ای که با صورت قبلی انجام شده، نتیجه‌ی یکسانی با انجام همان بررسی در صورت جدید دارد.


\begin{thm}
	(قضیه درستی): اگر $\mathsf{P}$ یک برنامه، $\mathsf{R}$ یک عبارت منظم و $\underline{\rho}$ یک محیط اولیه باشند، آنگاه داریم:
	$$\mathsf{P},\underline{\rho} \models_r \mathsf{R} \Rightarrow
	\mathsf{P},\underline{\rho} \models \mathsf{R}$$
\end{thm}
\begin{proof}
	طبق تعریف دو صورت، باید با فرض اینکه داریم:
	$$\{\underline{\rho}\}\times\mathcal{S}^* \llbracket \mathsf{P} \rrbracket \subseteq 
	\mathcal{M}\langle \underline{\rho},\mathsf{R} \rangle \mathcal{S}^* \llbracket \mathsf{P} \rrbracket$$
	ثابت کنیم:
	$$\{\underline{\rho}\}\times\mathcal{S}^* \llbracket \mathsf{P} \rrbracket \subseteq 
	\mathsf{prefix(\mathcal{S}^r \llbracket R \bullet (?:\mathit{T})^*} \rrbracket)$$
	در این راستا، می‌توانیم گزاره‌ی زیر را ثابت کنیم که از گزاره‌ی قبلی قوی‌تر است و آن را نتیجه می‌دهد:
	$$\mathcal{M}\langle \underline{\rho},\mathsf{R} \rangle \mathcal{S}^* \llbracket \mathsf{P} \rrbracket \subseteq \mathsf{prefix(\mathcal{S}^r \llbracket R \bullet (?:\mathit{T})^*} \rrbracket)$$
	
	در مورد عبارت منظم $\mathsf{R}$، فرض می‌کنیم
	$\mathsf{dnf(R)=R_1+R_2+...+R_n}$،
	که هر $\mathsf{R}_i$ ذکر شده عضو $\mathbb{R}^\nmid$ است. 
	فرض می‌کنیم
	$\langle \underline{\rho}, \pi \rangle \in \mathcal{M}\langle \underline{\rho},\mathsf{R} \rangle \mathcal{S}^* \llbracket \mathsf{P} \rrbracket$،
	آنگاه وجود دارد $k$ و $\mathsf{R'}$ که \break
	$\langle \pi , \mathsf{R'} \rangle  \in \mathcal{M}^\nmid\langle \underline{\rho},\mathsf{R}_k \rangle \mathcal{S}^* \llbracket \mathsf{P} \rrbracket$
	که طبق تعریف $\mathcal{M}^\nmid$ خواهیم داشت:
	$$\pi \in \mathcal{S}^* \llbracket \mathsf{P} \rrbracket \land \mathcal{M}^t \langle \underline{\rho},\mathsf{R}_k \rangle \pi = \langle \mathit{T},\mathsf{R'} \rangle$$
	طرف چپ گزاره‌ی عطفی بالا در فرض بود. در ادامه‌ی کار با طرف راست این گزاره پیش می‌رویم:
	$$\mathcal{M}^t \langle \underline{\rho},\mathsf{R}_k\rangle \pi = \langle \mathit{T},\mathsf{R'} \rangle$$

در مورد $\mathsf{R}_k$ دو حالت داریم، یا 
$\mathsf{R}_k \Bumpeq \varepsilon$
برقرار است، یا اینگونه نیست( 
$\mathsf{R}_k \not\Bumpeq \varepsilon$).
$$\blacktriangleright\mathsf{R}_k \Bumpeq \varepsilon:$$
در این صورت می‌توانیم، ثابت کنیم: 
$$\mathsf{prefix(\mathcal{S}^r \llbracket \mathsf{R} \bullet (?:\mathit{T})^* \rrbracket)}=
	\{\underline{\rho}\} \times \mathfrak{S}^+$$
با توجه به پخش پذیری عملگر چسباندن روی عملگر انتخاب، که پیش‌تر ثابت کردیم، داریم:
$$\mathsf{R}\bullet(?:\mathit{T})^* \Bumpeq \mathsf{(R_1 + R_2 + ... +R_m)}\bullet(?:\mathit{T})^*$$
$$\Bumpeq \mathsf{R_1}\bullet (?:\mathit{T})^* + \mathsf{R_2} \bullet (?:\mathit{T})+...+ \mathsf{R_n} \bullet (?:\mathit{T})^*$$
چون $\mathsf{R}_k \Bumpeq \varepsilon$، داریم:
$$\mathsf{R_1}\bullet (?:\mathit{T})^* + \mathsf{R_2} \bullet (?:\mathit{T})+...+ \mathsf{R}_k \bullet (?:\mathit{T})^* +...+ \mathsf{R_n} \bullet (?:\mathit{T})^*$$
$$\Bumpeq \mathsf{R_1}\bullet (?:\mathit{T})^* + \mathsf{R_2} \bullet (?:\mathit{T})+...+ \varepsilon \bullet (?:\mathit{T})^* +...+ \mathsf{R_n} \bullet (?:\mathit{T})^*$$
و از طرف دیگر داریم:
$$\varepsilon \bullet (?:\mathit{T})^* \Bumpeq (?:\mathit{T})^* = (\{\underline{\rho}\} \times \mathfrak{S}^+)$$  

پس 
$\mathsf{prefix(\mathcal{S}^r \llbracket \mathsf{R} \bullet (?:\mathit{T})^* \rrbracket)}$
مجموعه‌ی 
$\{\underline{\rho}\} \times \mathfrak{S}^+$
را به عنوان زیرمجموعه در درون خود دارد و عضوی بیش از این هم طبق تعریفش نمی‌تواند داشته باشد، پس:
$$\mathsf{prefix(\mathcal{S}^r \llbracket \mathsf{R} \bullet (?:\mathit{T})^* \rrbracket)}=
\{\underline{\rho}\} \times \mathfrak{S}^+$$
که این گزاره نتیجه می‌دهد:
$$\langle \underline{\rho}, \pi \rangle \in \mathsf{prefix(\mathcal{S}^r \llbracket \mathsf{R} \bullet (?:\mathit{T})^* \rrbracket)}$$

$$\blacktriangleright \mathsf{R}_k \not\Bumpeq \varepsilon:$$
همان طور که پیش‌تر اشاره شد، این فرض یعنی
$\mathsf{R}_k \in \mathsf{R}^+ \cap \mathsf{R}^\nmid$.
پس مجاز هستیم از تابع سر و دم استفاده کنیم. فرض می‌کنیم 
$\mathsf{fstnxt(R}_k) = \langle \mathsf{L}_k^1 : \mathsf{B}_k^1 , \mathsf{R}_k^1 \rangle$.
همین طور فرض می‌کنیم:
$$\pi = \langle l_0,\rho_0 \rangle \langle l_1 , \rho_1 \rangle \langle l_2 , \rho_2 \rangle ... \langle l_l , \rho_l \rangle$$
قرار دهید:
$$\pi(i)=\langle l_i , \rho_i\rangle \langle l_{i+1},\rho_{i+1}\rangle,...,\langle l_l,\rho_l \rangle$$
داریم:
$$\mathcal{M}^t \langle \underline{\rho},\mathsf{R}_k \rangle \pi = \langle \mathit{T} , \mathsf{R'} \rangle \Rightarrow \forall \mathsf{R''} \in \mathbb{R}: \mathcal{M}^t \langle \underline{\rho},\mathsf{R}_k \rangle \pi \neq \langle \mathit{F},\mathsf{R}'' \rangle$$
پس لاجرم تساوی زیر برقرار است( با توجه به سر و دم $\mathsf{R}_k$):
$$\mathcal{M}^t \langle \underline{\rho}, \mathsf{R}_k \rangle \pi= \mathcal{M}^t \langle \underline{\rho}, \mathsf{R}_k^1 \rangle \pi(1)$$
بدون کاستن از کلیت( چون ممکن است $\mathsf{R}_k^1 \Bumpeq \varepsilon$)، فرض می‌کنیم کاری که انجام دادیم را می‌توانیم روی دم خروجی عبارت منظم $\mathsf{R}_k$( یعنی $\mathsf{R}_k^1$) تکرار کنیم:
$$\mathcal{M}^t \langle \underline{\rho}, \mathsf{R}_k^1 \rangle \pi(1)= \mathcal{M}^t \langle \underline{\rho}, \mathsf{R}_k^2 \rangle \pi(2)\;\;\; \mathsf{where\;fstnxt}(\mathsf{R}_k^1)=\langle \mathsf{L}_k^2 : \mathsf{B}_k^2 , \mathsf{R}_k^2 \rangle$$

باز هم بدون کاستن از کلیت، می‌توانیم فرض کنیم که این رویه را تا $h$ مرحله ادامه دادیم، یعنی:
$$\mathcal{M}^t \langle \underline{\rho},\mathsf{R}_k^{h-1} \rangle \pi(h-1)=\mathcal{M}^t\langle \underline{\rho},\mathsf{R}_k^{h} \rangle \pi(h) \;\;\;
\mathsf{where\;fstnxt}(\mathsf{R}_k^{h-1})=\langle \mathsf{L}_k^{h} : \mathsf{B}_k^{h} , \mathsf{R}_k^{h} \rangle$$
در حالیکه، 
$\mathsf{R_k^h} \Bumpeq \varepsilon$. 
اگر گزاره‌ی $\mathsf{R_k^h} \Bumpeq \varepsilon$ هیچ زمانی برقرار نشود، یعنی بتوانیم این سلسله را تا بی‌نهایت ادامه دهیم، مطمئن خواهیم بود که در معنای $\mathsf{R_k}$ حتما ردهای پیشوندی نامتناهی حضور دارند. چنین چیزی با تعریف معنای عبارات منظم در تناقض است، چون در معنای عبارات منظم رد پیشوندی نامتناهی حضور ندارد.
تا اینجا، می‌توانیم بگوییم:
$$\mathsf{R_k} \Bumpeq \mathsf{L_k^1 : B_k^1 \bullet L_k^2:B_k^2 \bullet ... \bullet L_k^h:B_k^h}$$

حال بسته به اینکه $h<l$ برقرار باشد یا نباشد، می‌توانیم مسئله را به دو حالت افراز کنیم:
$$\blacktriangleright\blacktriangleright h<l:$$
در این صورت، داریم:
$$\mathcal{M}^t \langle \underline{\rho}, \mathsf{R}_k^h \rangle \pi(h+1)=\langle \mathit{T} , \varepsilon \rangle$$

که این یعنی داریم:
$$\forall j\;(1\leq j \leq h \rightarrow \langle \underline{\rho},\langle l_j, \rho_j \rangle\rangle \in \mathcal{S}^r \llbracket \mathsf{L}_j : \mathsf{B}_j \rrbracket)$$
$$\Rightarrow \langle \underline{\rho}, \pi \rangle \in \mathcal{S}^r \llbracket \mathsf{R}_k \bullet (?:\mathit{T})^* \rrbracket \Rightarrow \langle \underline{\rho},\pi\rangle \in \mathsf{prefix(\mathcal{S}^r \llbracket R \bullet (?:\mathit{T})^*} \rrbracket).$$ 

$$\blacktriangleright\blacktriangleright h\geq l:$$
در این صورت داریم:
$$\mathcal{M}^t \langle\underline{\rho},\mathsf{R_k}\rangle \pi = \langle \mathit{T}, \mathsf{R}_k^l\rangle$$
که یعنی:
$$\forall j\;(1\leq j \leq l \rightarrow \langle \underline{\rho},\langle l_j, \rho_j \rangle\rangle \in \mathcal{S}^r \llbracket \mathsf{L}_j : \mathsf{B}_j \rrbracket)$$
$$\Rightarrow \langle \underline{\rho}, \pi \rangle \in \mathsf{prefix}(\mathcal{S}^r \llbracket \mathsf{R}_k\rrbracket) \Rightarrow \langle \underline{\rho},\pi\rangle \in \mathsf{prefix(\mathcal{S}^r \llbracket R \bullet (?:\mathit{T})^*} \rrbracket).$$ 

پس در کل می‌توانیم، بگوییم
$$\langle \underline{\rho},\pi\rangle \in \mathsf{prefix(\mathcal{S}^r \llbracket R \bullet (?:\mathit{T})^*} \rrbracket)$$
و اثبات قضیه تمام می‌شود.
\end{proof}

حال به اثبات تمامیت می‌پردازیم.

\begin{thm}

	(قضیه تمامیت): اگر $\mathsf{P}$ یک برنامه، $\mathsf{R}$ یک عبارت منظم و $\underline{\rho}$ یک محیط اولیه باشند، آنگاه داریم:
$$\mathsf{P},\underline{\rho} \models \mathsf{R} \Rightarrow
\mathsf{P},\underline{\rho} \models_r \mathsf{R}$$
\end{thm}
\begin{proof}
	با برهان خلف این قضیه را ثابت می‌کنیم. شکل اثبات تا حدی شبیه به اثبات درستی است.
	$$\{\underline{\rho}\}\times \mathcal{S}^* \llbracket \mathsf{P} \rrbracket \nsubseteq \mathcal{M}\langle \underline{\rho}, \mathsf{R} \rangle \mathcal{S}^* \llbracket \mathsf{P} \rrbracket \Rightarrow
	\exists \pi\;(\langle \underline{\rho}, \pi \rangle \in \{\rho\}\times \mathcal{S}^* \llbracket \mathsf{P}\rrbracket \land \langle \underline{\rho}, \pi \rangle \notin \mathcal{M}\langle \underline{\rho},\mathsf{R}\rangle\mathcal{S}^* \llbracket \mathsf{P} \rrbracket)$$
	اگر فرض کنیم
	$\mathsf{dnf(R)=R_1+R_2+...+R_n}$
	و علاوه بر این، با توجه به آنچه در اثبات درستی گفتیم، فرض کنیم:
	$$\mathsf{R}_i \Bumpeq \mathsf{L_i^1:B_i^1 \bullet L_i^2:B_i^2 \bullet ... \bullet L_i^n:B_i^n}$$ 
	و
	$$\pi=\langle l_1 ,\rho_1 \rangle \langle l_2 ,\rho_2 \rangle ... \langle l_l ,\rho_l \rangle$$
	داریم:
	$$\forall i\;(1 \leq i \leq n \rightarrow 
	\langle \underline{\rho}, \pi \rangle \notin \mathcal{M}^\nmid \langle \underline{\rho}, \mathsf{R}_i \rangle \mathcal{S}^* \llbracket \mathsf{P} \rrbracket)$$
	$$\Rightarrow\forall i \;(1 \leq i \leq n \rightarrow\exists \mathsf{R'}_i (\mathcal{M}^t \langle \underline{\rho},\mathsf{R}_i \rangle \pi = \langle \mathit{F},\mathsf{R'}_i \rangle)).$$
	بنابراین:
	$$\exists j\; (\langle \underline{\rho},\langle l_j,\rho_j \rangle \rangle \notin \mathcal{S}^r \llbracket \mathsf{L_i^j:B_i^j} \rrbracket
	\Rightarrow \langle \underline{\rho},\pi \rangle \notin \mathcal {S}^r \llbracket \mathsf{R}_i \bullet (?:\mathit{T})^* \rrbracket).$$
	از نتیجه‌ی آخر می‌توانیم ثابت کنیم 
	$\langle \underline{\rho},\pi \rangle \notin \mathsf{prefix}(\mathcal {S}^r \llbracket \mathsf{R}_i \bullet (?:\mathit{T})^* \rrbracket)$.
	چون اگر غیر از این باشد، یعنی اگر فرض کنیم 
	$\langle \underline{\rho}, \pi' \rangle$
	 عضو 
	 $\mathsf{prefix}(\mathcal {S}^r \llbracket \mathsf{R}_i \bullet (?:\mathit{T})^* \rrbracket)$
	 هست، اما عضو \break
	 $\mathcal {S}^r \llbracket \mathsf{R}_i \bullet (?:\mathit{T})^* \rrbracket$
	 نیست، در آن صورت، اگر طول $\pi'$ بزرگتر یا مساوی $j$ باشد، به خاطر وجود 
	 $\langle l_j,\rho_j \rangle$
	 در $\pi$، خواهیم داشت $\pi \neq \pi'$ و اگر طول $\pi'$ کمتر از $j$ باشد، چون طول $\pi$ قطعا بزرگتر یا مساوی $j$ است( نتیجه از گزاره‌ی بالایی که دارای سور وجودی است)، پس باز هم 
	 $\pi \neq \pi'$.
	 توجه شود که 
	 $\langle \underline{\rho},\pi \rangle \notin \mathsf{prefix}( \mathcal {S}^r \llbracket \mathsf{R}_i \bullet (?:\mathit{T})^* \rrbracket)$
	 با توجه به آنچه گفتیم، با فرض در تناقض است و حکم ثابت می‌شود.
\end{proof}
















% !TeX root = UT-Thesis-Template.tex
\chapter{وارسی مدل ساختارمند}

در این فصل، به ادامه‌ی ساختارمندتر کردن کار می‌پردازیم. در فصل گذشته، ساختار عبارات منظم را به تعریف وارسی مدل اضافه کردیم و حالا می‌خواهیم، ساختار زبانمان را به کار اضافه کنیم. این آخرین تلاش \cite{calcul} برای گسترش کار بوده است. یعنی وارسی مدل در صورت جدید تعریف شده است و معادل بودن آن با صورت قبلی وارسی مدل ثابت شده است و پس از آن کار پایان می‌پذیرد. 

چون تعریف صورت جدید روی ساختار زبان انجام گرفته است، جزئیات بسیار طولانی‌ای دارد. همین موضوع باعث شده است، تا اثبات‌ برابری این صورت با صورت قبلی هم بسیار مفصل و حجیم باشد. این اثبات در \cite{calcul} به طور کامل حین معرفی هر حالت تعریف بیان شده است. بنابراین، از ارائه‌ی دوباره‌ی این جزئیات خودداری کرده‌ایم.

\begin{defn}
	
تابع
$\mathcal{\hat{M}}$
 را از نوع
 $\mathbb{(\underline{EV} \times R)} \rightarrow  \mathit{P}({\mathfrak{S}^{+\infty})}
\rightarrow ( \mathbb{\underline{EV}} \times \mathit{P}(\mathfrak{S}^{+\infty})) $
\gls{strucmodelc} می‌نامیم( ضابطه‌ی تابع در ادامه‌ی متن آمده است).
\end{defn}
در ادامه، ممکن است به‌جای 
$\mathcal{\hat{M}}\langle \underline{\rho}, \mathsf{R} \rangle \mathcal{S}^* \llbracket \mathsf{P} \rrbracket$
از 
$\mathcal{\hat{M}}\langle \underline{\rho}, \mathsf{R} \rangle \llbracket \mathsf{P} \rrbracket$
استفاده کرده باشیم، یعنی در اشاره به تابع $\mathcal{S}^*$ به براکت‌ها
$\llbracket \; \rrbracket$
قناعت کرده باشیم.



تعریف روی ساختار مجموعه‌ی 
$\mathbb{P \cup Sl \cup S}$
انجام شده است. تقریبا کاری شبیه به اثبات لمی که در بحث تصمیم ناپذیری در فصل سوم داشتیم.
در ادامه، قسمت‌های مختلف تعریف $\mathcal{\hat{M}}$ را به ازای برنامه‌ی $\mathsf{P}$، محیط اولیه‌ی $\underline{\rho}$ و عبارت منظم $\mathsf{R}$تعریف می‌کنیم. یعنی در حال تعریف 
$\mathcal{\hat{M}} \langle \underline{\rho} , \mathsf{R} \rangle \mathcal{S}^* 
\llbracket \mathsf{P} \rrbracket$
هستیم، روی ساختار برنامه‌ها یعنی $\mathsf{P}$.
$$\blacktriangleleft \mathsf{P=Sl:}$$
$$\mathcal{\hat{M}} \langle \underline{\rho} , \mathsf{R} \rangle \mathcal{S}^* \llbracket \mathsf{P} \rrbracket=
\bigcup_{i=1}^n \{\langle \underline{\rho} , \pi \rangle | \exists \mathsf{R'} \in \mathbb{R}, \; \langle \pi , \mathsf{R'} \rangle \mathcal{\hat{M}^\nmid}
\langle \underline{\rho}, \mathsf{R}_i \rangle \mathcal{S}^* \llbracket \mathsf{Sl} \rrbracket \}
$$  
$$\mathsf{where\; dnf(R)=R_1 + R_2 + ... + R_n}$$
اثبات برابری این قسمت از تابع با صورت فصل قبل با اینکه 
$\mathcal{\hat{M}^\nmid} \langle \underline{\rho}, \mathsf{R} \rangle \llbracket \mathsf{Sl} \rrbracket$
هنوز تعریف نشده است، در \cite{calcul} آمده است. استدلال این است که برابری 
$\mathcal{\hat{M}} \langle \underline{\rho}, \mathsf{R} \rangle \llbracket \mathsf{Sl} \rrbracket=
\mathcal{{M}} \langle \underline{\rho}, \mathsf{R} \rangle \llbracket \mathsf{Sl} \rrbracket$
در فرض استقرا آمده است. کلیت اثبات هم این است که از باز کردن تعریف
$\mathcal{{M}} \langle \underline{\rho}, \mathsf{R} \rangle \llbracket \mathsf{P} \rrbracket$
با استفاده‌ی مستقیم تعاریف و بدون تکنیک خاصی به 
$\mathcal{\hat{M}} \langle \underline{\rho}, \mathsf{R} \rangle \llbracket \mathsf{P} \rrbracket$
رسیده است.

در ادامه با توجه به تعریف قبل، به بیان تعریف 
$\mathcal{\hat{M}^\nmid}$
پرداخته شده است. این تنها بخش تابع $\mathcal{\hat{M}}$ است که معرفی نشده است و با مشخص شدن آن معنای 
$\mathcal{\hat{M}}$
به ازای برنامه‌های مختلف مشخص می‌شود. 

این نکته را در نظر داریم که 
$\mathcal{\hat{M}^\nmid}$
در عمل روی مجوعه‌ی معنای برنامه‌ها تعریف می‌شود. مثلا، به ازای  
$\Pi \subseteq \mathfrak{S}^{+\infty}$
دلخواه که مساوی معنای یک برنامه نباشد، اینکه این تابع با یک محیط اولیه و یک عبارت منظم چه خروجی‌ای دارد، برای ما اهمیتی ندارد. در واقع، تعریف تابع اصلا به ازای چنین ورودی‌ای خروجی ندارد. به عبارت دیگر، \gls*{partialf} است. مشابه 
$\mathcal{M^\nmid}$
خروجی 
$\mathcal{\hat{M}^\nmid}$
هم یک زوج مرتب شامل $\pi$ ای است که $\mathsf{R}$ را ارضا کرده است، به همراه یک عبارت منظم بدون $+$ که بخشی از $\mathsf{R}$ را نشان می‌دهد که با $\pi$ تطابق داده نشده است.
$$\blacktriangleleft\mathcal{\hat{M}^\nmid} \langle \underline{\rho}, \mathsf{\varepsilon} \rangle \llbracket \mathsf{S} \rrbracket
=
\{\langle \pi , \varepsilon \rangle | \pi \in \mathcal{S}^* \llbracket \mathsf{S} \rrbracket \}$$

برای $\mathsf{Sl=SL' \; S}$ و $\mathsf{R} \in \mathbb{R^\nmid} \cap \mathbb{R^+}$ داریم:

$$\blacktriangleleft\mathcal{\hat{M}^\nmid} \langle \underline{\rho}, \mathsf{R} \rangle \llbracket \mathsf{Sl} \rrbracket
=
\mathcal{\hat{M}^\nmid} \langle \underline{\rho}, \mathsf{R} \rangle \llbracket \mathsf{Sl'} \rrbracket
\cup$$ 
$$\{ \langle \pi \langle \mathsf{at}\llbracket \mathsf{S} \rrbracket, \rho \rangle \pi'
, \mathsf{R''} \rangle | \langle \pi \langle at \llbracket \mathsf{S}\rrbracket, \rho \rangle, \mathsf{R'} \rangle \in \mathcal{\hat{M}^\nmid}  \langle \underline{\rho} , \mathsf{R} \rangle \llbracket \mathsf{Sl'} \rrbracket\land$$
 $$\langle \langle at \llbracket \mathsf{S}\rrbracket , \rho \rangle \pi' , \mathsf{R''} \rangle \in \mathcal{\hat{M}^\nmid}
\langle \underline{\rho},\mathsf{R'} \rangle   \llbracket \mathsf{S} \rrbracket
 \}$$
از اینجا به بعد، با تعاریف طویل‌تری از آنچه تا حالا داشتیم، روبرو هستیم. هرچند که مفهوم چندان پیچیده‌ای پشت این تعاریف نیست. به ‌طور خلاصه، تعریف بالا می‌گوید، تعریف 
$\mathcal{\hat{M}^\nmid} \langle \underline{\rho}, \mathsf{R} \rangle \llbracket \mathsf{Sl} \rrbracket$
وابسته به تعریف 
$\mathcal{\hat{M}^\nmid} \langle \underline{\rho}, \mathsf{R} \rangle \llbracket \mathsf{S} \rrbracket$
و
$\mathcal{\hat{M}^\nmid} \langle \underline{\rho}, \mathsf{R} \rangle \llbracket \mathsf{Sl'} \rrbracket$
است. ردهای پیشوندی‌ای که داخل این دو مجموعه هستند، به یکدیگر چسبانده می‌شوند، به‌ طوریکه اول ردهای داخل 
$\mathcal{\hat{M}^\nmid} \langle \underline{\rho}, \mathsf{R} \rangle \llbracket \mathsf{Sl'} \rrbracket$
قرار می‌گیرند و بعد ردهای داخل 
$\mathcal{\hat{M}^\nmid} \langle \underline{\rho}, \mathsf{R} \rangle \llbracket \mathsf{S} \rrbracket$.
ردهای داخل 
$\mathcal{\hat{M}^\nmid} \langle \underline{\rho}, \mathsf{R} \rangle \llbracket \mathsf{Sl'} \rrbracket$
به تنهایی نیز داخل 
$\mathcal{\hat{M}^\nmid} \langle \underline{\rho}, \mathsf{R} \rangle \llbracket \mathsf{Sl} \rrbracket$
قرار می‌گیرند. این تعریف بر اساس تعریف 
$\mathcal{S}^* \llbracket \mathsf{Sl} \rrbracket$
ارائه شده است. اساس گرفتن تعریف تابع 
$\mathcal{S}^*$
در کنار توجه به تعریف تابع 
$\mathcal{M}$
که در فصل پیش ارائه شد، در ادامه‌ی تعریف 
$\mathcal{\hat{M^\nmid}}$
	 نیز حضور دارد.

برای $\mathsf{Sl=\epsilon}$ و $\mathsf{R} \in \mathbb{R^\nmid} \cap \mathbb{R^+}$ داریم:

$$\blacktriangleleft\mathcal{\hat{M}^\nmid} \langle \underline{\rho}, \mathsf{R} \rangle \llbracket \mathsf{Sl} \rrbracket
=$$
$$\{ \langle \langle at \llbracket \mathsf{Sl} \rrbracket , \rho \rangle , \mathsf{R'} \rangle | \langle \underline{\rho} , \langle at \llbracket \mathsf{Sl} \rrbracket, \rho \rangle \rangle \in \mathcal{S}^r \llbracket \mathsf{L:B} \rrbracket
\}$$
$$\mathsf{where \; fstnxt(R)=\langle L:B,R' \rangle}$$
یعنی خروجی تابع به ازای این ورودی مجموعه‌ای است، شامل همه‌ی ردهای پیشوندی تک عضوی‌ای که محیط آن‌ها اولین سر عبارت منظم$(\mathsf{L:B})$ را ارضا می‌کند. به عبارت دیگر، هر محیطی که این اتم را ارضا کند، برچسب این لیست عبارت‌های دستوری را در این مجموعه می‌آورد( به همراه ادامه‌ی عبارت منظم).

برای $\mathsf{S=\;x \doteq A;}$ و $\mathsf{R} \in \mathbb{R^\nmid} \cap \mathbb{R^+}$ داریم:
$$\blacktriangleleft\mathcal{\hat{M}^\nmid} \langle \underline{\rho}, \mathsf{R} \rangle \llbracket \mathsf{S} \rrbracket
=$$
$$\{
\langle \langle at \llbracket \mathsf{S} \rrbracket, \rho \rangle , \mathsf{R'} \rangle |
\langle \underline{\rho}, \langle at \llbracket \mathsf{S} \rrbracket, \rho \rangle \rangle \in \mathcal{S}^r \llbracket \mathsf{L:B} \rrbracket \}$$
$$\cup \{ \langle \langle at \llbracket \mathsf{S} \rrbracket , \rho \rangle  \langle aft \llbracket \mathsf{S} \rrbracket, 
\rho [\mathsf{x}\leftarrow \mathcal{A}\llbracket \mathsf{A} \rrbracket \rho ]\rangle , \varepsilon \rangle | \mathsf{R'} \in \mathbb{R_\varepsilon} \land$$
$$
\langle \underline{\rho}, \langle at \llbracket \mathsf{S} \rrbracket , \rho \rangle \rangle \in 
\mathcal{S}^r \llbracket \mathsf{L:B} \rrbracket
\}$$
$$\cup \{
\langle \langle at \llbracket \mathsf{S} \rrbracket , \rho \rangle \langle aft \llbracket \mathsf{S} \rrbracket , \rho [\mathsf{x} \leftarrow \mathcal \llbracket \mathsf{A}\rrbracket \rho]\rangle, \mathsf{R''}\rangle | \mathsf{R'} \notin \mathbb{R_\varepsilon} \land$$
$$\langle \underline{\rho},\langle at \llbracket \mathsf{S} \rrbracket , \rho \rangle \rangle \in \mathcal{S}^r \llbracket \mathsf{L:B} \rrbracket \land \langle \mathsf{L':B',R''} \rangle = 
\mathsf{fstnxt(R') }\land$$
$$\langle \underline{\rho}, \langle aft \llbracket \mathsf{S} \rrbracket,\rho 
[\mathsf{x} \leftarrow \mathcal{A}\llbracket \mathsf{A} \rrbracket \rho]\rangle \rangle \in 
\mathcal{S}^r \llbracket \mathsf{L':B'} \rrbracket
\}$$
$$\mathsf{where \; fstnxt(R)=\langle L:B,R' \rangle}$$

این تابع عبارت‌ دستوری را به همراه زوج مرتبی شامل محیط اولیه و عبارت منظم می‌گیرد، همان خروجی‌ای که در حالت قبلی برمی‌گردانْد را برمی‌گردانَد، سپس نسبت به اینکه پس از تغییر در محیط‌ها (در اثر اجرای عبارت‌ دستوری مقدار دهی) یک رد پیشوندی با ادامه‌ی عبارت منظم سازگار باشد یا نباشد، زوج‌هایی را متشکل از رد پیشوندی و عبارت منظم به خروجی اضافه می‌کند.

از این ۴ حالت تنها اثبات حالت آخر در \cite{calcul} آورده شده است. اثبات دیگر حالات را هم می‌توان در همین اثبات که مفصل‌تر است، دید. اثبات سر راست است و در آن از جایگذاری تساوی‌های واضح استفاده شده است و جزئیات کافی دارد.

برای عبارت منظم 
$\mathsf{R} \in \mathbb{R}^\nmid \cup \mathbb{R^+}$
و 
$\mathsf{S= \; if \; (B) \; S_t}$
داریم:
$$\blacktriangleleft\mathcal{\hat{M}^\nmid} \langle \underline{\rho},\mathsf{R} \rangle \llbracket \mathsf{S} \rrbracket=$$
$$\{\langle \langle at \llbracket \mathsf{S} \rrbracket , \rho \rangle , \mathsf{R'} \rangle | \langle \underline{\rho} , \langle at \llbracket \mathsf{S} \rrbracket , \rho \rangle \rangle \in \mathcal{S}^r \llbracket \mathsf{L':B'} ]\rrbracket \}$$
$$\cup \{\langle \langle at \llbracket \mathsf{S} \rrbracket , \rho \rangle \langle at \llbracket \mathsf{S_t} \rrbracket, \rho \rangle \pi, \mathsf{R''} | \mathcal{B}\llbracket\mathsf{B} \rrbracket \rho = \mathit{T} \land$$
$$\langle \underline{\rho},\langle at \llbracket \mathsf{S} \rrbracket , \rho \rangle \rangle \in \mathcal{S}^r \llbracket \mathsf{L':B'} \rrbracket \land$$
$$\langle \langle at \llbracket \mathsf{S_t} \rrbracket , \rho \rangle \pi,\mathsf{R''} \rangle \in \mathcal{\hat{M}^\nmid} \langle \underline{\rho}, \mathsf{R'} \rangle\llbracket \mathsf{S_t} \rrbracket \}$$
$$\cup \{\langle \langle at \llbracket \mathsf{S} \rrbracket, \rho \rangle \langle aft \llbracket \mathsf{S} \rrbracket , \rho \rangle , \varepsilon \rangle | \mathcal{B}\llbracket \mathsf{B} \rrbracket \rho = \mathit{F} \land \mathsf{R'} \in \mathbb{R_\varepsilon} \land$$
$$\langle \underline{\rho} , \langle \llbracket \mathsf{S} \rrbracket , \rho \rangle \rangle \in \mathcal{S}^r \llbracket \mathsf{L':B'} \rrbracket\}$$
$$\cup \{\langle \langle at \llbracket \mathsf{S} \rrbracket , \rho \rangle \langle aft \llbracket \mathsf{S} \rrbracket , \rho \rangle , \mathsf{R''} \rangle | \mathcal{B} \llbracket \mathsf{B} \rrbracket \rho = \mathit{F} \land \mathsf{R'} \notin \mathbb{R_\varepsilon} \land$$
$$\langle \underline{\rho}, \langle at \llbracket \mathsf{S} \rrbracket , \rho \rangle \rangle \in \mathcal{S}^r\llbracket \mathsf{L':B'} \rrbracket \land \langle \mathsf{L'':B''} , \mathsf{R''} \rangle = \mathsf{fstnxt(R')} \land$$
$$\langle \underline{\rho}, \langle aft \llbracket \mathsf{S} \rrbracket , \rho \rangle \rangle \in \mathcal{S}^r \llbracket \mathsf{L'':B''} \rrbracket \}$$
$$\mathsf{where\; fstnxt(R)=\langle L':B', R' \rangle}$$


برای عبارت منظم 
$\mathsf{R} \in \mathbb{R}^\nmid \cup \mathbb{R^+}$
و 
$\mathsf{S= \; if \; (B) \; S_t \; else \; S_f}$
داریم:

$$\blacktriangleleft\mathcal{\hat{M}^\nmid} \langle \underline{\rho},\mathsf{R} \rangle \llbracket \mathsf{S} \rrbracket=$$
$$\{\langle \langle at \llbracket \mathsf{S} \rrbracket , \rho \rangle , \mathsf{R'} \rangle | \langle \underline{\rho} , \langle at \llbracket \mathsf{S} \rrbracket , \rho \rangle \rangle \in \mathcal{S}^r \llbracket \mathsf{L':B'} ]\rrbracket \}$$
$$\cup \{\langle \langle at \llbracket \mathsf{S} \rrbracket , \rho \rangle \langle at \llbracket \mathsf{S_t} \rrbracket, \rho \rangle \pi, \mathsf{R''} | \mathcal{B}\llbracket\mathsf{B} \rrbracket \rho = \mathit{T} \land$$
$$\langle \underline{\rho},\langle at \llbracket \mathsf{S} \rrbracket , \rho \rangle \rangle \in \mathcal{S}^r \llbracket \mathsf{L':B'} \rrbracket \land$$
$$\langle \langle at \llbracket \mathsf{S_t} \rrbracket , \rho \rangle \pi,\mathsf{R''} \rangle \in \mathcal{\hat{M}^\nmid} \langle \underline{\rho}, \mathsf{R'} \rangle  \llbracket \mathsf{S_t} \rrbracket \}$$
$$\cup \{\langle \langle at \llbracket \mathsf{S} \rrbracket , \rho \rangle \langle at \llbracket \mathsf{S_f} \rrbracket, \rho \rangle \pi, \mathsf{R''} | \mathcal{B}\llbracket\mathsf{B} \rrbracket \rho = \mathit{F} \land$$
$$\langle \underline{\rho},\langle at \llbracket \mathsf{S} \rrbracket , \rho \rangle \rangle \in \mathcal{S}^r \llbracket \mathsf{L':B'} \rrbracket \land$$
$$\langle \langle at \llbracket \mathsf{S_f} \rrbracket , \rho \rangle \pi,\mathsf{R''} \rangle \in \mathcal{\hat{M}^\nmid}  \langle \underline{\rho}, \mathsf{R'} \rangle \llbracket \mathsf{S_f} \rrbracket \}$$
$$\mathsf{where\; fstnxt(R)=\langle L':B', R' \rangle}$$
دو قسمت بالا در مورد عبارت‌های دستوری شرطی هستند. در مورد نوع اولی شرطی، یک رد پیشوندی را در معنای $\mathsf{S}$ در نظر بگیرید. بسته به اینکه طبق محیط حاضر در اولین وضعیت این رد پیشوندی، عبارت بولی مقدار صحیح یا غلط داشته باشد، حضور این رد پیشوندی( در یک زوج، به همراه یک عبارت منظم) داخل 
$\mathcal{\hat{M}^\nmid} \langle \underline{\rho},\mathsf{R} \rangle \llbracket \mathsf{S} \rrbracket$
تعیین می‌شود. اگر عبارت بولی در محیط مذکور مقدار صحیح داشته باشد، 
در معنای نوع اول عبارت‌ دستوری شرطی، پس از تطبیق سر عبارت منظم با اولین وضعیت هر رد پیشوندی، بر اساس اینکه در دومین وضعیت رد پیشوندی، عبارت بولی برقرار باشد یا نباشد، وابسته به این می‌شود که آیا اگر از وضعیت دوم به بعد این رد پیشوندی( که خود یک رد پیشوندی است) در 
$\mathcal{\hat{M}^\nmid} \langle \underline{\rho},\mathsf{R'} \rangle \llbracket \mathsf{S_t} \rrbracket$
( که $\mathsf{R'}$ نیز دم $\mathsf{R}$ است) حضور دارد یا خیر. اگر عبارت بولی در محیط مذکور ارزش غلط داشته باشد، حضور رد پیشوندی در
$\mathcal{\hat{M}^\nmid} \langle \underline{\rho},\mathsf{R} \rangle \llbracket \mathsf{S} \rrbracket$
وابسته به سازگاری وضعیت دوم رد پیشوندی با سر $\mathsf{R'}$ خواهد بود.

در نوع دوم عبارت‌ دستوری شرطی نیز، تعریف شبیه به نوع اول است، با این تفاوت که اگر عبارت بولی در محیط اولین وضعیت رد پیشوندی ارزش غلط داشته باشد، اتفاقی شبیه به حالت درست می‌افتد.


 برای عبارت منظم 
 $\mathsf{R} \in \mathbb{R}^\nmid \cup \mathbb{R^+}$
 و 
 $\mathsf{S= \; break;}$
 داریم:
$$\blacktriangleleft \mathcal{\hat{M}^\nmid} \langle \underline{\rho}, \mathsf{R} \rangle 
\llbracket \mathsf{S} \rrbracket= $$
$$\{\langle \langle at \llbracket \mathsf{S} \rrbracket,\rho \rangle , \mathsf{R'} \rangle | \langle \underline{\rho} , \langle at \llbracket \mathsf{S} \rrbracket , \rho \rangle \rangle \in \mathcal{S}^r \llbracket \mathsf{L:B} \rrbracket\}$$
$$\cup \{ \langle \langle at \llbracket \mathsf{S} \rrbracket , \rho \rangle \langle brk-to \llbracket \mathsf{S} \rrbracket , \rho \rangle , \varepsilon \rangle | \mathsf{R'} \in \mathbb{R_\varepsilon} \land$$
$$\langle \underline{\rho}, \langle at \llbracket \mathsf{S} \rrbracket, \rho \rangle \rangle \in \mathcal{S}^r \llbracket \mathsf{L:B} \rrbracket \}$$
$$\cup \{ \langle \langle at \llbracket \mathsf{S} \rrbracket , \rho \rangle \langle brk-to \llbracket \mathsf{S} \rrbracket,\rho \rangle ,\mathsf{R''} \rangle | \mathsf{R'} \notin \mathbb{R_\varepsilon} \land$$
$$ \langle \underline{\rho},\langle at \llbracket \mathsf{S} \rrbracket , \rho \rangle \rangle \in \mathcal{S}^r \llbracket \mathsf{L:B} \rrbracket \land \langle \mathsf{L':B',R''} \rangle = \mathsf{fstnxt(R') \land}$$
$$\langle \underline{\rho}, \langle brk-to \llbracket \mathsf{S} \rrbracket, \rho \rangle \rangle \in \mathcal{S}^r \llbracket \mathsf{L':B'} \rrbracket \}$$
$$\mathsf{where\; fstnxt(R)=\langle L':B', R' \rangle}$$
 با توجه به موارد قبلی، اینکه این قسمت از تعریف چه معنایی دارد و به چه علت به این شکل است، قابل درک است.
 
 برای عبارت منظم 
 $\mathsf{R} \in \mathbb{R}^\nmid \cup \mathbb{R^+}$
 و 
 $\mathsf{S= \; while \; (B) \; S_b:}$
 داریم:
 
$$\blacktriangleleft\mathcal{\hat{M}^\nmid} \langle \underline{\rho}, \mathsf{R} \rangle \llbracket \mathsf{S} \rrbracket = lfp^\subseteq \; (\mathcal{\hat{F^\nmid}} \langle \underline{\rho}, \mathsf{R} \rangle \llbracket \mathsf{S} \rrbracket)$$
$$\mathsf{where} \;  \mathcal{\hat{F^\nmid}} \langle \underline{\rho}, \mathsf{R} \rangle X = \{ \langle \langle at \llbracket \mathsf{S} \rrbracket , \rho \rangle , \mathsf{R'} \rangle | \rho \in \mathbb{EV} \land \langle \underline{\rho}, \langle at \llbracket \mathsf{S} \rrbracket , \rho \rangle \rangle \in \mathcal{S}^r \llbracket \mathsf{L':B'} \rrbracket \}$$
$$\cup \{\langle \pi_2 \langle at \llbracket \mathsf{S} \rrbracket , \rho \rangle \langle aft \llbracket \mathsf{S}\rrbracket , \rho \rangle , \varepsilon \rangle | \langle \pi_2 \langle at \llbracket \mathsf{S} \rrbracket , \rho \rangle , \varepsilon \rangle \in X \land$$
$$\mathcal{B}\llbracket \mathsf{B} \rrbracket \rho = \mathit{F} \}$$
$$\cup \{ \langle \pi_2 \langle at \llbracket \mathsf{S} \rrbracket , \rho \rangle \langle aft \llbracket \mathsf{S} \rrbracket , \rho \rangle , \varepsilon \rangle | \langle \pi_2 \langle at \llbracket \mathsf{S} \rrbracket , \rho \rangle , \mathsf{R''} \rangle \in X \land $$
$$\mathcal{B}\llbracket \mathsf{B} \rrbracket \rho = \mathit{F} \land \mathsf{R''} \notin \mathbb{R_\varepsilon} \land \langle \mathsf{L':B',R'} \rangle = \mathsf{fstnxt(R'')} \land \mathsf{R'} \in \mathbb{R_\varepsilon} \land$$
$$\langle \underline{\rho} ,\langle at \llbracket \mathsf{S} \rrbracket , \rho \rangle \rangle \in \mathcal{S}^r \llbracket \mathsf{L':B'} \rrbracket\}$$
$$\cup \{\langle \pi_2 \langle at \llbracket \mathsf{S} \rrbracket , \rho \rangle \langle aft \llbracket \mathsf{S}\rrbracket , \rho \rangle , \mathsf{R'} \rangle | \langle \pi_2 \langle at \llbracket \mathsf{S} \rrbracket , \rho \rangle , \mathsf{R''} \rangle \in X \land$$
$$\mathcal{B}\llbracket \mathsf{B} \rrbracket \rho = \mathit{F} \land \mathsf{R''} \notin \mathbb{R_\varepsilon} \land \langle \mathsf{L':B',R'} \rangle = \mathsf{fstnxt(R'')} \land \mathsf{R'''} \notin \mathbb{R_\varepsilon} \land$$
$$\langle \underline{\rho}, \langle at \llbracket \mathsf{S} \rrbracket , \rho \rangle \rangle \in \mathcal{S}^r \llbracket \mathsf{L':B'} \rrbracket \land \langle \mathsf{L'':B'' , R'} \rangle = \mathsf{fstnxt(R''')} \land$$
$$\langle \underline{\rho}, \langle aft \llbracket \mathsf{S} \rrbracket , \rho \rangle \rangle \in \mathsf{S}^r \llbracket \mathsf{L'':B''} \rrbracket \}$$
$$\cup \{\langle \pi_2 \langle at \llbracket \mathsf{S} \rrbracket , \rho \rangle \langle at \llbracket \mathsf{S_b} \rrbracket , \rho \rangle \pi_3, \varepsilon \rangle | \langle \pi_2 \langle at \llbracket \mathsf{S} \rrbracket , \rho \rangle , \varepsilon \rangle \in X \land$$
$$\mathcal{B}\llbracket \mathsf{B} \rrbracket \rho = \mathit{T}\land \langle at \llbracket \mathsf{S_b} \rrbracket , \rho \rangle \pi_3 \in \mathcal{S}^* \llbracket \mathsf{S_b} \rrbracket \}$$
$$\cup\{\langle \pi_2 \langle at \llbracket \mathsf{S} \rrbracket , \rho \rangle \langle at \llbracket \mathsf{S_b} \rrbracket , \rho \rangle \pi_3 , \varepsilon \rangle | \langle \pi_2 \langle at \llbracket \mathsf{S} \rrbracket , \rho \rangle , \mathsf{R''} \rangle \in X \land$$
$$\mathcal{B} \llbracket \mathsf{B} \rrbracket \rho = \mathit{T} \land \mathsf{R''} \notin \mathbb{R_\varepsilon} \land \langle \mathsf{L:B} , \varepsilon \rangle = \mathsf{fstnxt(R'')} \land$$	
$$\langle \underline{\rho} ,\langle at \llbracket \mathsf{S} \rrbracket, \rho \rangle \rangle \in \mathcal{S}^r \llbracket \mathsf{L:B} \rrbracket \land \langle at \llbracket \mathsf{S_b} \rrbracket , \rho \rangle \pi_3 \in \mathcal{S}^* \llbracket \mathsf{S_b} \rrbracket\}$$ 
$$\cup \{ \langle \pi_2 \langle at \llbracket \mathsf{S} \rrbracket , \rho \rangle \langle at \llbracket \mathsf{S_b} \rrbracket , \rho \rangle \pi_3 , \mathsf{R'} \rangle | \langle \pi_2 \langle at \llbracket \mathsf{S} \rrbracket , \rho \rangle , \mathsf{R''} \rangle \in X \land$$
$$\mathcal{B} \llbracket \mathsf{B} \rrbracket \rho = \mathit{T} \land \mathsf{R''} \notin \mathbb{R_\varepsilon} \land \langle \mathsf{L:B,R''''} \rangle = \mathsf{fstnxt(R'')} \land$$
$$\langle \underline{\rho} , \langle at \llbracket \mathsf{S}\rrbracket , \rho \rangle \rangle \in \mathcal{S}^r \llbracket \mathsf{L:B} \rrbracket \land \mathsf{R''''} \notin \mathbb{R_\varepsilon} \land$$
$$\langle \mathsf{L':B',R'''} \rangle = \mathsf{fstnxt(R'''')} \land \langle \underline{\rho} , \langle at \llbracket \mathsf{S_b}\rrbracket , \rho \rangle \rangle \in \mathcal{S}^r \llbracket \mathsf{L':B'} \rrbracket \land$$
$$\langle \langle at \llbracket \mathsf{S_b} \rrbracket, \rho \rangle \pi_3 , \mathsf{R'} \rangle \in \mathcal{\hat{M}^\nmid} \langle \underline{\rho} , \mathsf{R'''} \rangle \llbracket \mathsf{S_b} \rrbracket \}$$
$$ \mathsf{where \; \langle L':B' , R' \rangle = fstnxt(R)}$$

تفاوت تعریف 
$\mathcal{\hat{M}}^\nmid$
 برای عبارت‌ دستوری حلقه با سایر عبارت‌های دستوری، حضور یک تابع به نام $\mathcal{\hat{F}^\nmid}$ در تعریف معنای آن است. در واقع، وارسی مدل به صورت کوچک‌ترین نقطه‌ثابت این تابع تعریف می‌شود. این همان کاری است که در تعریف معنای اجزای زبان هم انجام شد و چون می‌خواهیم ساختار زبان را به صورت وارسی مدل اضافه کنیم، انتظار داریم که عملگر نقطه‌ثابت هم در تعریف حضور پیدا کند. تابع $\mathcal{\hat{F}^\nmid}$ مثل یک دور اجرای حلقه‌ عمل می‌کند، منتها در همین حین، سازگاری ردهای پیشوندی را با عبارت منظم بررسی می‌کند و ردهای پیشوندی‌ای را که با عبارت منظم ناسازگار هستند، مجموعه‌ی خروجی‌اش حذف می‌کند.


برای عبارت منظم 
$\mathsf{R} \in \mathbb{R}^\nmid \cup \mathbb{R^+}$
و 
$\mathsf{S= \; ; \;:}$
داریم:
$$\blacktriangleleft\mathcal{\hat{M}^\nmid}  \langle \underline{\rho}, \mathsf{R} \rangle \llbracket \mathsf{S} \rrbracket = \{\langle \langle at \llbracket \mathsf{S} \rrbracket , \rho \rangle , \mathsf{R'} \rangle | \langle \underline{\rho}, \langle at \llbracket \mathsf{S} \rrbracket , \rho \rangle \rangle \in \mathcal{S}^r \llbracket \mathsf{L:B}\rrbracket\}$$
$$\mathsf{where \; fstnxt (R) = \langle L:B,R' \rangle}$$
	
	

برای عبارت منظم 
$\mathsf{R} \in \mathbb{R}^\nmid \cup \mathbb{R^+}$
و 
$\mathsf{S= \; \{Sl\} \;:}$
داریم:
$$\blacktriangleleft \mathcal{\hat{M}^\nmid} \langle \underline{\rho}, \mathsf{R} \rangle \llbracket \{\mathsf{Sl}\} \rrbracket =
 \mathcal{\hat{M}^\nmid} \langle \underline{\rho}, \mathsf{R} \rangle \llbracket \mathsf{Sl} \rrbracket$$


همین طور، درست بودن یک ویژگی $\mathsf{R} \in \mathbb{R}$ را برای برنامه‌ی $\mathsf{P}$ و محیط اولیه‌ی $\underline{\rho}$ با 
$$\mathsf{P} , \underline{\rho} \models_s \mathsf{R}$$
نشان می‌دهیم و برقرار بودن این شرط به شکل زیر تعریف می‌شود:
$$\mathsf{P} , \underline{\rho} \models_s \mathsf{R} \iff \{\underline{\rho}\} \times \mathcal{S}^* \llbracket \mathsf{P} \rrbracket \subseteq \hat{\mathcal{M}} \langle \underline{\rho} , \mathsf{R} \rangle \mathcal{S}^* \llbracket \mathsf{P} \rrbracket$$

در اینجا تعریف توابع مربوط به وارسی مدل ساختارمند به پایان می‌رسد.

 
 
 
 
 


\chapter{نتیجه گیری}

\chapter*{واژه‌نامه فارسی به انگلیسی}
\chapter*{ واژه‌نامه انگلیسی به فارسی}

\newpage
\begin{latin}
	\bibliographystyle{abbrv}
	\bibliography{references}
\end{latin}

\begin{latin}
\begin{abstract}
Abstract goes here...
\end{abstract}
\newpage

%Title page ---------------------------------------------
\begin{figure}
\centering
\includegraphics[height=2.5cm]{UT-Logo.pdf}
\end{figure}
\begin{center}

College of Science\\
School of Mathematics, Statistics, and Computer Science
\end{center}

\begin{center}
%%%%%%%%%%%%%
\end{center}

\begin{center}
\huge{Thesis Title}
\end{center}

\begin{center}
%%%
\end{center}

\begin{center}
\textbf{
Author name
}
\end{center}

\begin{center}
\begin{tabular}{rr}
Supervisor: & name\\
Co-Supervisor: & name\\
Advisor: & name
\end{tabular}
\end{center}

\vspace{3cm}
\begin{center}
A thesis submitted to Graduate Studies Office\\
in partial fulfillment of the requirements for the degree of \\
B.Sc./Master of Science/Doctor of Philosophy in\\
Pure Mathematics/ Applied Mathematics/ Statistics/ Computer Science
\end{center}

\begin{center}
yyyy
\end{center}

%\pagestyle{empty}
%\pagenumbering{}

\end{latin}

\end{document}
